% Lois de Newton et conservation de l'énergie mécanique — Physique Tle D — Cours signé OnBuch+
% Programme officiel MINESEC. Compiler avec : tectonic P05-lois-de-newton-energie-mecanique.tex
\newcommand{\DOCMATIERE}{Physique}
\newcommand{\DOCNIVEAU}{Tle D}
\newcommand{\DOCMODULE}{Module 2 : Mouvements et interactions — évolutions temporelles des systèmes mécaniques}
\newcommand{\DOCLECON}{P05}
\newcommand{\DOCTITRE}{Lois de Newton et conservation de l'énergie mécanique}
% OnBuch+ — préambule « cours signés » ENRICHI (compilé avec tectonic / XeLaTeX)
% Système visuel « Pop » d'OnBuch+ : fond crème (jamais blanc), bordures encre
% épaisses, ombres « dures » décalées, titres Archivo Black en capitales.
% Version MATHÉMATIQUES : graphes (pgfplots/tikz), boîtes pédagogiques
% (définition, propriété, méthode, attention, le savais-tu, exercice résolu,
% exercice, TP, bilan), page de garde et signature OnBuch+ ancrée en bas.
%
% Macros attendues dans main.tex AVANT \input{preamble} :
%   \DOCMATIERE  \DOCNIVEAU  \DOCMODULE  \DOCLECON (numéro)  \DOCTITRE
\documentclass[11pt]{article}

\usepackage{fontspec}
\usepackage[french]{babel}
\usepackage[a4paper,margin=2cm,top=3.4cm,bottom=2.8cm,headheight=1.4cm,footskip=1.3cm]{geometry}
\usepackage{amsmath,amssymb,amsfonts}
\usepackage{mathtools} % \xmapsto, \coloneqq... souvent utilisés par les modèles
\usepackage{cancel} % simplifications barrées (bilans de forces)
% \abs{z} (module, valeur absolue) et \norm{u} (norme), utilisés par les modèles
\providecommand{\abs}{}\let\abs\relax\DeclarePairedDelimiter{\abs}{\lvert}{\rvert}
\providecommand{\norm}{}\let\norm\relax\DeclarePairedDelimiter{\norm}{\lVert}{\rVert}
\usepackage[table]{xcolor} % [table] : fournit aussi \rowcolors (alternance de lignes)
\usepackage{pagecolor}
\usepackage{tikz}
\usetikzlibrary{arrows.meta,positioning,calc,shapes.geometric,shapes.misc,shapes.arrows,decorations.pathmorphing,decorations.markings,patterns,shadows,mindmap,trees,fit,backgrounds,matrix,intersections}
\usepackage{pgfplots}
\usepackage[version=4]{mhchem} % équations nucléaires \ce{^{235}_{92}U}
\usepackage[european,siunitx]{circuitikz} % schémas électriques
\pgfplotsset{compat=1.18}
\usepgfplotslibrary{fillbetween} % aires entre courbes (calcul intégral)
\usepackage{enumitem}
\usepackage{fancyhdr}
% [most] charge toutes les bibliothèques tcolorbox (listings, minted, raster,
% documentation...) et gonfle inutilement la mémoire TeX sur un document long
% (~300 boîtes) : on ne charge que ce qui est réellement utilisé.
\usepackage[skins,breakable]{tcolorbox}
\usepackage{graphicx}
\usepackage{colortbl}
\usepackage{booktabs}
\usepackage{tabularx}
\usepackage{diagbox} % case d'en-tête diagonale des tableaux à double entrée
% Colonnes extensibles centrée / gauche / droite (C, L, R), souvent utilisées par les modèles
\newcolumntype{C}{>{\centering\arraybackslash}X}
\newcolumntype{L}{>{\raggedright\arraybackslash}X}
\newcolumntype{R}{>{\raggedleft\arraybackslash}X}
\usepackage{array}
\usepackage{multicol}
\usepackage{siunitx}
% parse-numbers=false : accepte aussi \SI{2,5 \times 10^{-3}}{...}
\sisetup{locale=FR,output-decimal-marker={,},group-minimum-digits=4,parse-numbers=false}
\usepackage{qrcode}
% \degree : commande fréquemment utilisée par les modèles pour un angle ou une
% température en dehors de \SI{}{} (non fournie par les paquets chargés).
\providecommand{\degree}{^{\circ}}
% \qed (fin de démonstration), utilisé par les modèles sans amsthm
\providecommand{\qed}{\hfill$\square$}
% \barre : double barre des valeurs interdites dans un tableau de variations (tkz-tab)
\providecommand{\barre}{\ensuremath{\|}}
% environnement proof (amsthm non chargé) utilisé par les modèles
\ifdefined\proof\else\newenvironment{proof}[1][Démonstration]{\par\noindent\textit{#1.}\ }{\qed\par}\fi
\usepackage{lastpage}
\usepackage{hyperref}
\hypersetup{hidelinks}

% ============================================================
% Palette « Pop » OnBuch+ — mêmes hex que la classe P (plus_ui.dart)
% ============================================================
\definecolor{poporange}{HTML}{F59321}
\definecolor{popdark}{HTML}{E07A0C}
\definecolor{popcream}{HTML}{FFF6E8}
\definecolor{popink}{HTML}{1C1714}
\definecolor{poppurple}{HTML}{7A5AE0}
\definecolor{popgreen}{HTML}{2E9E6A}
\definecolor{poppink}{HTML}{E0457B}
\definecolor{popblue}{HTML}{2F7DE1}
\definecolor{popgold}{HTML}{C98A12}
\definecolor{popmuted}{HTML}{8A7F76}
\definecolor{popline}{HTML}{EADFD2}
\definecolor{popgray}{HTML}{8A7F76}
\colorlet{popgrey}{popgray}
% Alias tolérants (le modèle nomme parfois les couleurs différemment)
\colorlet{popwhite}{white}
\colorlet{popblack}{popink}
\colorlet{popred}{poppink}
\colorlet{popyellow}{popgold}
% Teintes claires pour fonds de tableaux / zones
\colorlet{popblueL}{popblue!10!white}
\colorlet{poporangeL}{poporange!14!white}
\colorlet{popgreenL}{popgreen!12!white}
\colorlet{poppurpleL}{poppurple!12!white}
\colorlet{poppinkL}{poppink!12!white}
\colorlet{popgoldL}{popgold!14!white}

\pagecolor{popcream}

% ============================================================
% Typographie
% ============================================================
\defaultfontfeatures{Ligatures=TeX}
\setmainfont{PlusJakartaSans-Medium.ttf}[Path=../../../fonts/,BoldFont=PlusJakartaSans-Bold.ttf]
% Maths en sans-serif (Fira Math) avec chiffres et lettres droites de Plus
% Jakarta, pour que les formules se fondent dans le texte.
\usepackage[math-style=ISO,bold-style=ISO]{unicode-math}
\setmathfont{FiraMath-Regular.otf}[Path=../../../fonts/]
\setmathfont{PlusJakartaSans-Medium.ttf}[Path=../../../fonts/,range={up/num,up/{Latin,latin},\mathup}]
\usepackage{icomma} % virgule décimale sans espace en mode math
% Caractères mathématiques Unicode absents des polices de texte (Plus Jakarta,
% Archivo Black) : rendus par leur équivalent mathématique, en texte comme en
% formule — sinon ils s'affichent en carré vide (ex. « Arithmétique dans ℤ »).
\usepackage{newunicodechar}
\newunicodechar{ℝ}{\ensuremath{\mathbb{R}}}
\newunicodechar{ℤ}{\ensuremath{\mathbb{Z}}}
\newunicodechar{ℂ}{\ensuremath{\mathbb{C}}}
\newunicodechar{ℕ}{\ensuremath{\mathbb{N}}}
\newunicodechar{ℚ}{\ensuremath{\mathbb{Q}}}
\newunicodechar{𝔻}{\ensuremath{\mathbb{D}}}
\newunicodechar{α}{\ensuremath{\alpha}}
\newunicodechar{β}{\ensuremath{\beta}}
\newunicodechar{γ}{\ensuremath{\gamma}}
\newunicodechar{δ}{\ensuremath{\delta}}
\newunicodechar{ε}{\ensuremath{\varepsilon}}
\newunicodechar{θ}{\ensuremath{\theta}}
\newunicodechar{λ}{\ensuremath{\lambda}}
\newunicodechar{μ}{\ensuremath{\mu}}
\newunicodechar{σ}{\ensuremath{\sigma}}
\newunicodechar{τ}{\ensuremath{\tau}}
\newunicodechar{φ}{\ensuremath{\varphi}}
\newunicodechar{ψ}{\ensuremath{\psi}}
\newunicodechar{ω}{\ensuremath{\omega}}
\newunicodechar{Σ}{\ensuremath{\Sigma}}
% majuscules produites par \MakeUppercase dans les titres (α-aminés -> Α)
\newunicodechar{Α}{\ensuremath{\alpha}}
\newunicodechar{Β}{\ensuremath{\beta}}
\newunicodechar{Γ}{\ensuremath{\Gamma}}
\newunicodechar{Δ}{\ensuremath{\Delta}}
\newunicodechar{Ε}{\ensuremath{\varepsilon}}
\newunicodechar{Θ}{\ensuremath{\Theta}}
\newunicodechar{Λ}{\ensuremath{\Lambda}}
\newunicodechar{Μ}{\ensuremath{\mu}}
\newunicodechar{Τ}{\ensuremath{\tau}}
\newunicodechar{Φ}{\ensuremath{\Phi}}
\newunicodechar{Ψ}{\ensuremath{\Psi}}
\newunicodechar{Ω}{\ensuremath{\Omega}}
\newunicodechar{↦}{\ensuremath{\mapsto}}
\newunicodechar{∈}{\ensuremath{\in}}
\newunicodechar{∉}{\ensuremath{\notin}}
\newunicodechar{∧}{\ensuremath{\wedge}}
\newunicodechar{⇔}{\ensuremath{\Leftrightarrow}}
\newunicodechar{⟺}{\ensuremath{\Longleftrightarrow}}
\newunicodechar{⇒}{\ensuremath{\Rightarrow}}
\newunicodechar{⟹}{\ensuremath{\Longrightarrow}}
\newunicodechar{∘}{\ensuremath{\circ}}
\newunicodechar{∪}{\ensuremath{\cup}}
\newunicodechar{∩}{\ensuremath{\cap}}
\newunicodechar{≡}{\ensuremath{\equiv}}
\newunicodechar{⊂}{\ensuremath{\subset}}
\newunicodechar{⊥}{\ensuremath{\perp}}
\newunicodechar{∃}{\ensuremath{\exists}}
\newunicodechar{∀}{\ensuremath{\forall}}
\newunicodechar{∛}{\ensuremath{\sqrt[3]{\,}}}
\newunicodechar{∜}{\ensuremath{\sqrt[4]{\,}}}
\newunicodechar{⃗}{\ensuremath{{}^{\to}}}
\newunicodechar{ⁿ}{\ifmmode^{n}\else\textsuperscript{\ensuremath{n}}\fi}
\newunicodechar{ˣ}{\ifmmode^{x}\else\textsuperscript{\ensuremath{x}}\fi}
\newunicodechar{ᵇ}{\ifmmode^{b}\else\textsuperscript{\ensuremath{b}}\fi}
\newunicodechar{ʳ}{\ifmmode^{r}\else\textsuperscript{\ensuremath{r}}\fi}
\newunicodechar{ᵘ}{\ifmmode^{u}\else\textsuperscript{\ensuremath{u}}\fi}
\newunicodechar{ʸ}{\ifmmode^{y}\else\textsuperscript{\ensuremath{y}}\fi}
\newunicodechar{ᵃ}{\ifmmode^{a}\else\textsuperscript{\ensuremath{a}}\fi}
\newunicodechar{ᵉ}{\ifmmode^{e}\else\textsuperscript{\ensuremath{e}}\fi}
\newunicodechar{ᵏ}{\ifmmode^{k}\else\textsuperscript{\ensuremath{k}}\fi}
\newunicodechar{ᵖ}{\ifmmode^{p}\else\textsuperscript{\ensuremath{p}}\fi}
\newunicodechar{ᴺ}{\ifmmode^{N}\else\textsuperscript{\ensuremath{N}}\fi}
\newunicodechar{ᶜ}{\ifmmode^{c}\else\textsuperscript{\ensuremath{c}}\fi}
\newunicodechar{ʰ}{\ifmmode^{h}\else\textsuperscript{\ensuremath{h}}\fi}
\newunicodechar{ᵠ}{\ifmmode^{q}\else\textsuperscript{\ensuremath{q}}\fi}
\newunicodechar{ₙ}{\ifmmode_{n}\else\textsubscript{\ensuremath{n}}\fi}
\newunicodechar{ₐ}{\ifmmode_{a}\else\textsubscript{\ensuremath{a}}\fi}
\newunicodechar{ᵢ}{\ifmmode_{i}\else\textsubscript{\ensuremath{i}}\fi}
\newunicodechar{ᵣ}{\ifmmode_{r}\else\textsubscript{\ensuremath{r}}\fi}
\newunicodechar{ₖ}{\ifmmode_{k}\else\textsubscript{\ensuremath{k}}\fi}
\newunicodechar{ᵦ}{\ifmmode_{\beta}\else\textsubscript{\ensuremath{\beta}}\fi}
\newunicodechar{ₓ}{\ifmmode_{x}\else\textsubscript{\ensuremath{x}}\fi}
\newfontfamily\popdisplay{ArchivoBlack-Regular.ttf}[Path=../../../fonts/]
\newfontfamily\poplogo{SpaceGrotesk-Bold.ttf}[Path=../../../fonts/]
\newfontfamily\popsemi{PlusJakartaSans-SemiBold.ttf}[Path=../../../fonts/]
\newfontfamily\popxbold{PlusJakartaSans-ExtraBold.ttf}[Path=../../../fonts/]
\newfontfamily\popxboldls{PlusJakartaSans-ExtraBold.ttf}[Path=../../../fonts/,LetterSpace=3.0]

\setlist[enumerate]{leftmargin=1.7em,itemsep=5pt,topsep=4pt}
\setlist[itemize]{leftmargin=1.7em,itemsep=4pt,topsep=4pt,label={\color{poporange}\textbullet}}
\setlength{\parindent}{0pt}
\setlength{\parskip}{5pt}
\renewcommand{\arraystretch}{1.25}

% pgfplots : style Pop par défaut
\pgfplotsset{
  every axis/.append style={axis line style={popink,line width=1pt},tick style={popink},
    grid style={popline,line width=0.5pt},label style={font=\small},tick label style={font=\footnotesize}},
  popcurve/.style={poporange,line width=1.8pt,no markers,smooth},
}
% Même style disponible dans un \draw TikZ simple (hors pgfplots)
\tikzset{popcurve/.style={poporange,line width=1.8pt}}

% ============================================================
% Pill / squiggle
% ============================================================
\newcommand{\poppill}[3]{%
  \tikz[baseline=-0.55ex]\node[draw=popink,line width=1.2pt,rounded corners=2.6pt,%
    fill=#2,inner xsep=6.5pt,inner ysep=3pt]{%
    \popxboldls\fontsize{7}{7}\selectfont\color{#3}\MakeUppercase{#1}};%
}
\newcommand{\popsquiggle}[1]{%
  \tikz[baseline=-0.4ex]{
    \draw[#1,line width=2.6pt,line cap=round]
      plot [smooth] coordinates {(0,0.09) (0.34,0.01) (0.68,0.21) (1.02,0.01) (1.36,0.10)};
  }%
}
% Mot-clé surligné (marqueur orange sous le texte)
\newcommand{\cle}[1]{{\popxbold\color{popdark}#1}}

% ============================================================
% En-tête / pied
% ============================================================
\pagestyle{fancy}
\fancyhf{}
\renewcommand{\headrulewidth}{0pt}
\renewcommand{\footrulewidth}{0pt}
\lhead{\raisebox{-0.15ex}{\poplogo\fontsize{13}{13}\selectfont\color{popink}OnBuch\color{poporange}+}}
\rhead{\poppill{\DOCMATIERE\ \textbullet\ \DOCNIVEAU\ \textbullet\ Leçon \DOCLECON}{white}{popink}}
\lfoot{\popxbold\fontsize{7.6}{7.6}\selectfont\color{popmuted}\MakeUppercase{Cours signé OnBuch+}\ \textbullet\ {\popsemi\DOCTITRE}}
\rfoot{\tikz[baseline=-0.6ex]\node[rounded corners=8.5pt,fill=popink,minimum height=17pt,minimum width=17pt,inner xsep=5pt,inner sep=0]{\popxbold\fontsize{8}{8}\selectfont\color{popcream}\thepage\,/\,\pageref*{LastPage}};}

% ============================================================
% Titres
% ============================================================
\newcommand{\chaptitre}[2]{%
  \begin{tcolorbox}[enhanced,colback=poporange,colframe=popink,boxrule=2.6pt,arc=11pt,
      drop shadow=popink,left=15pt,right=15pt,top=12pt,bottom=11pt,before skip=3pt,after skip=18pt]
    \poppill{#2}{popink}{popcream}\par\vspace{9pt}
    {\raggedright\hyphenpenalty=10000\exhyphenpenalty=10000\popdisplay\fontsize{22}{24}\selectfont\color{popcream}\MakeUppercase{#1}\par}
  \end{tcolorbox}
}
% Section numérotée : pastille orange avec le numéro + titre Archivo
\newcounter{popsec}
\newcommand{\coursec}[1]{%
  \stepcounter{popsec}\setcounter{popsub}{0}%
  \par\vspace{16pt}\noindent
  \begin{minipage}{\linewidth}
  \tikz[baseline=-0.7ex]\node[draw=popink,line width=1.6pt,fill=poporange,rounded corners=4pt,minimum size=22pt,inner sep=0]{\popdisplay\fontsize{12}{12}\selectfont\color{popcream}\thepopsec};%
  \hspace{8pt}{\popdisplay\fontsize{15}{17}\selectfont\color{popink}\MakeUppercase{#1}}\par
  \vspace{4pt}\noindent{\color{poporange}\rule{36pt}{3.6pt}}
  \end{minipage}\par\nopagebreak\vspace{8pt}
}
\newcounter{popsub}
\newcommand{\courssub}[1]{%
  \stepcounter{popsub}%
  \par\vspace{9pt}\noindent{\popxbold\fontsize{11.5}{13}\selectfont\color{popink}{\color{poporange}\thepopsec.\thepopsub}\hspace{6pt}#1}\par\nopagebreak\vspace{4pt}
}

% ============================================================
% Boîtes pédagogiques — même recette : fond blanc, bordure épaisse colorée,
% ombre dure encre, badge plein. Titre optionnel [#1].
% ============================================================
\tcbset{popbase/.style={breakable,enhanced,colback=white,boxrule=2.3pt,arc=9pt,
  shadow={3pt}{-3pt}{0pt}{popink},left=14pt,right=14pt,top=11pt,bottom=11pt,before skip=11pt,after skip=15pt,
  fontupper=\popsemi\fontsize{10.6}{14.5}\selectfont\color{popink},
  fontlower=\popsemi\fontsize{10.6}{14.5}\selectfont\color{popink}}}
% Coche dessinée (le glyphe ✓ n'existe pas dans les polices du document)
\newcommand{\popcheck}[1]{\tikz[baseline=-0.5ex]\draw[#1,line width=1.6pt,line cap=round,line join=round](-0.13,0.0)--(-0.04,-0.09)--(0.13,0.1);}
\providecommand{\coche}{\popcheck{popgreen}}
\newcommand{\popboxhead}[3]{\poppill{#1}{#2}{white}\ifx\relax#3\relax\else\hspace{6pt}{\popxbold\fontsize{10.6}{12}\selectfont\color{popink}#3}\fi\par\vspace{7pt}}

\newtcolorbox{aretenir}[1][]{popbase,colframe=popgreen,before upper={\popboxhead{À retenir}{popgreen}{#1}}}
\newtcolorbox{exemplebox}[1][]{popbase,colframe=poporange,before upper={\popboxhead{Exemple}{poporange}{#1}}}
\newtcolorbox{definition}[1][]{popbase,colframe=popblue,before upper={\popboxhead{Définition}{popblue}{#1}}}
\newtcolorbox{propriete}[1][]{popbase,colframe=poppurple,before upper={\popboxhead{Propriété}{poppurple}{#1}}}
\newtcolorbox{methode}[1][]{popbase,colframe=popgold,before upper={\popboxhead{Méthode}{popgold}{#1}}}
\newtcolorbox{attention}[1][]{popbase,colframe=poppink,before upper={\popboxhead{Attention}{poppink}{#1}}}
\newtcolorbox{experience}[1][]{popbase,colframe=popblue,colback=popblueL,before upper={\popboxhead{Expérience}{popblue}{#1}}}
% « Le savais-tu ? » : carte violette pleine (contexte, culture, Cameroun)
\newtcolorbox{savaistu}[1][]{popbase,colback=poppurple,colframe=popink,
  fontupper=\popsemi\fontsize{10.4}{14.2}\selectfont\color{white},
  before upper={\poppill{Le savais-tu\,?}{white}{poppurple}\ifx\relax#1\relax\else\hspace{6pt}{\popxbold\color{white}#1}\fi\par\vspace{7pt}}}
% Exercice résolu : énoncé en haut, \tcblower puis la solution sur fond crème
\newcounter{popexo}\newcounter{popexr}
\newtcolorbox{exoresolu}[1][]{popbase,colframe=poporange,colbacklower=popcream,
  segmentation style={popink,line width=1.2pt,dashed},
  before upper={\stepcounter{popexr}\popboxhead{Exercice résolu \thepopexr}{poporange}{#1}},
  before lower={\poppill{Solution}{popink}{popcream}\par\vspace{6pt}}}
% Exercice (non résolu en place) — la correction est dans la partie Corrigés
\newtcolorbox{exercice}[1][]{popbase,colframe=popink,
  before upper={\stepcounter{popexo}\popboxhead{Exercice \thepopexo}{popink}{#1}}}
% Corrigé
\newtcolorbox{corrige}[1][]{popbase,colframe=popgreen,colback=popgreenL,
  before upper={\popboxhead{Corrigé}{popgreen}{#1}}}
% Objectifs / prérequis / bilan
\newtcolorbox{objectifs}{popbase,colframe=popink,colback=poporangeL,
  before upper={\popboxhead{Objectifs de la leçon}{popink}{}}}
\newtcolorbox{prerequis}{popbase,colframe=popink,colback=popblueL,
  before upper={\popboxhead{Prérequis}{popblue}{}}}
\newtcolorbox{bilan}{popbase,colframe=popink,colback=white,boxrule=2.8pt,
  before upper={\popboxhead{Fiche bilan}{poporange}{L'essentiel à connaître pour l'examen}}}
% Figure encadrée avec légende
\newtcolorbox{popfigure}[1][]{enhanced,breakable=false,colback=white,colframe=popline,boxrule=1.4pt,arc=8pt,
  left=10pt,right=10pt,top=10pt,bottom=8pt,before skip=10pt,after skip=12pt,halign=center,
  lower separated=false,halign lower=center,
  fontlower=\popsemi\fontsize{9}{11.5}\selectfont\color{popmuted},
  #1}
\newcommand{\legende}[1]{\par\vspace{6pt}{\popsemi\fontsize{9}{11.5}\selectfont\color{popmuted}#1}}

% ============================================================
% Signature OnBuch+
% ============================================================
\newcommand{\poplogoinline}[1]{{\poplogo\fontsize{#1}{#1}\selectfont\color{popink}OnBuch\color{poporange}+}}
\newcommand{\signatureonbuch}{%
  \begin{tcolorbox}[enhanced,colback=popink,colframe=popink,boxrule=2.2pt,arc=10pt,
      drop shadow=poporange,left=14pt,right=14pt,top=10pt,bottom=10pt,before skip=0pt,after skip=0pt]
    \tikz[baseline=-0.7ex]\node[circle,fill=popgreen,draw=popcream,line width=1.3pt,minimum size=22pt,inner sep=0]{\popcheck{white}};%
    \hspace{9pt}\begin{minipage}[c]{0.62\linewidth}
      {\popxbold\fontsize{12}{13}\selectfont\color{popcream}Signé\ \textbullet\ {\poplogo OnBuch\color{poporange}+}}\par\vspace{2pt}
      {\raggedright\popsemi\fontsize{8}{10.5}\selectfont\color{popline}\MakeUppercase{Équipe pédagogique OnBuch+}\\\MakeUppercase{Conforme au programme officiel MINESEC}\par}
    \end{minipage}\hfill
    {\poplogo\fontsize{22}{22}\selectfont\color{popcream}OnBuch\color{poporange}+}
  \end{tcolorbox}%
}

% ============================================================
% Page de garde
% #1 titre · #2 matière · #3 niveau · #4 module · #5 n° leçon · #6 durée ·
% #7 description / objectifs (texte ou itemize)
% ============================================================
\newcommand{\pagegarde}[7]{%
  \thispagestyle{empty}
  \newgeometry{a4paper,margin=1.7cm,top=1.6cm,bottom=1.4cm}%
  \begin{tikzpicture}[remember picture,overlay]
    \fill[poporange!18] ([xshift=-3.2cm,yshift=-3.6cm]current page.north east) circle (4.6cm);
    \fill[poppurple!14] ([xshift=2.4cm,yshift=7.6cm]current page.south west) circle (3.8cm);
  \end{tikzpicture}%
  {\poplogo\fontsize{30}{30}\selectfont\color{popink}OnBuch\color{poporange}+}\hfill
  \raisebox{0.6ex}{\poppill{Cours signé \textbullet\ Premium}{popink}{popcream}}\par
  \vspace{1pt}\popsquiggle{poppurple}\par
  \vspace{8pt}
  \begin{tcolorbox}[enhanced,colback=poporange,colframe=popink,boxrule=2.8pt,arc=13pt,
      drop shadow=popink,left=18pt,right=18pt,top=12pt,bottom=13pt,after skip=10pt]
    \poppill{#2 \textbullet\ #3}{popink}{popcream}\hspace{5pt}\poppill{#4}{white}{popink}\par\vspace{8pt}
    {\popdisplay\fontsize{14}{14}\selectfont\color{popink}LEÇON #5}\par\vspace{4pt}
    {\raggedright\hyphenpenalty=10000\exhyphenpenalty=10000\popdisplay\fontsize{25}{27}\selectfont\color{popcream}\MakeUppercase{#1}\par}
    \vspace{8pt}
    {\popxbold\fontsize{9.5}{9.5}\selectfont\color{popink}\MakeUppercase{Durée conseillée : #6}}
  \end{tcolorbox}
  \begin{tcolorbox}[enhanced,colback=white,colframe=popink,boxrule=2.2pt,arc=10pt,
      drop shadow=popink,left=16pt,right=16pt,top=10pt,bottom=10pt,after skip=8pt]
    \poppill{Dans cette leçon}{poppurple}{white}\par\vspace{5pt}
    \popsemi\fontsize{9.3}{12.6}\selectfont\color{popink}#7
  \end{tcolorbox}
  \noindent\begin{minipage}[t]{0.487\linewidth}
    \begin{tcolorbox}[enhanced,colback=popgreen,colframe=popink,boxrule=2.2pt,arc=10pt,
        drop shadow=popink,left=11pt,right=11pt,top=10pt,bottom=10pt,height=3.1cm,valign=center]
      \tcbox[colback=white,colframe=popink,boxrule=1.3pt,arc=5pt,boxsep=0pt,nobeforeafter,
        left=3pt,right=3pt,top=3pt,bottom=3pt,valign=center]{\qrcode[height=1.9cm]{https://play.google.com/store/apps/details?id=com.luvvix.onbuch&hl=fr}}%
      \hspace{8pt}\begin{minipage}[c]{0.5\linewidth}
        {\popdisplay\fontsize{11}{12.5}\selectfont\color{white}\MakeUppercase{Télécharge OnBuch}}\par\vspace{4pt}
        {\raggedright\popsemi\fontsize{8.4}{11}\selectfont\color{white}Gratuit \textbullet\ Google Play\\Tuteur IA, annales, concours, résultats\par}
      \end{minipage}
    \end{tcolorbox}
  \end{minipage}\hfill
  \begin{minipage}[t]{0.487\linewidth}
    \begin{tcolorbox}[enhanced,colback=poppurple,colframe=popink,boxrule=2.2pt,arc=10pt,
        drop shadow=popink,left=11pt,right=11pt,top=10pt,bottom=10pt,height=3.1cm,valign=center]
      \tikz[baseline=-0.7ex]\node[circle,fill=white,draw=popink,line width=1.3pt,minimum size=1.6cm,inner sep=0]{\popdisplay\fontsize{24}{24}\selectfont\color{poppurple}+};%
      \hspace{8pt}\begin{minipage}[c]{0.6\linewidth}
        {\popdisplay\fontsize{11}{12.5}\selectfont\color{white}\MakeUppercase{Passe à OnBuch+}}\par\vspace{4pt}
        {\raggedright\popsemi\fontsize{8.4}{11}\selectfont\color{white}Tous les cours signés, Tuteur IA illimité, fiches PDF — onglet OnBuch+ de l'app\par}
      \end{minipage}
    \end{tcolorbox}
  \end{minipage}
  \vspace{8pt}
  \popcontenu
  \vfill
  \signatureonbuch
  \restoregeometry
  \newpage
}

% Rangée « Ce cours contient » (page de garde)
\newcommand{\poptile}[3]{% #1 couleur #2 gros texte #3 légende
  \begin{tcolorbox}[enhanced,colback=#1,colframe=popink,boxrule=2pt,arc=9pt,shadow={2.5pt}{-2.5pt}{0pt}{popink},
      left=6pt,right=6pt,top=9pt,bottom=9pt,halign=center,valign=center,height=1.75cm,width=\linewidth,nobeforeafter]
    {\popdisplay\fontsize{12.5}{13}\selectfont\color{white}\MakeUppercase{#2}}\par\vspace{3pt}
    {\centering\popsemi\fontsize{7.8}{9.5}\selectfont\color{white}#3\par}
  \end{tcolorbox}}
\newcommand{\popcontenu}{%
  \par\noindent{\popxboldls\fontsize{7.5}{7.5}\selectfont\color{popmuted}CE COURS SIGNÉ CONTIENT}\par\vspace{7pt}
  \noindent\begin{minipage}[t]{0.235\linewidth}\poptile{popblue}{Cours}{complet, illustré, pas à pas}\end{minipage}\hfill
  \begin{minipage}[t]{0.235\linewidth}\poptile{poporange}{Méthodes}{et exercices résolus}\end{minipage}\hfill
  \begin{minipage}[t]{0.235\linewidth}\poptile{poppink}{Exercices}{type examen + corrigés}\end{minipage}\hfill
  \begin{minipage}[t]{0.235\linewidth}\poptile{popgreen}{Fiche bilan}{l'essentiel à retenir}\end{minipage}\par}

% Sommaire
\newtcolorbox{sommaire}{enhanced,colback=white,colframe=popink,boxrule=2.2pt,arc=10pt,
  drop shadow=popink,left=18pt,right=18pt,top=13pt,bottom=13pt,before skip=6pt,after skip=20pt,
  before upper={\poppill{Sommaire}{poppurple}{white}\par\vspace{9pt}\popsemi\fontsize{10.6}{14.5}\selectfont\color{popink}}}

% Signature de fin de cours — poussée tout en bas de la dernière page
\newcommand{\findecours}{%
  \par\vfill\nopagebreak
  \begin{center}\popsquiggle{poporange}\end{center}\vspace{4pt}
  \signatureonbuch
}

%% FIGFIX-BEGIN v5 — correctifs globaux des figures (docs/onbuch-generation/tools/figfix)
\usepackage{adjustbox}\usepackage{etoolbox}\tcbuselibrary{raster}
% toute figure TikZ/pgfplots plus large que la ligne est réduite à la largeur disponible
\BeforeBeginEnvironment{tikzpicture}{\begin{adjustbox}{max width=\linewidth}}
\AfterEndEnvironment{tikzpicture}{\end{adjustbox}}
% légendes pgfplots lisibles par-dessus les courbes
\pgfplotsset{every axis/.append style={legend style={fill=white,fill opacity=0.92,text opacity=1,draw=black!15,inner sep=3pt}}}
% étiquettes d'arêtes (syntaxe edge["..."]) avec fond blanc
\tikzset{every edge quotes/.append style={fill=white,inner sep=1.2pt}}
%% FIGFIX-END
\begin{document}

\pagegarde{Lois de Newton et conservation de l'énergie mécanique}{Physique}{Tle D}{Module 2 : Mouvements et interactions — évolutions temporelles des systèmes mécaniques}{P05}{8 h}{Cette leçon établit les lois fondamentales de la mécanique : les trois lois de Newton appliquées au point matériel et au solide en rotation, puis l'approche énergétique fondée sur le théorème de l'énergie cinétique et la conservation de l'énergie mécanique. Elle apprend à choisir la méthode la plus efficace pour résoudre un problème de mécanique, compétence clé de l'épreuve de physique au Baccalauréat C.
\vspace{4pt}
\begin{multicols}{2}\small
\begin{itemize}
\item À la fin de cette leçon, je sais caractériser un mouvement dans un repère cartésien et dans le repère de Frenet (accélérations tangentielle et normale)
\item À la fin de cette leçon, je sais reconnaître un référentiel galiléen et énoncer le principe d'inertie
\item À la fin de cette leçon, je sais appliquer le théorème du centre d'inertie ΣF = m·a\_G à un point matériel ou à un solide en translation
\item À la fin de cette leçon, je sais énoncer et exploiter la troisième loi de Newton (principe des actions réciproques)
\item À la fin de cette leçon, je sais mettre en équation le mouvement d'un solide en rotation autour d'un axe fixe grâce à ΣM = J·θ''
\item À la fin de cette leçon, je sais appliquer le théorème de l'énergie cinétique en translation et en rotation
\end{itemize}\end{multicols}}

\begin{sommaire}
\begin{multicols}{2}
\begin{enumerate}
\item Rappels de cinématique et repère de Frenet
\item Référentiels galiléens et première loi de Newton
\item Deuxième et troisième lois de Newton
\item Dynamique du solide en rotation autour d'un axe fixe
\item Théorème de l'énergie cinétique
\item Énergie potentielle, énergie mécanique et conservation
\item Activité d'intégration
\item Méthodes et exercices résolus
\item Exercices
\item Corrigés des exercices
\item Fiche bilan
\end{enumerate}
\end{multicols}
\end{sommaire}

\begin{objectifs}
\begin{itemize}
\item À la fin de cette leçon, je sais caractériser un mouvement dans un repère cartésien et dans le repère de Frenet (accélérations tangentielle et normale)
\item À la fin de cette leçon, je sais reconnaître un référentiel galiléen et énoncer le principe d'inertie
\item À la fin de cette leçon, je sais appliquer le théorème du centre d'inertie ΣF = m·a\_G à un point matériel ou à un solide en translation
\item À la fin de cette leçon, je sais énoncer et exploiter la troisième loi de Newton (principe des actions réciproques)
\item À la fin de cette leçon, je sais mettre en équation le mouvement d'un solide en rotation autour d'un axe fixe grâce à ΣM = J·θ''
\item À la fin de cette leçon, je sais appliquer le théorème de l'énergie cinétique en translation et en rotation
\item À la fin de cette leçon, je sais utiliser la conservation de l'énergie mécanique pour déterminer une vitesse ou une position
\item À la fin de cette leçon, je sais détecter la non-conservation de l'énergie mécanique en présence de frottements et calculer l'énergie dissipée
\item À la fin de cette leçon, je sais choisir entre la méthode dynamique (lois de Newton) et la méthode énergétique selon la question posée
\end{itemize}
\end{objectifs}

\begin{prerequis}
\begin{itemize}
\item Vecteurs : composantes, norme, somme vectorielle (mathématiques et physique de Première)
\item Notions de force, masse, poids P = m·g (classes antérieures)
\item Dérivation d'une fonction et équations différentielles simples x'' = constante (mathématiques Tle)
\item Travail d'une force W(F) = F·AB·cos α (Première C)
\item Mouvement rectiligne uniforme et uniformément varié (Première C)
\end{itemize}
\end{prerequis}

\newpage

\coursec{Rappels de cinématique et repère de Frenet}

Un moto-taxi (benskin) aborde un virage relevé à la sortie de Yaoundé : pourquoi peut-il négocier le virage plus vite sans déraper ? Pourquoi le moteur d'un poids lourd grimpant la côte de Mbanga chauffe-t-il alors que sa vitesse reste constante ? Et comment un lanceur de marteau parvient-il à communiquer une énorme énergie à son engin en tournant sur lui-même ? Ces situations, comme le mouvement d'une pirogue sur le Wouri ou d'une toupie sur le sable, s'expliquent par trois lois universelles énoncées par Newton et par un principe puissant : la conservation de l'énergie mécanique. Cette leçon fournit les deux outils — dynamique et énergétique — pour analyser et prévoir le mouvement de tout objet de notre environnement.

Mais avant de manier les forces et les énergies, il faut un langage précis pour \emph{décrire} le mouvement : c'est la cinématique. Cette première section rassemble et approfondit les outils vus en Première, et introduit un outil nouveau et décisif : le repère de Frenet, qui permet de comprendre pourquoi un virage, même parcouru à vitesse constante, est un mouvement \emph{accéléré}.

\courssub{Système, référentiel et relativité du mouvement}

\textbf{Le système étudié.} En mécanique, on commence toujours par délimiter ce que l'on étudie : le \cle{système}. C'est un corps (ou un ensemble de corps) que l'on isole \emph{par la pensée} du reste de l'univers. Tout ce qui n'appartient pas au système constitue le \cle{milieu extérieur}. Dans cette leçon, le système sera le plus souvent modélisé par un \cle{point matériel} (tout le corps est réduit à son centre d'inertie $G$, affecté de la masse totale $m$) ou par un \cle{solide} indéformable.

\textbf{Le référentiel.} Dire qu'un objet « bouge » n'a de sens que par rapport à autre chose. Assis dans le car de l'agence de voyage Douala--Bafoussam, tu es immobile par rapport à ton voisin de siège, mais tu files à \SI{80}{km/h} par rapport à la route. Le mouvement est \emph{relatif} : il dépend du solide de référence choisi.

\begin{definition}[Référentiel et repères]
Un \cle{référentiel} est un solide de référence par rapport auquel on décrit le mouvement d'un système. Pour décrire quantitativement ce mouvement, on associe au référentiel :
\begin{itemize}
\item un \cle{repère d'espace} $(O\,;\vec{i},\vec{j},\vec{k})$, permettant de repérer la position du mobile par ses coordonnées $(x,y,z)$ ;
\item un \cle{repère de temps} : une horloge et une origine des dates $t_0 = 0$.
\end{itemize}
Référentiels usuels : le référentiel \emph{terrestre} (lié au sol, pour les mouvements quotidiens), le référentiel \emph{géocentrique} (lié au centre de la Terre, pour les satellites), le référentiel \emph{héliocentrique} (lié au Soleil, pour les planètes).
\end{definition}

\begin{attention}[La première phrase de tout exercice de mécanique]
Toute étude de mouvement commence par la déclaration du système et du référentiel : « Système étudié : la moto, assimilée à un point matériel $G$ ; référentiel : terrestre, supposé galiléen ». Oublier cette précision est une erreur de rédaction sanctionnée au Baccalauréat, car les lois de Newton ne s'appliquent que dans certains référentiels (voir section 2).
\end{attention}

\courssub{Vecteurs position, vitesse et accélération}

\begin{definition}[Vecteurs cinématiques]
Dans un repère $(O\,;\vec{i},\vec{j},\vec{k})$, la position du point $M$ à la date $t$ est donnée par le \cle{vecteur position} :
\[ \overrightarrow{OM}(t) = x(t)\,\vec{i} + y(t)\,\vec{j} + z(t)\,\vec{k} \]
Le \cle{vecteur vitesse} est la dérivée du vecteur position par rapport au temps :
\[ \vec{v}(t) = \dfrac{d\overrightarrow{OM}}{dt} = \dot{x}\,\vec{i} + \dot{y}\,\vec{j} + \dot{z}\,\vec{k} \]
Le \cle{vecteur accélération} est la dérivée du vecteur vitesse :
\[ \vec{a}(t) = \dfrac{d\vec{v}}{dt} = \ddot{x}\,\vec{i} + \ddot{y}\,\vec{j} + \ddot{z}\,\vec{k} \]
Unités SI : position en mètre (m), vitesse en \si{m.s^{-1}}, accélération en \si{m.s^{-2}}.
\end{definition}

Quelques propriétés géométriques essentielles, à retenir car elles guident tous les schémas :
\begin{itemize}
\item le vecteur vitesse $\vec{v}$ est toujours \cle{tangent à la trajectoire} et orienté dans le sens du mouvement ;
\item le vecteur accélération $\vec{a}$ est toujours dirigé vers l'\emph{intérieur de la concavité} de la trajectoire ;
\item si la trajectoire est rectiligne, $\vec{v}$ et $\vec{a}$ sont colinéaires à cette droite.
\end{itemize}

Vérifions l'homogénéité par \cle{analyse dimensionnelle} : $[v] = L \cdot T^{-1}$ car on dérive une longueur par rapport au temps ; $[a] = L \cdot T^{-2}$. La relation $a = \dfrac{dv}{dt}$ est donc dimensionnellement cohérente.

\courssub{Mouvements rectilignes : MRU et MRUV}

Quand la trajectoire est une droite, un seul axe $(Ox)$ suffit. On note $x(t)$ l'abscisse, $v = \dot{x}$ et $a = \ddot{x}$.

\begin{propriete}[Équations horaires du MRU et du MRUV]
\textbf{Mouvement rectiligne uniforme (MRU)} : $a = 0$, la vitesse est constante :
\[ x(t) = v\,t + x_0 \]
\textbf{Mouvement rectiligne uniformément varié (MRUV)} : $a = a_0$ constante :
\[ v(t) = a_0\,t + v_0 \qquad ; \qquad x(t) = \tfrac{1}{2}\,a_0\,t^2 + v_0\,t + x_0 \]
On en déduit la relation indépendante du temps, très utile en exercice :
\[ v^2 - v_0^2 = 2\,a_0\,(x - x_0) \]
\end{propriete}

Ces relations s'obtiennent par intégrations successives : $a$ constante $\Rightarrow$ $v = \int a\,dt = a_0 t + v_0$ $\Rightarrow$ $x = \int v\,dt = \tfrac{1}{2}a_0 t^2 + v_0 t + x_0$, les constantes étant fixées par les conditions initiales.

\begin{exemplebox}[Freinage d'un benskin]
Un moto-taxi roule à $v_0 = \SI{54}{km/h} = \SI{15}{m.s^{-1}}$ et freine avec une décélération constante $a_0 = -\SI{3,0}{m.s^{-2}}$ jusqu'à l'arrêt. Distance de freinage : la relation $v^2 - v_0^2 = 2a_0(x - x_0)$ avec $v = 0$ donne
\[ x - x_0 = \dfrac{-v_0^2}{2a_0} = \dfrac{-15^2}{2\times(-3,0)} = \SI{37,5}{m} \]
soit près de \SI{38}{m} : plus d'une quinzaine de fois la longueur de la moto !
\end{exemplebox}

\begin{popfigure}
\begin{center}
\begin{tikzpicture}
\begin{axis}[width=0.9\linewidth,height=5.2cm, xlabel={$t$ (s)}, ylabel={}, xmin=0, xmax=5.5, ymin=-4, ymax=20, legend pos=north east, grid=both, grid style={popline!40}, legend style={font=\small}]
\addplot[popcurve,domain=0:5,samples=100,color=popblue]{-1.5*x^2+15*x};
\addlegendentry{$x(t)$ (m)}
\addplot[popcurve,domain=0:5,samples=50,color=poporange]{-3*x+15};
\addlegendentry{$v(t)$ (\si{m.s^{-1}})}
\addplot[popcurve,domain=0:5,samples=2,color=popgreen]{-3};
\addlegendentry{$a(t)$ (\si{m.s^{-2}})}
\end{axis}
\end{tikzpicture}
\end{center}
\legende{Équations horaires du freinage du benskin (MRUV) : $x(t) = -\num{1,5}t^2 + 15t$, $v(t) = -3t + 15$, $a(t) = -3$. La parabole $x(t)$ atteint son maximum quand $v$ s'annule ($t = \SI{5}{s}$) : la moto s'arrête.}
\end{popfigure}

\courssub{Mouvement circulaire : grandeurs angulaires}

Pour un point $M$ décrivant un cercle de centre $O$ et de rayon $R$, la position la plus naturelle n'est pas $(x,y)$ mais l'\cle{abscisse angulaire} $\theta(t)$, angle entre l'axe $(Ox)$ et $\overrightarrow{OM}$, mesuré en radians (rad).

\begin{definition}[Grandeurs angulaires]
La \cle{vitesse angulaire} est $\omega = \dot{\theta} = \dfrac{d\theta}{dt}$ (en \si{rad.s^{-1}}) ; l'\cle{accélération angulaire} est $\ddot{\theta} = \dfrac{d\omega}{dt}$ (en \si{rad.s^{-2}}). La vitesse linéaire (norme de $\vec{v}$) est reliée à la vitesse angulaire par :
\[ v = R\,\omega \]
\end{definition}

Cette relation se comprend intuitivement : en une période $T$, le point parcourt le périmètre $2\pi R$ pendant que l'angle balaye $2\pi$ rad, donc $v = \dfrac{2\pi R}{T} = R\,\omega$. Analyse dimensionnelle : $R\omega$ s'exprime en $\si{m}\times\si{rad.s^{-1}} = \si{m.s^{-1}}$ (le radian est sans dimension) : la formule est homogène.

\begin{attention}[Le radian est obligatoire]
La relation $v = R\omega$ n'est valable que si $\omega$ est en \si{rad.s^{-1}}. Convertis toujours les tours par minute : $\omega\,(\si{rad.s^{-1}}) = N\,(\si{tr/min}) \times \dfrac{2\pi}{60}$.
\end{attention}

\courssub{Le repère de Frenet : l'outil des trajectoires courbes}

Pour une trajectoire quelconque (virage, parabole d'un ballon, spirale), le repère cartésien fixe est mal adapté. L'idée de Frenet : attacher au mobile un repère \emph{mobile} qui « suit » la trajectoire.

\begin{definition}[Repère de Frenet et composantes de l'accélération]
En tout point $M$ d'une trajectoire orientée, on définit la \cle{base de Frenet} $(\vec{T},\vec{N})$ :
\begin{itemize}
\item $\vec{T}$ : vecteur unitaire \emph{tangent} à la trajectoire, orienté dans le sens du mouvement ;
\item $\vec{N}$ : vecteur unitaire \emph{normal} (perpendiculaire à $\vec{T}$), orienté vers le \emph{centre de courbure} (l'intérieur du virage).
\end{itemize}
Dans cette base, le vecteur vitesse s'écrit $\vec{v} = v\,\vec{T}$ et l'accélération se décompose en :
\[ \vec{a} = a_T\,\vec{T} + a_N\,\vec{N} \qquad \text{avec} \qquad a_T = \dfrac{dv}{dt} \quad \text{et} \quad a_N = \dfrac{v^2}{\rho} \]
où $\rho$ est le \cle{rayon de courbure} de la trajectoire en $M$ (pour un cercle, $\rho = R$). La norme de l'accélération est $a = \sqrt{a_T^2 + a_N^2}$.
\end{definition}

\textbf{Justification qualitative.} Le vecteur accélération mesure la variation du vecteur vitesse. Or $\vec{v}$ peut varier de deux façons indépendantes :
\begin{itemize}
\item sa \emph{norme} peut changer (on accélère ou on freine) : c'est la composante tangentielle $a_T = \dfrac{dv}{dt}$, qui ne dépend que de la variation de la valeur de la vitesse ;
\item sa \emph{direction} peut changer (on tourne) : c'est la composante normale $a_N = \dfrac{v^2}{\rho}$, d'autant plus grande que la vitesse est élevée et que le virage est serré ($\rho$ petit).
\end{itemize}

\begin{popfigure}
\begin{center}
\begin{tikzpicture}[scale=1.2, >=Stealth]
% Trajectoire : une arche symétrique (concavité vers le bas)
\draw[popline!60, thick] (0,0) .. controls (1.5, 2.5) and (3.5, 2.5) .. (5, 0);
% Point M au sommet de l'arche (environ)
\coordinate (M) at (2.5, 2.0);
\fill[popink] (M) circle (2.2pt) node[above, popdark]{$M$};
% Vecteur tangent T (vers la droite, légèrement descendant)
\draw[->, popblue, thick] (M) -- ++(1.5, -0.3) node[midway, above, popblue]{$\vec{T}$};
\draw[->, popblue, very thick] (M) -- ++(2.0, -0.4) node[right, popblue]{$\vec{v}=v\,\vec{T}$};
% Vecteur normal N (vers le bas, centre de courbure)
\draw[->, poporange, thick] (M) -- ++(0.2, -1.0) node[below, poporange]{$\vec{N}$};
% Composante tangentielle a_T (dans le sens de T, mouvement accéléré)
\draw[->, popgreen, very thick] (M) -- ++(1.2, -0.24) node[below right, popgreen]{$a_T\,\vec{T}$};
% Composante normale a_N (dans le sens de N)
\draw[->, poppurple, very thick] (M) -- ++(0.15, -0.75) node[below, poppurple]{$a_N\,\vec{N}$};
% Vecteur accélération résultant
\draw[->, popdark, very thick] (M) -- ++(1.35, -0.99) node[below right, popdark]{$\vec{a}$};
% Lignes de construction (parallélogramme)
\draw[popmuted, dashed] (M) ++(1.2, -0.24) -- ++(0.15, -0.75);
\draw[popmuted, dashed] (M) ++(0.15, -0.75) -- ++(1.2, -0.24);
\node[popmuted, font=\small] at (2.5, -0.6) {trajectoire};
\end{tikzpicture}
\end{center}
\legende{Base de Frenet en un point $M$ d'une trajectoire curviligne : $\vec{T}$ tangent dans le sens du mouvement, $\vec{N}$ vers le centre de courbure. L'accélération $\vec{a}$ est la somme de $a_T\vec{T}$ (variation de la norme de la vitesse) et $a_N\vec{N}$ (variation de sa direction).}
\end{popfigure}

\courssub{Cas du mouvement circulaire uniforme (MCU)}

\begin{propriete}[Accélération du mouvement circulaire uniforme]
Dans un mouvement circulaire \emph{uniforme} (vitesse angulaire $\omega$ constante, rayon $R$), la norme de la vitesse est constante donc $a_T = \dfrac{dv}{dt} = 0$ : l'accélération est \emph{purement normale}, dirigée vers le centre (on dit \cle{centripète}) :
\[ a = a_N = \dfrac{v^2}{R} = R\,\omega^2 \]
L'égalité des deux expressions vient de $v = R\omega$. Cette propriété est admise en Terminale ; elle se démontre en dérivant deux fois le vecteur position $\overrightarrow{OM} = R\cos(\omega t)\,\vec{i} + R\sin(\omega t)\,\vec{j}$, ce qui donne $\vec{a} = -\omega^2\,\overrightarrow{OM}$ : $\vec{a}$ est bien colinéaire à $\vec{N}$, dirigée vers $O$, de norme $R\omega^2$.
\end{propriete}

\begin{attention}[Vitesse constante $\neq$ accélération nulle]
C'est l'erreur la plus fréquente ! « Uniforme » signifie que la \emph{norme} de la vitesse est constante, pas le \emph{vecteur} vitesse. Dans un MCU, la direction de $\vec{v}$ change en permanence, donc $\vec{a} \neq \vec{0}$ : $a = \dfrac{v^2}{R}$. Une voiture qui prend un rond-point à \SI{30}{km/h} constants est bel et bien \emph{accélérée} — c'est d'ailleurs pour cela que les passagers se sentent poussés vers l'extérieur.
\end{attention}

\begin{exoresolu}[Moto dans un virage à la sortie de Yaoundé]
Une moto négocie un virage circulaire de rayon $R = \SI{50}{m}$ à la vitesse constante $v = \SI{72}{km/h}$. Calculer son accélération.
\tcblower
\textbf{Conversion.} $v = \dfrac{72}{3,6} = \SI{20}{m.s^{-1}}$.

\textbf{Analyse.} Le mouvement est circulaire uniforme : $a_T = 0$ et l'accélération est purement centripète :
\[ a = a_N = \dfrac{v^2}{R} = \dfrac{20^2}{50} = \SI{8,0}{m.s^{-2}} \]
C'est une accélération considérable, proche de $g = \SI{9,8}{m.s^{-2}}$ ! Dirigée vers le centre du virage, elle devra être fournie par le frottement des pneus sur la chaussée (ou par l'inclinaison du virage relevé) : c'est toute la physique du virage, que la deuxième loi de Newton nous permettra de quantifier dans les sections suivantes.
\end{exoresolu}

\courssub{Nature d'un mouvement : lire les signes de $a_T$ et $a_N$}

La décomposition de Frenet donne un critère simple et puissant pour qualifier un mouvement, sans résoudre aucune équation.

\begin{methode}[Déterminer la nature d'un mouvement]
\begin{enumerate}
\item \textbf{Examiner $a_N = \dfrac{v^2}{\rho}$.} Si $a_N = 0$ en tout point, la trajectoire est \emph{rectiligne} (rayon de courbure infini). Si $a_N \neq 0$, la trajectoire est \emph{curviligne}.
\item \textbf{Examiner le signe de $a_T = \dfrac{dv}{dt}$} (ou le produit $a_T \cdot v$, ou encore l'angle entre $\vec{a}$ et $\vec{v}$) :
\begin{itemize}
\item $a_T \cdot v > 0$ ($\vec{a}$ et $\vec{v}$ font un angle aigu) : mouvement \cle{accéléré} ;
\item $a_T \cdot v < 0$ (angle obtus) : mouvement \cle{décéléré} (ou retardé) ;
\item $a_T = 0$ en permanence : mouvement \cle{uniforme}.
\end{itemize}
\item \textbf{Conclure} en combinant les deux : par exemple, $a_N \neq 0$ et $a_T = 0$ caractérisent un mouvement \emph{curviligne uniforme}.
\end{enumerate}
\end{methode}

\begin{exemplebox}[Le lanceur de marteau]
Pendant que le lanceur fait tourner le marteau de plus en plus vite avant de le lâcher, la boule décrit un cercle avec $a_N = \dfrac{v^2}{R} \neq 0$ (elle tourne) \emph{et} $a_T > 0$ (sa vitesse augmente) : le mouvement est \emph{circulaire accéléré}. Le vecteur accélération pointe vers l'avant et vers le centre. À l'instant du lâcher, l'athlète cesse d'exercer sa traction : nous verrons avec le principe d'inertie que le marteau part alors tangentiellement, avec l'énorme vitesse acquise.
\end{exemplebox}

\begin{aretenir}[L'essentiel de la section]
\begin{itemize}
\item Tout mouvement se décrit dans un \cle{référentiel}, avec un repère d'espace et un repère de temps ; le mouvement est relatif au référentiel choisi.
\item $\vec{v} = \dfrac{d\overrightarrow{OM}}{dt}$ est tangent à la trajectoire ; $\vec{a} = \dfrac{d\vec{v}}{dt}$ pointe vers la concavité.
\item MRUV : $v = a_0 t + v_0$, $x = \tfrac{1}{2}a_0 t^2 + v_0 t + x_0$, et $v^2 - v_0^2 = 2a_0(x-x_0)$.
\item Mouvement circulaire : $v = R\omega$ avec $\omega$ en \si{rad.s^{-1}}.
\item Base de Frenet : $\vec{a} = \dfrac{dv}{dt}\,\vec{T} + \dfrac{v^2}{\rho}\,\vec{N}$. La composante tangentielle gouverne la variation de la \emph{norme} de la vitesse, la composante normale celle de sa \emph{direction}.
\item MCU : $\vec{a}$ est centripète, $a = \dfrac{v^2}{R} = R\omega^2 \neq 0$ même si la vitesse est constante en norme.
\item Nature du mouvement : signe de $a_T \cdot v$ (accéléré, décéléré, uniforme) et nullité ou non de $a_N$ (rectiligne ou curviligne).
\end{itemize}
\end{aretenir}

\begin{savaistu}[Pourquoi les virages sont-ils relevés ?]
Sur les autoroutes et les vélodromes, les virages sont inclinés (relevés) : la réaction de la chaussée, perpendiculaire à la piste, possède alors une composante horizontale dirigée vers le centre du virage, qui fournit tout ou partie de l'accélération centripète $\dfrac{v^2}{R}$ sans solliciter le frottement des pneus. C'est exactement le même principe qui permet au benskin de la sortie de Yaoundé de négocier le virage relevé plus vite et plus sûrement. Nous le démontrerons quantitativement avec la deuxième loi de Newton dans la section 3.
\end{savaistu}

\begin{exercice}[Grande roue de la fête foraine]
Une nacelle d'une grande roue de rayon $R = \SI{12}{m}$ effectue un tour complet en $T = \SI{24}{s}$, à vitesse constante.
\begin{enumerate}
\item Calculer la vitesse angulaire $\omega$ et la vitesse linéaire $v$ de la nacelle.
\item Le mouvement est-il accéléré ? Justifier et calculer la norme de l'accélération.
\item Comparer cette accélération à $g = \SI{9,8}{m.s^{-2}}$.
\end{enumerate}
\end{exercice}

\begin{corrige}[Exercice : Grande roue de la fête foraine]
\begin{enumerate}
\item $\omega = \dfrac{2\pi}{T} = \dfrac{2\pi}{24} \approx \SI{0,262}{rad.s^{-1}}$ ; $v = R\omega = 12 \times 0,262 \approx \SI{3,14}{m.s^{-1}}$.
\item Oui : la norme de la vitesse est constante mais sa direction change, donc $a_T = 0$ et $a_N \neq 0$. $a = \dfrac{v^2}{R} = \dfrac{3,14^2}{12} \approx \SI{0,822}{m.s^{-2}}$.
\item $a \approx 0,084\,g$ : l'accélération est faible devant $g$, c'est pourquoi le manège est confortable.
\end{enumerate}
\end{corrige}

\coursec{Référentiels galiléens et première loi de Newton}

Dans la section précédente, nous avons appris à \emph{décrire} un mouvement : repère, vecteurs position, vitesse et accélération, repère de Frenet (accélérations tangentielle et normale). Mais la cinématique seule ne répond pas à une question fondamentale : \emph{pourquoi} un corps se met-il en mouvement, pourquoi s'arrête-t-il, pourquoi sa trajectoire se courbe-t-elle ? Répondre à cette question, c'est entrer dans la \cle{dynamique}, la science des causes du mouvement. Et la dynamique commence par une constatation troublante : la description d'un mouvement dépend de l'observateur. Un passager assis dans le car Yaoundé--Douala roulant à \SI{90}{km.h^{-1}} est immobile pour son voisin de siège, mais file à \SI{90}{km.h^{-1}} pour le mototaxiste arrêté au bord de la route. Qui a raison ? Les deux ! Tout dépend du \cle{référentiel} choisi. Avant d'énoncer les lois de Newton, il faut donc préciser dans quels référentiels ces lois sont valables : ce seront les \cle{référentiels galiléens}.

\courssub{La notion de référentiel galiléen}

\textbf{Intuition et observation.}\par
Imagine un palet posé sur une table parfaitement horizontale et parfaitement lisse. Tu le lâches sans vitesse : il reste immobile. Tu le lances doucement : il glisse en ligne droite, pratiquement sans ralentir. Rien d'étonnant... sauf si l'expérience est réalisée dans un train qui freine brusquement : le même palet, sans que personne ne le touche, se met à accélérer vers l'avant du wagon ! Le principe « un corps sur lequel aucune force ne s'exerce garde un mouvement rectiligne uniforme » est vrai dans la gare, faux dans le train qui freine. Il existe donc une classe privilégiée de référentiels dans lesquels les lois de la mécanique prennent leur forme la plus simple.

\begin{definition}[Référentiel galiléen]
Un \cle{référentiel galiléen} (ou référentiel d'inertie) est un référentiel dans lequel le \cle{principe d'inertie} est vérifié : tout corps mécaniquement isolé (soumis à aucune force, ou à des forces qui se compensent) est soit au repos, soit animé d'un mouvement rectiligne uniforme.
\end{definition}

Cette définition est remarquable : ce n'est pas une propriété géométrique du référentiel, mais une propriété \emph{dynamique}. Un référentiel est galiléen si, et seulement si, le principe d'inertie s'y vérifie expérimentalement.

\begin{propriete}[Référentiels en translation uniforme]
Tout référentiel en translation rectiligne uniforme par rapport à un référentiel galiléen est lui-même galiléen. Il existe donc une infinité de référentiels galiléens, tous en mouvement rectiligne uniforme les uns par rapport aux autres.
\end{propriete}

En revanche, tout référentiel \emph{accéléré} par rapport à un référentiel galiléen (train qui freine, manège en rotation, ascenseur qui démarre) est \cle{non galiléen} : le principe d'inertie n'y est pas vérifié.

\courssub{Les trois référentiels usuels}

En pratique, on utilise une hiérarchie de trois référentiels, de plus en plus « proches » de nous mais de moins en moins rigoureusement galiléens.

\begin{definition}[Référentiels héliocentrique, géocentrique et terrestre]
\begin{itemize}
\item Le \cle{référentiel héliocentrique} (ou de Copernic) a pour origine le centre d'inertie du Soleil et des axes pointant vers trois étoiles lointaines « fixes ». C'est le meilleur référentiel galiléen connu ; on l'utilise pour étudier le mouvement des planètes et des sondes interplanétaires.
\item Le \cle{référentiel géocentrique} a pour origine le centre de la Terre et des axes pointant vers les mêmes étoiles fixes. Il est en translation quasi circulaire autour du Soleil. On l'utilise pour les satellites terrestres et la Lune.
\item Le \cle{référentiel terrestre} (ou référentiel du laboratoire) est lié à la surface de la Terre : ta salle de classe, le stade, la route. Il est entraîné dans la rotation de la Terre sur elle-même et autour du Soleil.
\end{itemize}
\end{definition}

\begin{popfigure}
\begin{tikzpicture}[scale=0.95]
% Héliocentrique
\begin{scope}[shift={(0,0)}]
\fill[popgold] (0,0) circle (0.28);
\draw[popline, dashed] (0,0) circle (0.9);
\fill[popblue] (0.9,0) circle (0.09);
\draw[->, thick, popdark] (0,0) -- (0.7,0.5) node[above, scale=0.7]{étoile fixe};
\node[scale=0.8, popdark, align=center] at (0,-1.5) {Référentiel\\\textbf{héliocentrique}};
\node[scale=0.65, popmuted] at (0,-2.15) {planètes, sondes};
\end{scope}
% Géocentrique
\begin{scope}[shift={(5,0)}]
\fill[popblue] (0,0) circle (0.22);
\draw[popline, dashed] (0,0) circle (0.75);
\fill[popmuted] (0.75,0) circle (0.07);
\draw[->, thick, popdark] (0,0) -- (0.6,0.55) node[above, scale=0.7]{étoile fixe};
\node[scale=0.8, popdark, align=center] at (0,-1.5) {Référentiel\\\textbf{géocentrique}};
\node[scale=0.65, popmuted] at (0,-2.15) {satellites, Lune};
\end{scope}
% Terrestre
\begin{scope}[shift={(10,0)}]
\fill[popblue!25] (0,0) circle (0.55);
\draw[popgreen, thick] (-0.5,-0.15) .. controls (-0.2,0.1) and (0.2,-0.2) .. (0.5,0.1);
\draw[->, thick, poporange] (0,0) -- (0,0.95) node[right, scale=0.7]{axe des pôles};
\draw[->, thick, popdark] (0.55,0.1) -- (1.1,0.1);
\node[scale=0.8, popdark, align=center] at (0,-1.5) {Référentiel\\\textbf{terrestre}};
\node[scale=0.65, popmuted] at (0,-2.15) {expériences de laboratoire};
\end{scope}
\end{tikzpicture}
\legende{Les trois référentiels usuels. Le référentiel héliocentrique est le plus galiléen ; le référentiel terrestre ne l'est qu'approximativement.}
\end{popfigure}

\begin{attention}[Le référentiel terrestre n'est galiléen qu'approximativement]
À cause de la rotation de la Terre, un point de l'équateur possède une accélération centripète $a = \omega^2 R_T \approx \SI{0,034}{m.s^{-2}}$ (avec $\omega = \dfrac{2\pi}{\SI{86164}{s}}$ et $R_T = \SI{6,38e6}{m}$). Cette accélération est faible devant $g \approx \SI{9,8}{m.s^{-2}}$ : pour des expériences \emph{courtes} (quelques secondes à quelques minutes) et de \emph{faible amplitude} (quelques mètres), le référentiel terrestre peut être considéré comme galiléen. En revanche, il ne l'est plus pour les mouvements de grande durée ou de grande échelle (pendule de Foucault, cyclones, tirs balistiques longue portée). \textbf{Dans une copie du Baccalauréat, précise toujours le référentiel d'étude} : « on se place dans le référentiel terrestre supposé galiléen ». C'est un point de rédaction souvent sanctionné.
\end{attention}

\courssub{Centre d'inertie d'un système}

Avant d'énoncer la première loi, il faut un objet mathématique capable de résumer le mouvement d'ensemble d'un système : c'est le \cle{centre d'inertie}.

\begin{definition}[Centre d'inertie]
Le centre d'inertie $G$ d'un système de points matériels $A_i$ de masses $m_i$ est le barycentre de ces points affectés de leurs masses :
\[
\overrightarrow{OG} = \frac{\sum_i m_i \, \overrightarrow{OA_i}}{\sum_i m_i} = \frac{1}{M}\sum_i m_i \, \overrightarrow{OA_i},
\]
où $M = \sum_i m_i$ est la masse totale du système. $G$ est aussi appelé \emph{centre de masse}.
\end{definition}

Pour un \cle{solide homogène} (masse volumique uniforme) possédant un centre de symétrie géométrique, $G$ coïncide avec ce centre : centre d'une sphère, d'un disque, d'un parallélépipède, milieu d'une tige homogène. Pour un solide homogène non symétrique, $G$ se détermine par découpage en éléments simples ou expérimentalement (recherche du point d'équilibre).

\begin{exemplebox}[Centre d'inertie d'un système de deux masses]
Une barre rigide de masse négligeable relie une boule de masse $m_1 = \SI{2,0}{kg}$ placée en $A$ et une boule de masse $m_2 = \SI{6,0}{kg}$ placée en $B$, avec $AB = \SI{40}{cm}$. En prenant $A$ comme origine :
\[
AG = \frac{m_1 \times 0 + m_2 \times AB}{m_1 + m_2} = \frac{6{,}0 \times 40}{8{,}0} = \SI{30}{cm}.
\]
$G$ est à \SI{30}{cm} de $A$, donc plus proche de la masse la plus lourde : le centre d'inertie « penche » toujours du côté des grandes masses.
\end{exemplebox}

\courssub{Systèmes isolés et pseudo-isolés}

\begin{definition}[Système isolé, système pseudo-isolé]
\begin{itemize}
\item Un système est \cle{isolé} s'il n'est soumis à \emph{aucune} force extérieure : $\sum \vec{F}_{\text{ext}} = \vec{0}$ car il n'y a pas de force du tout. C'est un cas idéal (un corps seul dans l'espace lointain).
\item Un système est \cle{pseudo-isolé} s'il est soumis à des forces extérieures dont la \emph{somme vectorielle est nulle} : les forces se compensent exactement.
\[
\sum \vec{F}_{\text{ext}} = \vec{0}.
\]
\end{itemize}
\end{definition}

L'exemple canonique est le \cle{mobile sur table à coussin d'air} : une soufflerie crée un mince film d'air sous le mobile, ce qui supprime quasiment les frottements. Le mobile est soumis à deux forces : son poids $\vec{P}$ (vertical, vers le bas) et la réaction de la table $\vec{R}$ (verticale, vers le haut, portée par le coussin d'air). La table étant horizontale et le mobile sans mouvement vertical, ces deux forces se compensent : $\vec{P} + \vec{R} = \vec{0}$. Le mobile est pseudo-isolé.

\begin{popfigure}
\begin{tikzpicture}[scale=1.0]
% Table
\fill[popline!40] (-3,-0.9) rectangle (3,-0.55);
\draw[popdark, thick] (-3,-0.55) -- (3,-0.55);
% Coussin d'air
\draw[popblue, dashed] (-1.1,-0.5) -- (1.1,-0.5);
\node[popblue, scale=0.7] at (2.2,-0.35) {coussin d'air};
% Mobile
\fill[poporange!80] (-1,-0.5) rectangle (1,0.3);
\draw[popdark] (-1,-0.5) rectangle (1,0.3);
\fill[popdark] (0,-0.1) circle (0.06);
\node[popdark, scale=0.8, above right=0.02 and 0.05 of {(0,-0.1)}] {$G$};
% Forces
\draw[->, very thick, poppink] (0,-0.1) -- (0,1.3) node[right]{$\vec{R}$};
\draw[->, very thick, poppurple] (0,-0.1) -- (0,-1.5) node[right]{$\vec{P}$};
% Vitesse
\draw[->, very thick, popgreen] (1.1,0.1) -- (2.4,0.1) node[above]{$\vec{v}_G$};
\node[scale=0.8, popmuted, align=center] at (0,-2.1) {$\vec{P} + \vec{R} = \vec{0}$ : le mobile est pseudo-isolé, $G$ est en MRU};
\end{tikzpicture}
\legende{Mobile pseudo-isolé sur table à coussin d'air : le poids $\vec{P}$ et la réaction $\vec{R}$ se compensent ; le centre d'inertie $G$ garde un mouvement rectiligne uniforme.}
\end{popfigure}

Autre exemple de la vie quotidienne : un véhicule roulant à \emph{vitesse constante} sur une route \emph{rectiligne et horizontale}. Il est soumis à son poids $\vec{P}$, à la réaction de la route $\vec{R}$, à la force motrice $\vec{F}_m$ du moteur et aux forces de frottement $\vec{f}$ (air + roulement). À vitesse constante, $\vec{F}_m + \vec{f} = \vec{0}$ et $\vec{P} + \vec{R} = \vec{0}$ : le véhicule est pseudo-isolé. C'est d'ailleurs ainsi que le principe d'inertie se vérifie au quotidien, contrairement à l'intuition aristotélicienne qui voudrait qu'« il faille une force pour maintenir un mouvement ».

\courssub{Première loi de Newton : le principe d'inertie}

\begin{experience}[Mise en évidence sur table à coussin d'air]
\textbf{Matériel :} table à coussin d'air horizontale, mobile autoporteur avec étincelage, soufflerie, papier enregistreur.\par
\textbf{Protocole :} on met la soufflerie en marche, on lance le mobile d'une pichenette et on enregistre les positions successives de $G$ à intervalles de temps égaux $\tau = \SI{40}{ms}$.\par
\textbf{Observations :} les points enregistrés sont \emph{alignés} et \emph{équidistants}.\par
\textbf{Interprétation :} des positions équidistantes à intervalles égaux signifient une vitesse constante ; l'alignement signifie une trajectoire rectiligne. Le mobile, pourtant soumis à son poids, se déplace en ligne droite à vitesse constante : c'est la somme des forces qui est nulle, pas « l'absence de mouvement ». Le principe d'inertie est vérifié.
\end{experience}

\begin{propriete}[Première loi de Newton — principe d'inertie]
Dans un référentiel galiléen, si la somme vectorielle des forces extérieures appliquées à un système est nulle, alors le centre d'inertie $G$ du système est au repos (s'il était au repos) ou animé d'un mouvement rectiligne uniforme :
\[
\sum \vec{F}_{\text{ext}} = \vec{0} \quad \Longrightarrow \quad \vec{v}_G = \overrightarrow{\text{const}}.
\]
\textbf{Réciproque :} dans un référentiel galiléen, si le centre d'inertie d'un système est au repos ou en mouvement rectiligne uniforme, alors la somme des forces extérieures est nulle :
\[
\vec{v}_G = \overrightarrow{\text{const}} \quad \Longrightarrow \quad \sum \vec{F}_{\text{ext}} = \vec{0}.
\]
Le principe d'inertie est donc une \emph{équivalence} : $\sum \vec{F}_{\text{ext}} = \vec{0} \iff \vec{v}_G = \overrightarrow{\text{const}}$.
\end{propriete}

\begin{aretenir}[Ce qu'il faut retenir du principe d'inertie]
\begin{itemize}
\item Ce n'est pas la \emph{force} qui entretient le mouvement, c'est la \emph{variation de vitesse} qui exige une force résultante.
\item La réciproque est un outil d'analyse très puissant : observer un mouvement rectiligne uniforme permet de conclure immédiatement que les forces se compensent, et d'en déduire des valeurs (frottements, réaction, tension).
\item Le principe ne concerne que le centre d'inertie $G$ : un solide pseudo-isolé peut tourner sur lui-même à vitesse angulaire constante pendant que $G$ est en MRU.
\end{itemize}
\end{aretenir}

\begin{methode}[Appliquer le principe d'inertie (bilan de forces et nature du mouvement)]
\begin{enumerate}
\item \textbf{Choisir le référentiel} : l'énoncer explicitement (ex. « référentiel terrestre supposé galiléen »).
\item \textbf{Isoler le système} : définir clairement le système étudié (point matériel ou solide).
\item \textbf{Dresser le bilan des forces extérieures} : lister \emph{toutes} les forces (poids, réaction(s), tension(s), frottements, actions à distance...), les nommer, les orienter, les représenter sur un schéma.
\item \textbf{Calculer la somme vectorielle} $\sum \vec{F}_{\text{ext}}$ (souvent par projection sur des axes bien choisis).
\item \textbf{Conclure} :
\begin{itemize}
\item Si $\sum \vec{F}_{\text{ext}} = \vec{0}$ : le centre d'inertie $G$ est en MRU (ou au repos). La vitesse $\vec{v}_G$ est constante.
\item Si $\sum \vec{F}_{\text{ext}} \neq \vec{0}$ : le principe d'inertie ne s'applique pas ; il faut utiliser la 2\textsuperscript{e} loi de Newton ($\sum \vec{F}_{\text{ext}} = m\vec{a}_G$).
\end{itemize}
\end{enumerate}
\end{methode}

\begin{savaistu}[De Galilée à Newton]
C'est Galilée qui, au début du XVII\textsuperscript{e} siècle, comprit le premier qu'une bille lancée sur un plan horizontal de plus en plus poli ralentit de moins en moins, et qu'à la limite d'un plan parfait elle ne ralentirait \emph{jamais} : le mouvement uniforme n'a pas besoin de cause. En 1687, Isaac Newton reprit cette idée comme \emph{première loi} de ses \emph{Philosophiae Naturalis Principia Mathematica} : « Tout corps persévère dans son état de repos ou de mouvement rectiligne uniforme, à moins que des forces imprimées ne le contraignent d'en changer. » Cette loi renversait vingt siècles de physique aristotélicienne et fondait la mécanique moderne — celle qui, trois siècles plus tard, permettra de calculer les orbites des satellites de télécommunication qui couvrent le Cameroun.
\end{savaistu}

\courssub{Applications du principe d'inertie}

\textbf{Lecture inertielle d'une situation quotidienne : le freinage brutal.}\par
Un car de la ligne Douala--Bafoussam roule à $v_0 = \SI{72}{km.h^{-1}} = \SI{20}{m.s^{-1}}$. Le chauffeur freine brutalement. Les passagers debout, non retenus, sont « projetés » vers l'avant. En réalité, rien ne les projette : par inertie, chaque passager \emph{conserve} sa vitesse $\vec{v}_0$ tant qu'aucune force horizontale ne s'exerce sur lui, pendant que le car, lui, ralentit sous l'action des freins. Dans le référentiel terrestre (galiléen), le passager continue tout droit à \SI{20}{m.s^{-1}} et rattrape le plancher qui décélère. La ceinture de sécurité ou la poignée fournit alors la force nécessaire pour le faire décélérer avec le véhicule. Le principe d'inertie explique à la fois le danger et le remède.

\begin{exoresolu}[Solide pseudo-isolé sur table horizontale lisse]
Un solide $S$ de masse $m = \SI{500}{g}$ glisse sur une table horizontale parfaitement lisse (frottements négligeables). Lâché avec une vitesse initiale $v_0 = \SI{2,0}{m.s^{-1}}$, il parcourt une distance $d = \SI{3,0}{m}$ avant d'atteindre le bord de la table. On se place dans le référentiel terrestre supposé galiléen.\par
1. Faire le bilan des forces extérieures appliquées à $S$.\par
2. Déterminer la nature du mouvement de $G$.\par
3. Calculer la durée $\Delta t$ du parcours sur la table.
\tcblower
\textbf{1. Bilan des forces.} Le système étudié est le solide $S$. Il est soumis à :
\begin{itemize}
\item son poids $\vec{P} = m\vec{g}$, vertical, vers le bas, appliqué en $G$ ;
\item la réaction $\vec{R}$ de la table, verticale, vers le haut (table lisse : pas de composante tangentielle de frottement).
\end{itemize}
\textbf{2. Nature du mouvement.} Le solide ne quitte pas la table et ne s'enfonce pas : il n'y a pas de mouvement vertical, donc $\vec{P} + \vec{R} = \vec{0}$. Le solide est \cle{pseudo-isolé}. D'après le principe d'inertie, son centre d'inertie $G$ est animé d'un \textbf{mouvement rectiligne uniforme} à la vitesse $v_0 = \SI{2,0}{m.s^{-1}}$.\par
\textbf{3. Durée du parcours.} Pour un MRU, $d = v_0 \, \Delta t$, d'où :
\[
\Delta t = \frac{d}{v_0} = \frac{3{,}0}{2{,}0} = \SI{1,5}{s}.
\]
Le solide atteint le bord de la table au bout de $\Delta t = \SI{1,5}{s}$, sans avoir ralenti : sur une table idéalement lisse, rien ne modifie sa vitesse.
\end{exoresolu}

\begin{attention}[Pièges fréquents autour du principe d'inertie]
\begin{itemize}
\item \textbf{« Le corps est en mouvement, donc une force le pousse. »} Faux ! C'est l'intuition d'Aristote, réfutée par Galilée. Un mouvement rectiligne uniforme n'exige \emph{aucune} force résultante.
\item \textbf{Oublier la réaction du support.} Sur une table, le solide n'est pas soumis « à une seule force, son poids » : la réaction $\vec{R}$ compense $\vec{P}$. Un bilan de forces incomplet fausse toute la suite.
\item \textbf{Appliquer le principe dans un mauvais référentiel.} Dans le car qui freine (référentiel non galiléen), le passager accélère « sans force » : le principe d'inertie n'y est pas applicable directement.
\item \textbf{Confondre « isolé » et « pseudo-isolé ».} Un système pseudo-isolé subit bien des forces ; c'est leur \emph{somme vectorielle} qui est nulle.
\item \textbf{Croire que le principe interdit toute rotation.} Il ne porte que sur le centre d'inertie $G$ : un solide pseudo-isolé peut tourner autour de $G$ à vitesse angulaire constante.
\end{itemize}
\end{attention}

\begin{exercice}[À toi de jouer]
Un palet de hockey de masse $m = \SI{160}{g}$ glisse sur la glace d'une patinoire horizontale. On enregistre ses positions toutes les $\tau = \SI{50}{ms}$ : les distances parcourues sont toutes égales à \SI{75}{cm} et les points sont alignés.\par
1. Montrer que le palet est pseudo-isolé.\par
2. Calculer sa vitesse.\par
3. Le palet heurte ensuite la bande de la patinoire. Le principe d'inertie reste-t-il applicable pendant le choc ? Justifier.
\end{exercice}

\begin{corrige}[Exercice 1]
\textbf{1.} Les positions sont alignées (trajectoire rectiligne) et les distances parcourues pendant des durées égales sont égales (vitesse constante) : $G$ est en mouvement rectiligne uniforme dans le référentiel terrestre supposé galiléen. D'après la \emph{réciproque} du principe d'inertie, $\sum \vec{F}_{\text{ext}} = \vec{0}$ : le palet est pseudo-isolé (poids et réaction de la glace se compensent, frottements négligeables).\par
\textbf{2.} $v = \dfrac{d}{\tau} = \dfrac{\SI{0,75}{m}}{\SI{0,050}{s}} = \SI{15}{m.s^{-1}}$, soit \SI{54}{km.h^{-1}}.\par
\textbf{3.} Pendant le choc, la bande exerce sur le palet une force de contact intense et brève : la somme des forces n'est plus nulle, la vitesse change brutalement (direction et norme). Le principe d'inertie ne s'applique plus pendant le choc ; il redevient applicable après, si le palet repart en glissant sur la glace.
\end{corrige}

Nous savons désormais reconnaître les référentiels où la mécanique s'écrit simplement, et traiter le cas où les forces se compensent. Mais que se passe-t-il lorsque la somme des forces extérieures n'est \emph{pas} nulle ? C'est toute la puissance de la deuxième loi de Newton, le théorème du centre d'inertie $\sum \vec{F}_{\text{ext}} = m\,\vec{a}_G$, que nous étudierons dans la section suivante.

\coursec{Deuxième et troisième lois de Newton}

La première loi de Newton nous a appris à reconnaître les situations où un solide conserve son mouvement : c'est le cas lorsque la somme des forces extérieures est nulle. Mais que se passe-t-il lorsque cette somme n'est \emph{pas} nulle ? Une mototaxi qui accélère au feu vert, un ballon qui retombe après un dégagement, un camion qui ralentit dans la descente de la côte de Mbanga : dans tous ces cas, le mouvement \emph{change}. La deuxième loi de Newton donne la réponse quantitative : elle relie la cause (les forces) à l'effet (l'accélération). C'est la loi fondamentale de toute la dynamique du point, et l'outil principal de cette leçon.

\courssub{Le théorème du centre d'inertie (TCI)}

\textbf{Intuition et observation.}\par
Pousse doucement un cahier sur la table : il accélère faiblement. Pousse-le fort : il accélère davantage. Pousse maintenant, avec la même force, un cahier puis une lourde encyclopédie : l'encyclopédie accélère beaucoup moins. L'expérience suggère donc que l'accélération est \cle{proportionnelle à la force} appliquée et \cle{inversement proportionnelle à la masse}. C'est exactement le contenu de la deuxième loi de Newton.

\begin{propriete}[Théorème du centre d'inertie (deuxième loi de Newton)]
Dans un \cle{référentiel galiléen}, la somme vectorielle des \cle{forces extérieures} appliquées à un système de masse $m$ \cle{constante} est égale au produit de la masse par le vecteur accélération de son centre d'inertie $G$ :
\[ \sum \vec{F}_{\mathrm{ext}} = m\,\vec{a}_G \]
Cette loi fondamentale est \emph{admise} (elle ne se démontre pas) ; elle est vérifiée par l'ensemble des expériences de mécanique. Analyse dimensionnelle : $[m\,a] = \mathrm{kg\cdot m\cdot s^{-2}} = \mathrm{N}$, cohérent avec l'unité de force.
\end{propriete}

\begin{attention}[Conditions d'application à ne jamais oublier]
Le TCI ne s'applique que si les trois conditions suivantes sont réunies :
\begin{itemize}
\item le \cle{référentiel} d'étude est \textbf{galiléen} (référentiel terrestre pour des mouvements de courte durée, référentiel géocentrique pour les satellites) ;
\item la \cle{masse} du système est \textbf{constante} (le TCI sous cette forme ne vaut pas pour une fusée qui éjecte ses gaz) ;
\item le bilan des forces est \textbf{exhaustif} et ne contient que des forces \emph{extérieures} : les forces intérieures au système se compensent deux à deux (troisième loi) et n'apparaissent pas.
\end{itemize}
Oublier une force (souvent les frottements ou la réaction du support) est l'erreur la plus fréquente au Baccalauréat.
\end{attention}

\begin{methode}[Résoudre un problème de dynamique en 5 étapes]
\begin{enumerate}
\item \textbf{Système} : délimiter précisément le corps étudié (et sa masse $m$).
\item \textbf{Référentiel} : choisir un référentiel galiléen et le préciser dans la copie (ex. : « référentiel terrestre supposé galiléen »).
\item \textbf{Bilan des forces extérieures} : les lister et les représenter sur un schéma (poids, réaction, frottements, tension, poussée...).
\item \textbf{Projection du TCI} : choisir des axes adaptés au mouvement (souvent un axe le long du mouvement, un axe perpendiculaire) et projeter $\sum\vec{F}_{\mathrm{ext}} = m\vec{a}_G$ sur chacun.
\item \textbf{Résolution} : obtenir l'accélération, écrire l'équation différentielle, intégrer deux fois en utilisant les conditions initiales pour obtenir $v(t)$ et $x(t)$.
\end{enumerate}
\end{methode}

\courssub{Application 1 : la chute libre verticale}

\begin{definition}[Chute libre]
Un solide est en \cle{chute libre} lorsqu'il n'est soumis qu'à son \cle{poids}. C'est le cas d'un objet lâché dans le vide, ou dans l'air si les frottements de l'air sont négligeables (objet dense, faible vitesse, petite distance).
\end{definition}

Appliquons la méthode. Système : bille de masse $m$. Référentiel terrestre galiléen. Bilan des forces : le poids $\vec{P} = m\vec{g}$, seul. On choisit un axe vertical $(Oz)$ descendant, l'origine $O$ au point de lâcher. Le TCI s'écrit $m\vec{g} = m\vec{a}$, soit en projection :
\[ a_z = g \]
L'accélération est \textbf{constante} et \textbf{indépendante de la masse} : en chute libre, tous les corps tombent avec la même accélération, comme Galilée l'avait pressenti. En intégrant avec $v(0)=0$ et $z(0)=0$ :
\[ v(t) = g\,t \qquad\text{et}\qquad z(t) = \tfrac{1}{2}\,g\,t^{2} \]
En éliminant le temps entre ces deux relations ($t = v/g$), on obtient la \cle{vitesse d'impact} après une chute de hauteur $h$ :
\[ v^{2} = 2\,g\,h \qquad\Longleftrightarrow\qquad v = \sqrt{2\,g\,h} \]

\begin{exemplebox}[Chute d'une mangue]
Une mangue se détache d'un arbre à $h = \SI{5,0}{m}$ du sol. On prend $g = \SI{10}{m.s^{-2}}$. Sa vitesse au sol est :
\[ v = \sqrt{2\times 10\times 5{,}0} = \sqrt{100} = \SI{10}{m.s^{-1}} \quad (\text{soit } \SI{36}{km.h^{-1}} !) \]
Le temps de chute vaut $t = v/g = \SI{1,0}{s}$. C'est pourquoi il est dangereux de stationner sous un manguier chargé par grand vent.
\end{exemplebox}

\courssub{Application 2 : solide glissant sur un plan incliné}

\textbf{Sans frottements.}\par
Un bloc de masse $m$ glisse sur un plan incliné d'un angle $\alpha$ avec l'horizontale. Bilan des forces : le poids $\vec{P}$ et la réaction normale $\vec{R}_N$, perpendiculaire au plan (pas de frottements). On choisit l'axe $(Ox)$ le long de la ligne de plus grande pente, dirigé vers le bas, et $(Oy)$ perpendiculaire au plan. Le poids se décompose en :
\[ P_x = m g \sin\alpha \qquad ; \qquad P_y = -m g \cos\alpha \]
Le mouvement se fait selon $(Ox)$, donc $a_y = 0$. La projection du TCI donne :
\[ \text{sur } (Ox) : \quad m g \sin\alpha = m a \qquad ; \qquad \text{sur } (Oy) : \quad R_N - m g \cos\alpha = 0 \]
D'où l'accélération et la réaction normale :
\[ a = g\,\sin\alpha \qquad ; \qquad R_N = m g \cos\alpha \]
On retrouve bien $a = g$ pour $\alpha = 90^{\circ}$ (chute libre) et $a = 0$ pour $\alpha = 0$ (plan horizontal) : le résultat est cohérent aux limites.

\begin{popfigure}
\begin{center}
\begin{tikzpicture}[scale=1.1,>=Stealth]
% plan incliné
\coordinate (O) at (0,0);
\coordinate (A) at (6,0);
\coordinate (B) at (6,3.2);
\fill[popcream] (O)--(A)--(B)--cycle;
\draw[thick,popdark] (O)--(A)--(B)--cycle;
\draw[popdark] (4.6,0) arc[start angle=0,end angle=28.1,radius=1.4];
\node[popdark] at (5.0,0.35) {$\alpha$};
% bloc
\coordinate (G) at (4.35,2.05);
\begin{scope}[rotate around={28.1:(G)}]
\fill[popblueL,draw=popblue,thick] ($(G)+(-0.55,-0.35)$) rectangle ($(G)+(0.55,0.35)$);
\end{scope}
\fill[popdark] (G) circle (0.05) node[above left=1pt,popdark] {$G$};
% poids et composantes
\draw[->,very thick,popink] (G) -- ($(G)+(0,-1.9)$) node[below] {$\vec{P}$};
\draw[->,thick,poporange] (G) -- ($(G)+(-1.68,-0.89)$) node[below left] {$\vec{P}_x$};
\draw[->,thick,poporange] (G) -- ($(G)+(0.89,-1.68)$) node[below right] {$\vec{P}_y$};
\draw[dashed,popmuted] ($(G)+(-1.68,-0.89)$) -- ($(G)+(0,-1.9)$) -- ($(G)+(0.89,-1.68)$);
% réaction et frottement
\draw[->,very thick,popgreen] (G) -- ($(G)+(-0.72,1.36)$) node[above left] {$\vec{R}_N$};
\draw[->,very thick,poppurple] ($(G)+(-0.5,-0.1)$) -- ($(G)+(0.75,0.56)$) node[above right] {$\vec{f}$};
% axes
\draw[->,thick,popdark] ($(G)+(1.6,0.85)$) -- ($(G)+(0.6,0.32)$) node[midway,above,sloped] {$x$};
\draw[->,thick,popdark] ($(G)+(1.6,0.85)$) -- ($(G)+(1.32,1.66)$) node[midway,right] {$y$};
\end{tikzpicture}
\end{center}
\legende{Bloc sur un plan incliné : poids décomposé selon les axes, réaction normale $\vec{R}_N$ et force de frottement $\vec{f}$ opposée au glissement.}
\end{popfigure}

\textbf{Avec frottements solides.}\par
Lorsque le contact est rugueux, il apparaît une force de \cle{frottement} $\vec{f}$, tangente au plan et opposée au glissement. Pour un glissement effectif, sa norme est proportionnelle à la réaction normale :
\[ f = \mu\,R_N = \mu\,m g \cos\alpha \]
où $\mu$ est le \cle{coefficient de frottement} (sans unité), qui dépend des surfaces en contact. La projection du TCI sur $(Ox)$ devient :
\[ m g \sin\alpha - \mu\,m g \cos\alpha = m a \]
\begin{propriete}[Accélération sur un plan incliné avec frottements]
\[ a = g\left(\sin\alpha - \mu\cos\alpha\right) \]
Le mouvement n'est accéléré que si $\tan\alpha > \mu$ ; sinon le bloc reste immobile ou glisse à vitesse constante (cas limite $\tan\alpha = \mu$).
\end{propriete}

\begin{exoresolu}[Bloc sur un plan incliné]
Énoncé. Un bloc de masse $m = \SI{2,0}{kg}$ est lâché sans vitesse initiale sur un plan incliné de $\alpha = 30^{\circ}$, parfaitement lisse. On prend $g = \SI{10}{m.s^{-2}}$. (1) Calculer son accélération. (2) Quelle distance parcourt-il en $\SI{2,0}{s}$ ? (3) Mêmes questions si le contact présente un frottement de coefficient $\mu = 0{,}20$.
\tcblower
Solution. (1) Sans frottement : $a = g\sin\alpha = 10\times \sin 30^{\circ} = 10\times 0{,}50 = \SI{5,0}{m.s^{-2}}$. Remarque : la masse n'intervient pas, elle se simplifie.
(2) Mouvement rectiligne uniformément accéléré avec $v_0 = 0$ : $x(t) = \tfrac{1}{2}at^{2} = \tfrac{1}{2}\times 5{,}0\times (2{,}0)^{2} = \SI{10}{m}$.
(3) Avec frottements :
\[ a = g(\sin\alpha - \mu\cos\alpha) = 10\times\left(0{,}50 - 0{,}20\times 0{,}866\right) = 10\times 0{,}327 \approx \SI{3,3}{m.s^{-2}}. \]
La distance parcourue en $\SI{2,0}{s}$ devient $x = \tfrac{1}{2}\times 3{,}3\times (2{,}0)^{2} = \SI{6,5}{m}$. Les frottements réduisent l'accélération, conformément à l'expérience quotidienne.
\end{exoresolu}

\courssub{Application 3 : le véhicule en virage — la force centripète}

Pour qu'un véhicule de masse $m$ décrive un virage circulaire de rayon $r$ à la vitesse $v$, il faut une accélération \cle{normale} (centripète) $a_N = v^{2}/r$, dirigée vers le centre du virage (voir le repère de Frenet, section 1). D'après le TCI, il faut donc une \cle{force centripète} résultante :
\[ F_c = \dfrac{m v^{2}}{r} \]
Sur une route horizontale, cette force est fournie par le \cle{frottement} des pneus sur la chaussée : $f = m v^{2}/r$. Si la chaussée est glissante (pluie sur l'axe Yaoundé--Douala), $f$ est plafonnée par $\mu R_N = \mu m g$ : la vitesse maximale est alors $v_{\max} = \sqrt{\mu g r}$. Au-delà, le véhicule part en tout-droit.

Sur un \cle{virage relevé} (piste inclinée d'un angle $\alpha$), la réaction de la piste n'est plus verticale : sa composante horizontale $R_N\sin\alpha$ fournit la force centripète, tandis que sa composante verticale $R_N\cos\alpha$ équilibre le poids. En divisant membre à membre :
\[ \tan\alpha = \dfrac{v^{2}}{g\,r} \]
Le virage peut alors être négocié \emph{sans aucun frottement} à la vitesse $v = \sqrt{g r \tan\alpha}$.

\begin{popfigure}
\begin{center}
\begin{tikzpicture}[scale=1.15,>=Stealth]
% piste relevée
\fill[popcream] (0,0) -- (5.5,0) -- (5.5,2.4) -- cycle;
\draw[thick,popdark] (0,0) -- (5.5,0) -- (5.5,2.4) -- cycle;
\draw[popdark] (4.1,0) arc[start angle=0,end angle=23.6,radius=1.4];
\node[popdark] at (4.5,0.32) {$\alpha$};
% véhicule
\coordinate (G) at (3.6,1.35);
\begin{scope}[rotate around={23.6:(G)}]
\fill[popblueL,draw=popblue,thick] ($(G)+(-0.6,-0.3)$) rectangle ($(G)+(0.6,0.3)$);
\fill[popdark] ($(G)+(-0.35,-0.38)$) circle (0.09);
\fill[popdark] ($(G)+(0.35,-0.38)$) circle (0.09);
\end{scope}
\fill[popdark] (G) circle (0.05) node[above=2pt,popdark] {$G$};
% forces
\draw[->,very thick,popink] (G) -- ($(G)+(0,-1.7)$) node[below] {$\vec{P}$};
\draw[->,very thick,popgreen] (G) -- ($(G)+(-0.85,1.95)$) node[above left] {$\vec{R}_N$};
\draw[->,thick,poporange,dashed] (G) -- ($(G)+(-2.0,0)$) node[left] {$R_N\sin\alpha$};
\draw[->,thick,poporange,dashed] (G) -- ($(G)+(0,1.83)$) node[right] {$R_N\cos\alpha$};
% centre du virage
\draw[->,thick,poppurple] ($(G)+(-2.0,-0.55)$) -- ($(G)+(-1.0,-0.55)$) node[midway,below] {vers le centre $O$};
\end{tikzpicture}
\end{center}
\legende{Virage relevé : la composante horizontale de la réaction $\vec{R}_N$ joue le rôle de force centripète, dirigée vers le centre du virage.}
\end{popfigure}

\begin{savaistu}[Virages relevés : circuits et vélodromes]
Les concepteurs de circuits automobiles et de vélodromes inclinent les virages d'un angle calculé par $\tan\alpha = v^{2}/(gr)$ : à la vitesse prévue, les véhicules négocient le virage \emph{même sur piste glissante}, sans avoir besoin du frottement. Les patinoires de vitesse et certaines bretelles d'autoroute utilisent le même principe. À l'inverse, un virage \emph{déversé} vers l'extérieur (mauvaise construction) est extrêmement dangereux : la composante de la réaction écarte le véhicule du centre au lieu de l'y ramener.
\end{savaistu}

\courssub{Troisième loi de Newton : principe des actions réciproques}

\textbf{Observation.}\par
Un rameur dans une pirogue sur le Wouri pousse l'eau vers l'arrière avec sa pagaie : la pirogue avance vers l'avant. En marchant, tu pousses le sol vers l'arrière avec ton pied : le sol te pousse vers l'avant. À chaque fois qu'un corps $A$ exerce une force sur un corps $B$, le corps $B$ « répond » par une force sur $A$.

\begin{propriete}[Troisième loi de Newton — actions réciproques]
Si un corps $A$ exerce sur un corps $B$ une force $\vec{F}_{A/B}$, alors le corps $B$ exerce sur $A$ une force $\vec{F}_{B/A}$ telle que :
\[ \vec{F}_{A/B} = -\,\vec{F}_{B/A} \]
Ces deux forces ont \textbf{même support} (la droite d'action commune), \textbf{même norme}, des \textbf{sens opposés}, et s'exercent sur \textbf{deux corps différents}. Elles existent toujours simultanément, quels que soient les mouvements des corps.
\end{propriete}

\begin{popfigure}
\begin{center}
\begin{tikzpicture}[scale=1.0,>=Stealth]
% eau
\fill[popblueL] (-0.5,-1.6) rectangle (10.5,-0.6);
\node[popblue] at (9.6,-1.1) {eau};
% pirogue
\draw[thick,popdark,fill=popgoldL] (1.5,0) .. controls (2.5,-0.55) and (5.5,-0.55) .. (6.5,0) -- (6.2,0.25) -- (1.8,0.25) -- cycle;
\node[popdark] at (4,0.55) {pirogue};
% rameur
\fill[popdark] (3.2,0.55) circle (0.14);
\draw[thick,popdark] (3.2,0.4) -- (3.2,0.0);
% pagaie
\draw[very thick,poporange] (3.2,0.15) -- (4.6,-1.3);
\node[poporange] at (5.0,-1.35) {pagaie};
% forces
\draw[->,very thick,popink] (4.6,-1.0) -- (5.9,-1.0) node[right] {$\vec{F}_{\text{pagaie/eau}}$};
\draw[->,very thick,popgreen] (4.4,-0.75) -- (3.1,-0.75) node[left] {$\vec{F}_{\text{eau/pagaie}}$};
\draw[->,very thick,poppurple] (2.2,0.1) -- (0.9,0.1) node[left] {pirogue en mouvement};
\end{tikzpicture}
\end{center}
\legende{Actions réciproques : la pagaie pousse l'eau vers l'arrière, l'eau pousse la pagaie (donc le rameur et la pirogue) vers l'avant.}
\end{popfigure}

\textbf{Exemples d'actions réciproques.}\par
\begin{itemize}
\item \textbf{La marche} : le pied exerce $\vec{F}_{\text{pied/sol}}$ vers l'arrière ; le sol exerce $\vec{F}_{\text{sol/pied}}$ vers l'avant — c'est elle qui nous propulse.
\item \textbf{Interaction Terre--objet} : la Terre attire la mangue vers le bas (poids), et la mangue attire la Terre vers le haut avec une force de même norme. La Terre, de masse immense, n'acquiert qu'une accélération indétectable ($a = F/M_{\text{Terre}}$).
\item \textbf{Propulsion} : la pirogue et sa pagaie, le nageur qui repousse l'eau, le ballon qui se dégonfle en chassant l'air.
\end{itemize}

\begin{attention}[Action--réaction ou forces qui se compensent ?]
Deux forces de même norme et de sens opposés ne forment un couple action--réaction que si elles s'exercent sur \textbf{deux corps différents}. Contre-exemple classique : un livre posé sur une table est soumis à son poids $\vec{P}$ (sur le livre) et à la réaction $\vec{R}_N$ de la table (\emph{sur le livre aussi}). Ces deux forces se compensent car le livre est immobile, mais ce ne sont \textbf{pas} des actions réciproques ! La réciproque du poids est l'attraction du livre \emph{sur la Terre} ; la réciproque de $\vec{R}_N$ est la force du livre \emph{sur la table}. Critère décisif : un couple action--réaction ne peut jamais s'annuler « sur un même système », puisque les deux forces agissent sur des systèmes différents.
\end{attention}

\begin{exoresolu}[Démarrage d'une mototaxi]
Énoncé. Une mototaxi et son conducteur (masse totale $m = \SI{180}{kg}$) démarrent sur une route horizontale avec une accélération $a = \SI{2,0}{m.s^{-2}}$. (1) Quelle force motrice horizontale la route exerce-t-elle sur les pneus ? (2) Quelle est l'action réciproque de cette force ?
\tcblower
Solution. (1) Référentiel terrestre galiléen. Sur l'horizontale, seule la force motrice $\vec{F}$ (frottement de la route sur les pneus) intervient, le poids et la réaction normale se compensant verticalement. Le TCI projeté sur l'axe du mouvement donne :
\[ F = m\,a = 180\times 2{,}0 = \SI{360}{N} \]
(2) Réciproque : les pneus exercent sur la route une force $\vec{F}_{\text{pneus/route}} = -\vec{F}$, de norme $\SI{360}{N}$, dirigée vers l'arrière. C'est en repoussant la route vers l'arrière que la moto est poussée vers l'avant — exactement comme le rameur repousse l'eau.
\end{exoresolu}

\begin{exercice}[Virage sur route mouillée]
Une voiture de masse $m = \SI{900}{kg}$ aborde un virage circulaire horizontal de rayon $r = \SI{50}{m}$. Le coefficient de frottement pneus--chaussée vaut $\mu = 0{,}60$ sur route mouillée. On prend $g = \SI{10}{m.s^{-2}}$. (1) Calculer la vitesse maximale de négociation du virage, en $\mathrm{m\cdot s^{-1}}$ puis en $\mathrm{km\cdot h^{-1}}$. (2) Le conducteur roule à $\SI{72}{km.h^{-1}}$ : franchit-il le virage ? (3) Quel angle de relevage $\alpha$ permettrait de négocier ce virage à $\SI{72}{km.h^{-1}}$ sans frottement ?
\end{exercice}

\begin{corrige}[Exercice : virage sur route mouillée]
(1) $v_{\max} = \sqrt{\mu g r} = \sqrt{0{,}60\times 10\times 50} = \sqrt{300} \approx \SI{17,3}{m.s^{-1}}$, soit environ $\SI{62}{km.h^{-1}}$.
(2) $\SI{72}{km.h^{-1}} = \SI{20}{m.s^{-1}} > v_{\max}$ : la force de frottement maximale $\mu m g = \SI{5400}{N}$ est inférieure à la force centripète requise $mv^{2}/r = 900\times 400/50 = \SI{7200}{N}$. Le véhicule dérape vers l'extérieur.
(3) $\tan\alpha = v^{2}/(gr) = 400/(10\times 50) = 0{,}80$, d'où $\alpha \approx 39^{\circ}$ : un relevage important, réservé aux circuits.
\end{corrige}

\begin{aretenir}[L'essentiel de la section]
\begin{itemize}
\item Théorème du centre d'inertie : dans un référentiel galiléen, $\sum\vec{F}_{\mathrm{ext}} = m\,\vec{a}_G$ (masse constante, bilan exhaustif des forces extérieures).
\item Méthode en 5 étapes : système, référentiel, bilan des forces, projection, résolution.
\item Chute libre : $a = g$, $v = \sqrt{2gh}$. Plan incliné : $a = g(\sin\alpha - \mu\cos\alpha)$. Virage : force centripète $F_c = mv^{2}/r$ ; virage relevé : $\tan\alpha = v^{2}/(gr)$.
\item Troisième loi : $\vec{F}_{A/B} = -\vec{F}_{B/A}$, même support, même norme, \emph{deux corps différents} — à ne pas confondre avec deux forces qui s'équilibrent sur un même corps.
\end{itemize}
\end{aretenir}

\coursec{Dynamique du solide en rotation autour d'un axe fixe}

Dans la section précédente, nous avons établi le théorème du centre d'inertie $\sum \vec{F}_{\mathrm{ext}} = m\,\vec{a}_G$, qui gouverne le mouvement de \emph{translation} d'un système. Mais observe autour de toi : la roue de la mototaxi, la poulie du puits du village, le volant d'une machine à coudre, la turbine du barrage d'Edéa... Tous ces objets \emph{tournent} autour d'un axe fixe sans se déplacer globalement. Pour eux, le TCI ne suffit pas : il faut une loi adaptée à la rotation. C'est l'objet de cette section, l'une des plus importantes du programme de Terminale C.

\courssub{Description cinématique d'un solide en rotation autour d'un axe fixe}

\begin{definition}[Solide en rotation autour d'un axe fixe]
Un solide est en \cle{rotation autour d'un axe fixe} $\Delta$ si tous ses points décrivent des cercles centrés sur $\Delta$, dans des plans perpendiculaires à $\Delta$. La position du solide est repérée par l'\cle{angle} $\theta(t)$ que fait un plan lié au solide avec un plan de référence contenant $\Delta$.
\end{definition}

La conséquence fondamentale est la suivante : à un instant donné, \textbf{tous les points du solide balayent le même angle} $\theta$ pendant la même durée. Le solide possède donc \emph{une seule} coordonnée angulaire :

\begin{itemize}
\item une \cle{vitesse angulaire} commune : $\omega = \dot{\theta} = \dfrac{d\theta}{dt}$ (en $\mathrm{rad\cdot s^{-1}}$) ;
\item une \cle{accélération angulaire} commune : $\ddot{\theta} = \dot{\omega} = \dfrac{d\omega}{dt}$ (en $\mathrm{rad\cdot s^{-2}}$).
\end{itemize}

En revanche, les vitesses \emph{linéaires} des points diffèrent : un point $M$ situé à la distance $r$ de l'axe parcourt un arc de longueur $s = r\,\theta$, d'où, en dérivant par rapport au temps :

\begin{aretenir}[Vitesse et accélération d'un point du solide en rotation (Repère de Frenet)]
Pour un point $M$ situé à la distance $r$ de l'axe $\Delta$ :
\begin{align*}
\text{Vitesse :} \quad & \vec{v} = r\,\omega\,\vec{u}_t \quad \Rightarrow \quad v = r\,\omega \\
\text{Accélération tangentielle :} \quad & \vec{a}_t = r\,\ddot{\theta}\,\vec{u}_t \quad \Rightarrow \quad a_t = r\,\ddot{\theta} \\
\text{Accélération normale (centripète) :} \quad & \vec{a}_n = -r\,\omega^2\,\vec{u}_n \quad \Rightarrow \quad a_n = r\,\omega^2 = \dfrac{v^2}{r}
\end{align*}
Plus un point est éloigné de l'axe, plus sa vitesse et ses accélérations sont grandes. La périphérie d'une roue de mototaxi file bien plus vite que le voisinage du moyeu.
\end{aretenir}

\courssub{Moment d'une force par rapport à un axe}

\textbf{Intuition.} Pour ouvrir une porte, tu pousses loin des gonds, perpendiculairement à la porte : c'est plus efficace que de pousser près des gonds ou dans la direction de l'axe. L'efficacité d'une force à faire \emph{tourner} dépend donc de deux choses : son intensité $F$ et la distance de sa droite d'action à l'axe. Cette efficacité est mesurée par le \emph{moment}.

\begin{definition}[Moment d'une force par rapport à un axe]
Le \cle{moment} d'une force $\vec{F}$ par rapport à l'axe $\Delta$ est :
\[ \mathcal{M}_\Delta(\vec{F}) = \pm\, F \times d \]
où $d$ est le \cle{bras de levier} : distance \emph{perpendiculaire} de l'axe à la droite d'action de $\vec{F}$. Le signe est $+$ si $\vec{F}$ tend à faire tourner le solide dans le sens positif choisi (généralement trigonométrique), $-$ sinon. Unité : le newton-mètre ($\mathrm{N\cdot m}$).
\end{definition}

\begin{attention}[Bras de levier : la distance qui compte]
Le bras de levier $d$ est la distance de l'axe à la \textbf{droite d'action} de la force, \emph{pas} la distance de l'axe au point d'application. Deux cas à moment nul à connaître par c\oe ur :
\begin{itemize}
\item une force dont la droite d'action \textbf{coupe l'axe} ($d = 0$) : moment nul. C'est le cas du \textbf{poids} d'un solide dont le centre d'inertie $G$ est \emph{sur l'axe} : $\vec{P}$ appliqué en $G$ ne fait pas tourner le solide ;
\item une force \textbf{parallèle à l'axe} : moment nul (elle ne peut pas faire tourner autour de $\Delta$).
\end{itemize}
Oublier que le poids d'une poulie bien centrée a un moment nul est l'erreur la plus fréquente au Baccalauréat.
\end{attention}

\courssub{Moment d'inertie : l'inertie de rotation}

\textbf{Intuition.} Il est beaucoup plus facile de lancer en rotation une petite roue de vélo qu'un grand volant de machine, même de même masse : dans le volant, la masse est répartie \emph{loin de l'axe}. La résistance d'un solide à la mise en rotation dépend donc non seulement de sa masse, mais de la \emph{répartition} de cette masse autour de l'axe.

\begin{definition}[Moment d'inertie]
Le \cle{moment d'inertie} d'un solide par rapport à l'axe $\Delta$ est :
\[ J_\Delta = \sum_i m_i\, r_i^2 \quad \text{(système discret)} \qquad \text{ou} \qquad J_\Delta = \int r^2\,\mathrm{d}m \quad \text{(solide continu)} \]
où $m_i$ est la masse de chaque élément du solide et $r_i$ sa distance à l'axe. Unité SI : $\mathrm{kg\cdot m^2}$. Plus la masse est éloignée de l'axe, plus $J_\Delta$ est grand, plus le solide est difficile à mettre en rotation.
\end{definition}

\begin{aretenir}[Moments d'inertie usuels (admis)]
Pour un solide homogène de masse $M$, par rapport à son axe de symétrie :
\begin{center}
\begin{tabularx}{\linewidth}{|l|X|}
\hline
\rowcolor{popblueL} \textbf{Solide} & \textbf{Moment d'inertie} \\
\hline
Cerceau (ou anneau) de rayon $R$ & $J_\Delta = M R^2$ \\
\hline
Cylindre plein / disque de rayon $R$ & $J_\Delta = \dfrac{1}{2} M R^2$ \\
\hline
Tige de longueur $L$ (axe perpendiculaire en son milieu) & $J_\Delta = \dfrac{1}{12} M L^2$ \\
\hline
Sphère pleine de rayon $R$ & $J_\Delta = \dfrac{2}{5} M R^2$ \\
\hline
\end{tabularx}
\end{center}
Ces résultats sont \textbf{donnés} dans les énoncés du Baccalauréat : ils ne sont pas à démontrer, mais à savoir utiliser.
\end{aretenir}

\begin{popfigure}
\begin{center}
\begin{tikzpicture}[scale=1.0]
% Cerceau
\begin{scope}[xshift=-3.2cm]
\draw[line width=4pt, popblue] (0,0) circle (1.4);
\draw[popdark, dashed] (0,-1.9) -- (0,1.9);
\node[above, popdark] at (0,1.9) {$\Delta$};
\fill[popdark] (0,0) circle (0.05);
\node[below, align=center] at (0,-2.1) {\textbf{Cerceau} : masse au bord\\ $J = M R^2$};
\end{scope}
% Disque
\begin{scope}[xshift=3.2cm]
\fill[poporangeL] (0,0) circle (1.4);
\draw[line width=1.5pt, poporange] (0,0) circle (1.4);
\draw[popdark, dashed] (0,-1.9) -- (0,1.9);
\node[above, popdark] at (0,1.9) {$\Delta$};
\fill[popdark] (0,0) circle (0.05);
\node[below, align=center] at (0,-2.1) {\textbf{Disque} : masse répartie\\ $J = \dfrac{1}{2} M R^2$};
\end{scope}
\end{tikzpicture}
\end{center}
\legende{À masse $M$ et rayon $R$ égaux, le cerceau (toute la masse à la distance $R$) a un moment d'inertie \emph{double} de celui du disque (masse répartie du centre au bord). Il est donc deux fois plus difficile à lancer en rotation.}
\end{popfigure}

\courssub{Relation fondamentale de la dynamique de rotation}

Rappelons le théorème du centre d'inertie : $\sum \vec{F}_{\mathrm{ext}} = m\,\vec{a}_G$. Il relie la \emph{cause} (les forces) à l'\emph{effet} (l'accélération linéaire), la masse mesurant l'inertie de translation. Par analogie directe, pour la rotation autour d'un axe fixe, la cause devient le \emph{moment} des forces, l'effet devient l'\emph{accélération angulaire}, et l'inertie devient le \emph{moment d'inertie}.

\begin{propriete}[Relation fondamentale de la dynamique de rotation (RFD)]
Pour un solide en rotation autour d'un axe fixe $\Delta$, dans un référentiel galiléen :
\[ \sum \mathcal{M}_\Delta(\vec{F}_{\mathrm{ext}}) = J_\Delta\,\ddot{\theta} \]
La somme algébrique des moments des forces extérieures par rapport à $\Delta$ est égale au produit du moment d'inertie par l'accélération angulaire. Cette loi est \textbf{admise} ; elle est l'analogue rotationnel du TCI et s'applique avec la même méthode (bilan, projection, résolution).
\end{propriete}

\begin{attention}[Conditions d'application]
La RFD exige : un référentiel \textbf{galiléen}, un axe \textbf{fixe} dans ce référentiel, et des moments calculés par rapport à \emph{ce même axe}. Vérifie l'homogénéité : $[\mathcal{M}] = \mathrm{N\cdot m} = \mathrm{kg\cdot m^2 \cdot s^{-2}}$ et $[J\ddot{\theta}] = \mathrm{kg\cdot m^2 \cdot rad\cdot s^{-2}}$ ; le radian étant sans dimension, l'équation est bien homogène.
\end{attention}

\textbf{Cas particulier : l'équilibre de rotation.} Si $\sum \mathcal{M}_\Delta = 0$, alors $\ddot{\theta} = 0$ : la vitesse angulaire est \emph{constante}. Le solide est immobile ($\omega = 0$) ou en rotation uniforme. C'est l'analogue du principe d'inertie.

\begin{popfigure}
\begin{center}
\begin{tabularx}{\linewidth}{|l|X|X|}
\hline
\rowcolor{popblueL} \textbf{Grandeur} & \textbf{Translation} & \textbf{Rotation autour de $\Delta$} \\
\hline
Position & $x$ (m) & $\theta$ (rad) \\
\hline
Vitesse & $v = \dot{x}$ & $\omega = \dot{\theta}$ \\
\hline
Accélération & $a = \ddot{x}$ & $\ddot{\theta}$ \\
\hline
Inertie & masse $m$ (kg) & moment d'inertie $J_\Delta$ ($\mathrm{kg\cdot m^2}$) \\
\hline
Cause du mouvement & force $F$ (N) & moment $\mathcal{M}_\Delta = \pm F d$ ($\mathrm{N\cdot m}$) \\
\hline
Loi fondamentale & $\sum F = m\,a$ & $\sum \mathcal{M}_\Delta = J_\Delta\,\ddot{\theta}$ \\
\hline
Énergie cinétique & $\dfrac{1}{2} m v^2$ & $\dfrac{1}{2} J_\Delta\, \omega^2$ \\
\hline
\end{tabularx}
\end{center}
\legende{Dictionnaire translation--rotation : toute la dynamique de la rotation se déduit du TCI par simple substitution des grandeurs analogues.}
\end{popfigure}

\courssub{Application : la poulie et la masse suspendue}

Situation classique du Baccalauréat : une poulie assimilée à un disque homogène de masse $M$, de rayon $R$, mobile autour de son axe horizontal $\Delta$ ; un fil inextensible, non glissant sur la poulie, porte une masse $m$ qui descend en entraînant la rotation.

\begin{popfigure}
\begin{center}
\begin{tikzpicture}[scale=1.1]
% Support
\fill[pattern=north east lines, pattern color=popmuted] (-1.6,2.6) rectangle (1.6,2.9);
\draw[popdark] (-1.6,2.6) -- (1.6,2.6);
% Axe
\draw[popdark, line width=2pt] (0,2.6) -- (0,1.0);
% Poulie
\fill[popblueL] (0,0) circle (1.0);
\draw[popblue, line width=1.5pt] (0,0) circle (1.0);
\fill[popdark] (0,0) circle (0.06);
\node[popdark] at (0,0.35) {\small $G$};
\draw[popdark, ->] (0,0) -- (0.7,0.7) node[midway, above left] {\small $R$};
% Fil et masse
\draw[popdark, line width=1pt] (1.0,0) -- (1.0,-2.2);
\draw[fill=poporangeL, draw=poporange, line width=1pt] (0.65,-2.9) rectangle (1.35,-2.2);
\node at (1.0,-2.55) {$m$};
% Forces sur la masse
\draw[->, poppink, line width=1.3pt] (1.0,-2.9) -- (1.0,-3.9) node[below] {$\vec{P} = m\vec{g}$};
\draw[->, popgreen, line width=1.3pt] (1.0,-2.2) -- (1.0,-1.3) node[right] {$\vec{T}$};
% Force sur la poulie (tension tangentielle)
\draw[->, popgreen, line width=1.3pt] (1.0,-0.05) -- (1.0,-1.0) node[right] {$\vec{T}\,'$};
% Sens de rotation
\draw[->, poppurple, line width=1.2pt] (-0.75,0.65) arc (140:40:0.95) node[midway, above left] {\small $\ddot{\theta}$};
% Réaction et poids de la poulie
\draw[->, popmuted, line width=1.2pt] (0,0) -- (0,0.85) node[left] {\small $\vec{R}_{\mathrm{axe}}$};
\draw[->, popmuted, line width=1.2pt] (-0.12,0) -- (-0.12,-0.85) node[left] {\small $\vec{P}_{\mathrm{poulie}}$};
\end{tikzpicture}
\end{center}
\legende{Poulie--masse : le poids $\vec{P}_{\mathrm{poulie}}$ et la réaction $\vec{R}_{\mathrm{axe}}$ passent par l'axe (moments nuls) ; seule la tension $\vec{T}\,'$ exerce un moment moteur $T' R$. La masse $m$ est soumise à son poids et à la tension $\vec{T}$.}
\end{popfigure}

\begin{methode}[Coupler TCI et RFD avec la condition de non-glissement]
\begin{enumerate}
\item \textbf{Isoler la masse $m$} (translation) : bilan des forces ($\vec{P}$, $\vec{T}$), puis TCI projeté sur la verticale descendante : $m g - T = m a$.
\item \textbf{Isoler la poulie} (rotation) : le poids et la réaction de l'axe ont des moments nuls (droites d'action passant par $\Delta$) ; seule la tension $\vec{T}\,'$ du fil agit, avec le bras de levier $R$. RFD : $T' R = J_\Delta\,\ddot{\theta}$.
\item \textbf{Liaison du fil} : fil inextensible et masse du fil négligée $\Rightarrow T = T'$.
\item \textbf{Non-glissement} du fil sur la poulie : le point de contact a même vitesse que le fil, d'où $v = R\omega$ et en dérivant :
\[ a = R\,\ddot{\theta} \]
\item \textbf{Résoudre} le système de trois équations à trois inconnues $(a, \ddot{\theta}, T)$.
\end{enumerate}
\end{methode}

\begin{exoresolu}[Poulie entraînée par une masse suspendue]
Une poulie assimilée à un disque homogène de masse $M = \SI{2}{kg}$ et de rayon $R = \SI{20}{cm}$ tourne sans frottement autour de son axe horizontal. Un fil inextensible enroulé sur la poulie porte une masse $m = \SI{1}{kg}$ lâchée sans vitesse initiale. On prend $g = \SI{10}{m\cdot s^{-2}}$. Déterminer l'accélération $a$ de la masse, l'accélération angulaire de la poulie et la tension du fil.
\tcblower
\textbf{Moment d'inertie de la poulie} (disque) :
\[ J_\Delta = \frac{1}{2} M R^2 = \frac{1}{2} \times 2 \times (0,20)^2 = \SI{4e-2}{kg\cdot m^2} \]
soit $J_\Delta = 4\times 10^{-2}\ \mathrm{kg\cdot m^2}$.

\textbf{Équations.} TCI sur la masse (axe vertical descendant) : $m g - T = m a$. RFD sur la poulie : $T R = J_\Delta\,\ddot{\theta}$. Non-glissement : $\ddot{\theta} = a/R$, donc $T = \dfrac{J_\Delta\, a}{R^2}$. En reportant dans le TCI :
\[ m g = m a + \frac{J_\Delta}{R^2}\, a \quad\Longrightarrow\quad a = \frac{m g}{m + \dfrac{J_\Delta}{R^2}} \]

\textbf{Application numérique} : $\dfrac{J_\Delta}{R^2} = \dfrac{4\times 10^{-2}}{(0,20)^2} = \SI{1}{kg}$, d'où
\[ a = \frac{1 \times 10}{1 + 1} = \SI{5}{m\cdot s^{-2}} \]
\textbf{Accélération angulaire} : $\ddot{\theta} = \dfrac{a}{R} = \dfrac{5}{0,20} = \SI{25}{rad\cdot s^{-2}}$.

\textbf{Tension} : $T = m(g - a) = 1\times(10-5) = \SI{5}{N}$.

\textbf{Remarque physique} : l'accélération ($\SI{5}{m\cdot s^{-2}}$) est inférieure à $g$ car une partie du travail du poids sert à mettre la poulie en rotation. Si la poulie était de masse négligeable ($J_\Delta \to 0$), on retrouverait la chute libre $a = g$.
\end{exoresolu}

\begin{attention}[Pièges fréquents dans les exercices de poulie]
\begin{itemize}
\item Ne jamais écrire $T = m g$ : la masse est \emph{accélérée}, donc $T \neq mg$ (ce serait vrai uniquement à vitesse constante).
\item Ne pas oublier le facteur $\tfrac{1}{2}$ dans $J = \tfrac{1}{2}MR^2$ pour un disque : le confondre avec le cerceau $J = MR^2$ fausse tout le calcul.
\item Convertir le rayon en \textbf{mètres} avant tout calcul.
\item La condition $a = R\ddot{\theta}$ n'est valable que s'il y a \textbf{non-glissement} du fil : le vérifier dans l'énoncé.
\end{itemize}
\end{attention}

\courssub{Autres applications : volant d'inertie et treuil}

\begin{savaistu}[Le volant d'inertie, réservoir d'énergie de rotation]
Un \cle{volant d'inertie} est une roue de grand moment d'inertie $J_\Delta$ qu'on fait tourner à vitesse angulaire $\omega$ élevée : il stocke l'énergie cinétique $E_c = \tfrac{1}{2}J_\Delta\,\omega^2$ et la restitue en régulant les à-coups. On en trouve dans les moteurs (lissage du mouvement entre les explosions), les groupes électrogènes qui suppléent les délestages d'ENEO, et même dans des systèmes modernes de stockage d'électricité pour l'énergie solaire en zone rurale. Un grand $J_\Delta$ donne une grande « inertie » : d'après la RFD, à moment résistant donné, $\ddot{\theta}$ est faible, donc la vitesse varie peu.
\end{savaistu}

\begin{exemplebox}[Treuil de puits]
Un treuil est un cylindre de rayon $R$ sur lequel s'enroule une corde ; on le man\oe uvre par une manivelle de longueur $\ell$. Pour soulever un seau de masse $m$ à vitesse constante ($\ddot{\theta} = 0$), la RFD donne l'équilibre des moments : $F\,\ell = m g R$ (en négligeant les frottements), soit $F = m g \dfrac{R}{\ell}$. Avec $\ell = 5R$, l'effort est cinq fois plus faible que le poids du seau : c'est le principe du \emph{démultiplicateur d'effort}, utilisé dans les puits et les treuils de levage des chantiers de Douala.
\end{exemplebox}

\begin{exercice}[Pour t'entraîner]
Un cerceau de masse $M = \SI{0,5}{kg}$ et de rayon $R = \SI{10}{cm}$, mobile sans frottement autour de son axe horizontal, est entraîné par un fil enroulé sur sa jante auquel on suspend une masse $m = \SI{200}{g}$. Calculer l'accélération de la masse, l'accélération angulaire du cerceau et la tension du fil ($g = \SI{10}{m\cdot s^{-2}}$). Comparer avec le cas où le cerceau serait remplacé par un disque de mêmes masse et rayon.
\end{exercice}

\begin{corrige}[Exercice]
Cerceau : $J_\Delta = M R^2 = 0,5\times(0,10)^2 = 5\times 10^{-3}\ \mathrm{kg\cdot m^2}$, d'où $\dfrac{J_\Delta}{R^2} = M = \SI{0,5}{kg}$ et
\[ a = \frac{m g}{m + M} = \frac{0,2\times 10}{0,2 + 0,5} \approx \SI{2,9}{m\cdot s^{-2}} \]
$\ddot{\theta} = a/R \approx \SI{29}{rad\cdot s^{-2}}$ ; $T = m(g-a) \approx 0,2\times 7,1 \approx \SI{1,4}{N}$.

Disque : $\dfrac{J_\Delta}{R^2} = \dfrac{M}{2} = \SI{0,25}{kg}$, d'où $a = \dfrac{2}{0,45} \approx \SI{4,4}{m\cdot s^{-2}}$ : le disque, de moment d'inertie moitié moindre, s'accélère plus vite, conformément à l'intuition.
\end{corrige}

\begin{aretenir}[L'essentiel de la section]
\begin{itemize}
\item Solide en rotation autour d'un axe fixe : un seul angle $\theta(t)$, vitesse d'un point $v = r\omega$, accélérations $a_t = r\ddot{\theta}$, $a_n = r\omega^2$.
\item Moment d'une force : $\mathcal{M}_\Delta = \pm F d$ ($d$ = bras de levier) ; nul si la droite d'action coupe l'axe ou si la force est parallèle à l'axe.
\item Moment d'inertie $J_\Delta = \sum m_i r_i^2$ ($\mathrm{kg\cdot m^2}$) : mesure de l'inertie de rotation.
\item RFD : $\sum \mathcal{M}_\Delta(\vec{F}_{\mathrm{ext}}) = J_\Delta\,\ddot{\theta}$ ; équilibre de rotation si $\sum \mathcal{M}_\Delta = 0$.
\item Systèmes couplés : TCI pour la translation, RFD pour la rotation, reliés par le non-glissement $a = R\ddot{\theta}$.
\end{itemize}
\end{aretenir}

\coursec{Théorème de l'énergie cinétique}

Dans la section précédente, nous avons appris à mettre en équation le mouvement d'un solide en rotation autour d'un axe fixe grâce à la relation fondamentale de la dynamique : $\sum \mathcal{M}_\Delta = J_\Delta \ddot{\theta}$. Cette méthode, dite \emph{dynamique}, est puissante mais exige de résoudre des équations différentielles et de connaître l'évolution du système \emph{instant par instant}. Or, dans de nombreuses situations concrètes — quelle distance faut-il à une mototaxi pour s'arrêter ? quelle vitesse atteint un wagonnet en bas d'une descente ? —, seules les \textbf{positions de départ et d'arrivée} nous intéressent, pas la chronologie détaillée. Il existe alors un raccourci remarquable : le \cle{théorème de l'énergie cinétique}, qui relie directement la variation de vitesse au \emph{travail} des forces, sans passer par le temps. C'est l'objet de cette section.

\courssub{Énergie cinétique : translation et rotation}

\textbf{Observation.} Un camion lancé à \SI{90}{km/h} sur la route Yaoundé--Douala est bien plus difficile à arrêter qu'un vélo à \SI{20}{km/h} : il possède une grande « réserve » liée à son mouvement. Cette réserve, qui dépend de la masse et de la vitesse, est appelée \cle{énergie cinétique}.

\begin{definition}[Énergie cinétique]
L'\cle{énergie cinétique} d'un système est l'énergie qu'il possède du fait de son mouvement.
\begin{itemize}
\item Pour un solide en \textbf{translation} (ou un point matériel) de masse $m$ dont le centre d'inertie $G$ a la vitesse $v_G$ :
\[ E_c = \frac{1}{2}\, m\, v_G^2 \]
\item Pour un solide en \textbf{rotation} autour d'un axe fixe $\Delta$, de moment d'inertie $J_\Delta$ et de vitesse angulaire $\omega = \dot{\theta}$ :
\[ E_c = \frac{1}{2}\, J_\Delta\, \omega^2 \]
\end{itemize}
$E_c$ s'exprime en \cle{joules} (J), avec $m$ en kg, $v_G$ en m/s, $J_\Delta$ en $\mathrm{kg\,m^2}$ et $\omega$ en rad/s.
\end{definition}

Vérifions l'homogénéité : $[\tfrac12 m v^2] = \mathrm{kg \cdot m^2\,s^{-2}} = \mathrm{N \cdot m} = \mathrm{J}$. De même, $[\tfrac12 J\omega^2] = \mathrm{kg\,m^2 \cdot s^{-2}} = \mathrm{J}$. La structure des deux formules est identique : la masse $m$ (inertie de translation) est remplacée par le moment d'inertie $J_\Delta$ (inertie de rotation), et la vitesse linéaire $v$ par la vitesse angulaire $\omega$. Cette analogie translation/rotation, déjà rencontrée pour les lois de Newton, se poursuit ici.

\begin{attention}[Pièges de calcul]
Deux erreurs classiques : oublier que la vitesse intervient \textbf{au carré} (doubler $v$ quadruple $E_c$ !), et utiliser une vitesse en km/h sans la convertir : $v\ (\mathrm{m/s}) = v\ (\mathrm{km/h}) / 3{,}6$. Ainsi \SI{90}{km/h} $=$ \SI{25}{m/s}.
\end{attention}

\courssub{Rappels : travail d'une force et puissance}

Le théorème de l'énergie cinétique fait le bilan des \emph{travaux} des forces. Rappelons les expressions essentielles.

\begin{aretenir}[Travaux à connaître]
\begin{itemize}
\item \textbf{Force constante} $\vec{F}$ dont le point d'application se déplace de $A$ vers $B$ :
\[ W_{AB}(\vec{F}) = \vec{F} \cdot \overrightarrow{AB} = F \times AB \times \cos\alpha \]
où $\alpha$ est l'angle entre $\vec{F}$ et $\overrightarrow{AB}$. Le travail est \emph{moteur} ($W>0$) si $\alpha < 90°$, \emph{résistant} ($W<0$) si $\alpha > 90°$, \emph{nul} si $\vec{F} \perp \overrightarrow{AB}$.
\item \textbf{Travail du poids} entre deux points d'altitudes $z_A$ et $z_B$ (axe $Oz$ vertical ascendant) :
\[ W_{AB}(\vec{P}) = m\,g\,(z_A - z_B) \]
Il ne dépend que de la dénivellation, pas du chemin suivi : le poids est une force \emph{conservative}.
\item \textbf{Travail d'un couple} de moment $\mathcal{M}$ constant appliqué à un solide tournant de $\Delta\theta$ :
\[ W = \mathcal{M}\,\Delta\theta \qquad (\Delta\theta \text{ en radians}) \]
\end{itemize}
\end{aretenir}

\begin{definition}[Puissance]
La \cle{puissance} d'une force est le travail qu'elle fournit par unité de temps : $P = \dfrac{dW}{dt}$, en watts (W $=$ J/s).
\begin{itemize}
\item En translation : $P = \vec{F}\cdot\vec{v} = F\,v\cos\alpha$.
\item En rotation : $P = \mathcal{M}\,\omega$.
\end{itemize}
\end{definition}

\courssub{Énoncé et démonstration du théorème de l'énergie cinétique}

\begin{propriete}[Théorème de l'énergie cinétique (TEC)]
Dans un référentiel galiléen, la variation d'énergie cinétique d'un système entre deux positions $A$ et $B$ est égale à la somme des travaux de \textbf{toutes} les forces extérieures appliquées au système entre $A$ et $B$ :
\[ \Delta E_c = E_c(B) - E_c(A) = \sum W_{AB}(\vec{F}_{\mathrm{ext}}) \]
En translation : $\dfrac{1}{2}m v_B^2 - \dfrac{1}{2}m v_A^2 = \sum W_{AB}(\vec{F}_{\mathrm{ext}})$. \quad
En rotation autour d'un axe fixe : $\dfrac{1}{2}J_\Delta \omega_B^2 - \dfrac{1}{2}J_\Delta \omega_A^2 = \sum W(\text{couples}) = \sum \mathcal{M}\,\Delta\theta$.
\end{propriete}

\begin{experience}[Démonstration guidée du TEC en translation (à partir du TCI)]
Considérons un point matériel de masse $m$ se déplaçant de $A$ vers $B$ dans un référentiel galiléen.
\begin{enumerate}
\item \textbf{Point de départ : le TCI.} La deuxième loi de Newton s'écrit :
\[ m\,\vec{a} = \sum \vec{F}_{\mathrm{ext}} \]
\item \textbf{Multiplication scalaire par la vitesse.} Multiplions scalairement les deux membres par $\vec{v}$ :
\[ m\,\vec{a}\cdot\vec{v} = \sum \vec{F}_{\mathrm{ext}}\cdot\vec{v} \]
Le membre de droite est la somme des \emph{puissances} des forces : $\sum P$.
\item \textbf{Transformation du membre de gauche.} Remarquons que
\[ \vec{a}\cdot\vec{v} = \frac{d\vec{v}}{dt}\cdot\vec{v} = \frac{1}{2}\frac{d}{dt}\left(\vec{v}\cdot\vec{v}\right) = \frac{1}{2}\frac{d(v^2)}{dt} \]
donc
\[ m\,\vec{a}\cdot\vec{v} = \frac{d}{dt}\left(\frac{1}{2}m v^2\right) = \frac{dE_c}{dt} \]
\item \textbf{Intégration entre $A$ et $B$.} En intégrant de l'instant $t_A$ à l'instant $t_B$ :
\[ \int_{t_A}^{t_B}\frac{dE_c}{dt}\,dt = \int_{t_A}^{t_B} \sum P\,dt \quad\Longrightarrow\quad E_c(B)-E_c(A) = \sum W_{AB}(\vec{F}_{\mathrm{ext}}) \]
\end{enumerate}
Le théorème est ainsi une conséquence directe du TCI : il en est la forme \emph{intégrée par rapport à l'espace}. En rotation, la même démarche à partir de $J_\Delta\ddot{\theta} = \sum\mathcal{M}$, multipliée par $\dot{\theta}$ puis intégrée, donne $\Delta(\tfrac12 J_\Delta\omega^2) = \sum \mathcal{M}\Delta\theta$.
\end{experience}

\begin{methode}[Appliquer le TEC : quand et comment]
Le TEC est l'outil idéal lorsqu'on cherche une \textbf{vitesse} ou une \textbf{distance} \emph{sans connaître le temps} ni la durée du mouvement.
\begin{enumerate}
\item Définir le système, le référentiel galiléen, et les positions $A$ (départ) et $B$ (arrivée).
\item Faire le \textbf{bilan exhaustif} des forces extérieures (schéma !).
\item Calculer le travail de \emph{chaque} force entre $A$ et $B$, avec son signe (poids : $mg(z_A-z_B)$ ; frottements : $-f \times AB$ ; réaction normale : travail nul).
\item Écrire $\frac12 m v_B^2 - \frac12 m v_A^2 = \sum W$ et résoudre pour l'inconnue.
\item Vérifier la cohérence du signe et de l'unité du résultat.
\end{enumerate}
\end{methode}

\begin{attention}[Erreurs fréquentes]
\begin{itemize}
\item \textbf{Oublier le travail des frottements}, qui est toujours \emph{négatif} : $W(\vec{f}) = -f \times d$ si $\vec{f}$ s'oppose au déplacement. L'omettre conduit à des distances d'arrêt nulles ou des vitesses surestimées.
\item \textbf{Se tromper de signe pour le poids} : si le système \emph{descend} ($z_A > z_B$), le poids est moteur, $W = +mg(z_A - z_B) > 0$ ; s'il \emph{monte}, le poids est résistant, $W < 0$. Retenir : $W(\vec P) = mg(z_A - z_B)$, avec les \emph{altitudes}, pas la distance parcourue.
\item \textbf{Compter le travail de la réaction normale} sur un support : elle est perpendiculaire au déplacement, donc $W(\vec{R}_N) = 0$.
\end{itemize}
\end{attention}

\courssub{Applications}

\textbf{1. Distance de freinage d'un véhicule.} Une voiture de masse $m$ roule à la vitesse $v$ sur une route horizontale ; les freins exercent une force constante $\vec{f}$ opposée au mouvement. Entre le début du freinage ($A$, vitesse $v$) et l'arrêt ($B$, vitesse nulle), le poids et la réaction de la route ne travaillent pas (horizontale). Le TEC donne :
\[ 0 - \frac{1}{2}m v^2 = -f\,d \quad\Longrightarrow\quad \boxed{d = \frac{m\,v^2}{2f}} \]
La distance d'arrêt est proportionnelle au \textbf{carré} de la vitesse.

\begin{exemplebox}[Distance d'arrêt d'une voiture]
Une voiture de $m = \SI{1000}{kg}$ lancée à $v = \SI{20}{m/s}$ (soit \SI{72}{km/h}) freine avec une force $f = \SI{5000}{N}$. Alors :
\[ d = \frac{1000 \times 20^2}{2 \times 5000} = \frac{4\times 10^5}{10^4} = \SI{40}{m} \]
À \SI{40}{m/s} (\SI{144}{km/h}), il faudrait $d = \SI{160}{m}$ : la distance est \emph{quadruplée} !
\end{exemplebox}

\begin{popfigure}
\begin{center}
\begin{tikzpicture}
\begin{axis}[
width=0.85\linewidth, height=6.5cm,
xlabel={Vitesse $v$ (m/s)}, ylabel={Distance d'arrêt $d$ (m)},
xmin=0, xmax=42, ymin=0, ymax=180,
grid=major, grid style={popline!50},
legend pos=north west, legend style={font=\small},
xtick={0,10,20,30,40}, ytick={0,40,80,120,160},
]
\addplot[popcurve,domain=0:42,samples=200]{0.1*x^2};
\addlegendentry{$d = \dfrac{m v^2}{2f}$ \ ($m=1000$ kg, $f=5000$ N)}
\addplot[only marks, mark=*, mark size=2pt, poporange] coordinates {(20,40) (40,160)};
\node[poporange, anchor=west, font=\small] at (axis cs:21,32) {$(20\ ;\ 40)$};
\node[poporange, anchor=west, font=\small] at (axis cs:33,150) {$(40\ ;\ 160)$};
\end{axis}
\end{tikzpicture}
\end{center}
\legende{Distance d'arrêt en fonction de la vitesse : la loi en $v^2$ fait exploser la distance d'arrêt dès que la vitesse augmente.}
\end{popfigure}

\begin{savaistu}[Sécurité routière : pourquoi doubler la vitesse quadruple le danger]
Puisque $d = \dfrac{mv^2}{2f}$, la distance de freinage est proportionnelle à $v^2$ : \textbf{doubler la vitesse quadruple la distance d'arrêt}. Il faut y ajouter la distance parcourue pendant le temps de réaction du conducteur (environ 1 s). C'est pourquoi, sur les axes routiers du Cameroun, les campagnes de sécurité insistent sur la limitation de vitesse : à \SI{110}{km/h} au lieu de \SI{90}{km/h}, la distance de freinage augmente de près de \SI{50}{\%}, et l'énergie cinétique à dissiper lors d'un choc — donc sa violence — augmente dans la même proportion.
\end{savaistu}

\textbf{2. Vitesse en bas d'un plan incliné.} Un solide de masse $m$ glisse sans frottements sur un plan incliné, d'une hauteur $h$. Seul le poids travaille (la réaction est perpendiculaire au déplacement) :
\[ \frac{1}{2}m v_B^2 - 0 = mgh \quad\Longrightarrow\quad v_B = \sqrt{2gh} \]
Résultat remarquable : la vitesse ne dépend ni de la masse, ni de l'angle du plan, mais seulement de la dénivellation $h$ — comme en chute libre. Avec frottements $f$ le long du plan de longueur $L$, on aurait $\tfrac12 m v_B^2 = mgh - fL$.

\textbf{3. Volant soumis à un couple moteur.} Un volant de moment d'inertie $J_\Delta$, initialement au repos, est soumis à un couple moteur de moment $\mathcal{M}_m$ (algébrique, positif) et à un couple résistant $\mathcal{M}_r$ (algébrique, négatif). Après $N$ tours ($\Delta\theta = 2\pi N$) :
\[ \frac{1}{2}J_\Delta\omega^2 - 0 = (\mathcal{M}_m + \mathcal{M}_r)\, 2\pi N \quad\Longrightarrow\quad \omega = \sqrt{\frac{4\pi N (\mathcal{M}_m + \mathcal{M}_r)}{J_\Delta}} \]
(On note que $\mathcal{M}_m + \mathcal{M}_r = \mathcal{M}_m - |\mathcal{M}_r|$ car $\mathcal{M}_r < 0$.)

\begin{popfigure}
\begin{center}
\begin{tikzpicture}[scale=1.0]
% Axe
\fill[poppurple] (0,0) circle (0.06);
\node[below, font=\small] at (0,-0.1) {$\Delta$};
% Volant
\draw[very thick, popblue, fill=popblueL] (0,0) circle (1.5);
\draw[popblue] (0,0) -- (1.06,1.06);
\draw[popblue] (0,0) -- (-1.06,-1.06);
\fill[popblue] (1.06,1.06) circle (0.09);
\fill[popblue] (-1.06,-1.06) circle (0.09);
% Couple moteur (flèche courbe)
\draw[-{Stealth[length=3mm]}, very thick, popgreen] (2.1,-0.9) arc[start angle=-30, end angle=60, radius=2.28];
\node[popgreen, font=\small, align=center] at (3.1,0.9) {couple moteur\\ $\mathcal{M}_m > 0$};
% Couple résistant
\draw[-{Stealth[length=3mm]}, very thick, poporange] (-2.1,0.9) arc[start angle=150, end angle=240, radius=2.28];
\node[poporange, font=\small, align=center] at (-3.2,-1.0) {couple résistant\\ $\mathcal{M}_r < 0$};
% Rotation
\draw[-{Stealth[length=2.5mm]}, thick, popdark] (0.5,1.9) arc[start angle=75, end angle=15, radius=1.97];
\node[popdark, font=\small] at (1.15,2.15) {$\omega$};
\end{tikzpicture}
\end{center}
\legende{Volant de moment d'inertie $J_\Delta$ tournant autour de l'axe fixe $\Delta$, soumis à un couple moteur $\mathcal{M}_m$ et à un couple résistant $\mathcal{M}_r$. Le TEC en rotation s'écrit $\tfrac12 J_\Delta \omega^2 = (\mathcal{M}_m + \mathcal{M}_r)\,\Delta\theta$ (avec $\mathcal{M}_r$ négatif).}
\end{popfigure}

\begin{exoresolu}[Lancement d'un volant]
Un volant de machine, de moment d'inertie $J_\Delta = \SI{20}{kg.m^2}$, initialement au repos, est soumis à un couple moteur constant $\mathcal{M}_m = \SI{50}{N.m}$ et à un couple de frottement $\mathcal{M}_r = \SI{-10}{N.m}$. Déterminer sa vitesse angulaire après $N = 10$ tours, puis la puissance moyenne du moteur pendant cette phase si celle-ci a duré $\Delta t = \SI{8}{s}$.
\tcblower
\textbf{Solution.} Le système est le volant, en rotation autour de l'axe fixe $\Delta$, dans le référentiel galiléen de l'atelier. Le poids et la réaction de l'axe ont un moment nul par rapport à $\Delta$ : seuls les couples travaillent.

Angle balayé : $\Delta\theta = 2\pi N = 20\pi \approx \SI{62,8}{rad}$.

Application du TEC en rotation entre le départ ($\omega_A = 0$) et l'état final ($\omega_B = \omega$) :
\[ \frac{1}{2}J_\Delta \omega^2 - 0 = (\mathcal{M}_m + \mathcal{M}_r)\,\Delta\theta = (50 - 10)\times 62{,}8 = \SI{2513}{J} \]
\[ \omega = \sqrt{\frac{2 \times 2513}{20}} = \sqrt{251{,}3} \approx \SI{15,9}{rad/s} \]
soit environ $152$ tours/min.

Puissance moyenne du couple moteur : le travail moteur vaut $W_m = \mathcal{M}_m\,\Delta\theta = 50 \times 62{,}8 = \SI{3142}{J}$, d'où
\[ P_{\text{moy}} = \frac{W_m}{\Delta t} = \frac{3142}{8} \approx \SI{393}{W} \]
Vérification énergétique : sur les \SI{3142}{J} fournis par le moteur, $\mathcal{M}_r\Delta\theta = \SI{628}{J}$ sont dissipés par frottements et $\tfrac12 J_\Delta\omega^2 = \SI{2513}{J}$ sont stockés en énergie cinétique : $3142 = 628 + 2513$. Le bilan est cohérent.
\end{exoresolu}

\begin{aretenir}[L'essentiel]
\begin{itemize}
\item Énergie cinétique : $E_c = \tfrac12 m v_G^2$ (translation), $E_c = \tfrac12 J_\Delta \omega^2$ (rotation autour d'un axe fixe).
\item Théorème de l'énergie cinétique : $\Delta E_c = \sum W_{AB}(\vec{F}_{\mathrm{ext}})$, conséquence intégrée du TCI.
\item Le TEC relie vitesse et déplacement \emph{sans le temps} : c'est l'outil privilégié pour les distances d'arrêt, les vitesses en fonction de l'altitude, les vitesses angulaires après un nombre de tours donné.
\item Puissance : $P = \vec F\cdot\vec v$ (translation), $P = \mathcal{M}\omega$ (rotation), et $P = \dfrac{dE_c}{dt}$ pour le système.
\end{itemize}
\end{aretenir}

Ce théorème ouvre naturellement la voie à la section suivante : en identifiant les forces dont le travail ne dépend pas du chemin suivi (comme le poids), nous pourrons définir l'\emph{énergie potentielle}, puis l'\emph{énergie mécanique}, et énoncer sa conservation pour les systèmes isolés — un outil encore plus puissant pour étudier les mouvements.

\coursec{Énergie potentielle, énergie mécanique et conservation}

Le théorème de l'énergie cinétique, étudié dans la section précédente, nous a appris à relier la variation d'énergie cinétique d'un système au travail de \emph{toutes} les forces appliquées. C'est un outil puissant, mais il exige de calculer le travail de chaque force, une par une. Or certaines forces — le poids, la force élastique — possèdent une propriété remarquable : leur travail ne dépend pas du chemin suivi. Pour ces forces, on peut « stocker » le travail sous forme d'une nouvelle grandeur, l'\cle{énergie potentielle}, et faire apparaître une quantité qui se conserve : l'\cle{énergie mécanique}. C'est l'objet de cette dernière section, qui couronne tout le chapitre et fournit la méthode de résolution la plus rapide dans de nombreux problèmes du Baccalauréat.

\courssub{Forces conservatives et énergie potentielle}

\textbf{Intuition.} Lâche une bille d'une hauteur $h$ : elle accélère. D'où vient son énergie cinétique ? Elle était « en réserve » dans la position de la bille, du fait de la présence du poids. De même, un ressort comprimé peut projeter une bille : l'énergie était stockée dans la déformation. Ces réserves d'énergie liées à la \emph{position} du système sont des énergies potentielles.

\begin{definition}[Force conservative]
Une force est dite \cle{conservative} si son travail entre deux points $A$ et $B$ ne dépend pas du chemin suivi, mais seulement des positions de $A$ et $B$. On peut alors lui associer une \cle{énergie potentielle} $E_p$ telle que :
\[
W_{A\to B}(\vec{F}) = -\Delta E_p = E_p(A) - E_p(B)
\]
Le travail d'une force conservative fait \emph{diminuer} l'énergie potentielle. Une force dont le travail dépend du chemin (frottements, force de traction d'un moteur) est dite \cle{non conservative}.
\end{definition}

\textbf{Énergie potentielle de pesanteur.} Le travail du poids entre $A$ et $B$ vaut $W_{A\to B}(\vec{P}) = mg(z_A - z_B)$, où l'axe $(Oz)$ est vertical ascendant. Ce travail ne dépend que des altitudes : le poids est conservatif, et :
\[
E_{p} = mgz + \text{constante}
\]
La constante se fixe en choisissant une \cle{référence} : on pose $E_p = 0$ à une altitude arbitraire (souvent le sol ou le point le plus bas du mouvement). Alors $E_p = mgz$.

\begin{attention}[La référence est libre, le signe non]
L'énergie potentielle est définie à une constante additive près : seules ses \emph{variations} ont un sens physique. Mais une fois la référence choisie, il faut s'y tenir pendant tout l'exercice ! Si l'axe $z$ est orienté vers le haut, un point situé \emph{sous} la référence a $E_p < 0$ : c'est normal. Erreur classique : écrire $E_p = mgh$ avec $h$ « distance au sol » alors que la référence n'est pas le sol.
\end{attention}

\textbf{Énergie potentielle élastique (complément utile au Bac).} La force de rappel d'un ressort $\vec{F} = -kx\,\vec{u}_x$ est conservative. L'énergie potentielle élastique associée, avec la référence $E_p = 0$ à la position d'équilibre ($x = 0$), s'écrit :
\[
E_{p,\text{él}} = \dfrac{1}{2}kx^2
\]
où $k$ est la raideur (en $\mathrm{N\,m^{-1}}$) et $x$ l'allongement algébrique. Cette expression n'est pas exigible en soi en Terminale C, mais elle apparaît fréquemment dans les sujets de Bac comme donnée : sache l'utiliser sans hésiter.

\courssub{Énergie mécanique : définition et théorème}

\begin{definition}[Énergie mécanique]
L'\cle{énergie mécanique} d'un système est la somme de son énergie cinétique et de son (ou ses) énergie(s) potentielle(s) :
\[
E_m = E_c + E_p
\]
Pour un solide en translation soumis à son poids : $E_m = \dfrac{1}{2}mv^2 + mgz$. Elle s'exprime en joules (J).
\end{definition}

Comment évolue $E_m$ au cours du mouvement ? Répondons rigoureusement à partir du théorème de l'énergie cinétique.

\begin{propriete}[Théorème de l'énergie mécanique — démonstration guidée]
Entre deux points $A$ et $B$, la variation d'énergie mécanique est égale au travail des forces \textbf{non conservatives} :
\[
\Delta E_m = E_m(B) - E_m(A) = W_{A\to B}(\vec{F}_{\text{non conservatives}})
\]
\textbf{Démonstration.} Séparons les forces appliquées au système en deux familles : conservatives (indice $c$) et non conservatives (indice $nc$). Le théorème de l'énergie cinétique s'écrit :
\[
\Delta E_c = W(\vec{F}_c) + W(\vec{F}_{nc})
\]
Or par définition de l'énergie potentielle, $W(\vec{F}_c) = -\Delta E_p$. Donc :
\[
\Delta E_c = -\Delta E_p + W(\vec{F}_{nc})
\quad\Longrightarrow\quad
\Delta E_c + \Delta E_p = W(\vec{F}_{nc})
\quad\Longrightarrow\quad
\Delta E_m = W(\vec{F}_{nc}) \qquad \blacksquare
\]
Cette démonstration est exigible : elle tient en trois lignes et repose uniquement sur le TEC et la définition de $E_p$.
\end{propriete}

\begin{propriete}[Conservation de l'énergie mécanique]
Si un système n'est soumis qu'à des forces conservatives, ou à des forces dont le travail est nul (réaction normale d'un support sans frottement, tension d'un fil inextensible), alors son énergie mécanique se conserve :
\[
E_m = \text{constante} \qquad \text{c'est-à-dire} \qquad E_c(A) + E_p(A) = E_c(B) + E_p(B)
\]
On dit que le système est \cle{conservatif}. Il y a alors simple \emph{transfert} d'énergie entre forme cinétique et forme potentielle, sans création ni disparition.
\end{propriete}

\begin{attention}[Le piège n°1 du chapitre]
Ne jamais écrire « $E_m$ se conserve » par réflexe ! Avant toute application, fais le \textbf{bilan des forces} et demande-toi : quelles forces \emph{travaillent} ? Si l'énoncé mentionne des frottements, une piste rugueuse, une résistance de l'air, alors $E_m$ \textbf{ne se conserve pas} : il faut utiliser $\Delta E_m = W(\vec{f})$. Une réaction normale à un support, elle, ne travaille pas (force perpendiculaire au déplacement) : elle ne rompt pas la conservation.
\end{attention}

\courssub{Applications de la conservation}

\textbf{Chute libre.} Un objet de masse $m$ lâché sans vitesse initiale d'une hauteur $h$ n'est soumis qu'à son poids (frottements de l'air négligés). La conservation de $E_m$ entre le point de départ ($v = 0$, $z = h$) et le sol ($z = 0$) donne :
\[
mgh = \dfrac{1}{2}mv^2 \quad\Longrightarrow\quad \boxed{v = \sqrt{2gh}}
\]
Résultat remarquable : la vitesse d'arrivée ne dépend \emph{pas} de la masse. Vérifions l'homogénéité : $[gh] = \mathrm{m\,s^{-2}}\times\mathrm{m} = \mathrm{m^2\,s^{-2}}$, donc $[\sqrt{2gh}] = \mathrm{m\,s^{-1}}$. Correct.

\begin{popfigure}
\begin{center}
\begin{tikzpicture}
\begin{axis}[
  width=0.85\linewidth, height=6.2cm,
  xlabel={altitude $z$ (m)}, ylabel={énergie (J)},
  xmin=0, xmax=5.2, ymin=0, ymax=11,
  legend style={at={(0.03,0.97)}, anchor=north west, font=\small},
  grid=both, grid style={popline!50},
]
\addplot[popcurve, popblue, domain=0:5, samples=50] {2*x};
\addlegendentry{$E_p(z)=mgz$}
\addplot[popcurve, poporange, domain=0:5, samples=50] {10-2*x};
\addlegendentry{$E_c(z)=E_m-mgz$}
\addplot[popcurve, popgreen, domain=0:5, samples=2] {10};
\addlegendentry{$E_m=10$ J (constante)}
\end{axis}
\end{tikzpicture}
\end{center}
\legende{Chute libre d'une bille de $m = 0{,}204$ kg lâchée de $h = 5$ m ($E_m = mgh = 10$ J). En tombant, l'énergie potentielle se convertit intégralement en énergie cinétique ; leur somme reste constante.}
\end{popfigure}

\textbf{Pendule simple.} Un pendule de longueur $\ell$ écarté d'un angle $\alpha_0$ et lâché sans vitesse : la tension du fil, perpendiculaire au déplacement, ne travaille pas ; seul le poids travaille. L'énergie mécanique se conserve. En prenant $E_p = 0$ au point le plus bas, la hauteur initiale est $h = \ell(1-\cos\alpha_0)$, d'où la vitesse maximale au passage à la verticale :
\[
v_{\max} = \sqrt{2g\ell(1-\cos\alpha_0)}
\]

\begin{popfigure}
\begin{center}
\begin{tikzpicture}[scale=1.05]
\coordinate (O) at (0,0);
\coordinate (A) at (-2.6,-2.5);
\coordinate (B) at (0,-3.6);
\coordinate (C) at (2.6,-2.5);
\draw[popline] (-3.2,0) -- (3.2,0);
\fill[poppurple!40] (-3.2,0) rectangle (3.2,0.15);
\draw[popmuted, dashed] (O) -- (B);
\draw[popmuted, dashed] (O) -- (C);
\draw[thick, popblue] (O) -- (A);
\fill[poporange] (A) circle (0.16);
\fill[popblue] (B) circle (0.16);
\draw[popmuted, dashed] (A) -- ++(2.6,0) coordinate (H);
\draw[<->, popdark] ($(H)+(0.15,0)$) -- ($(B)+(0.15,0)$) node[midway, right, font=\small]{$h=\ell(1-\cos\alpha_0)$};
\draw[popgreen, ->, thick] (A) ++(0,-0.35) -- ++(0,-0.8) node[below, font=\small]{$\vec{P}$};
\draw[popgreen, ->, thick] (B) ++(0.25,0) -- ++(1.1,0) node[right, font=\small]{$\vec{v}_{\max}$};
\draw[popdark] (0,-1.1) arc[start angle=-90, end angle=-136, radius=1.1] node[midway, below left, font=\small]{$\alpha_0$};
\draw[->, poppurple, thick] (A) ++(0.15,-0.55) arc[start angle=160, end angle=245, radius=1.5];
\node[font=\small, popdark] at (-2.6,-3.6) {$E_p$ max, $E_c=0$};
\node[font=\small, popdark] at (1.6,-4.15) {$E_p=0$, $E_c$ max};
\fill (O) circle (0.06) node[above right, font=\small]{$O$};
\end{tikzpicture}
\end{center}
\legende{Pendule simple : aux positions extrêmes, toute l'énergie est potentielle ; au point bas, elle est entièrement cinétique. La tension du fil ne travaille jamais.}
\end{popfigure}

\textbf{Solide sur piste sans frottement et boucle (complément Bac).} Sur une piste parfaitement lisse, la réaction $\vec{R}_N$ ne travaille pas : $E_m$ se conserve. Pour une boucle verticale de rayon $r$, un mobile arrivant avec une vitesse $v_0$ à la base atteint le sommet avec $v_S$ telle que $\dfrac{1}{2}mv_0^2 = \dfrac{1}{2}mv_S^2 + 2mgr$. Pour que le mobile reste en contact au sommet, la deuxième loi de Newton projetée sur la normale impose $v_S^2 \geq gr$ (la réaction doit rester positive ou nulle). On en déduit la \cle{vitesse minimale} à la base : $v_{0,\min} = \sqrt{5gr}$. Ce résultat classique combine méthode énergétique (pour relier les vitesses) et méthode dynamique (pour la condition de contact) : un beau résumé de tout le chapitre.

\begin{exoresolu}[Bille lâchée de 5 m, avec et sans frottements]
Une bille de masse $m = 100$ g est lâchée sans vitesse initiale d'une hauteur $h = 5{,}0$ m. On prend $g = 10\ \mathrm{m\,s^{-2}}$.
\begin{enumerate}
\item En négligeant les frottements, calculer la vitesse d'arrivée au sol.
\item En réalité, l'air exerce une force de frottement moyenne $f = 0{,}5$ N sur tout le trajet. Recalculer la vitesse d'arrivée.
\end{enumerate}
\tcblower
\textbf{1. Sans frottements.} Seul le poids travaille : $E_m$ se conserve. Avec $E_p = 0$ au sol :
\[
mgh = \dfrac{1}{2}mv^2 \;\Longrightarrow\; v = \sqrt{2gh} = \sqrt{2\times 10\times 5{,}0} = \sqrt{100} = 10\ \mathrm{m\,s^{-1}}
\]
\textbf{2. Avec frottements.} $E_m$ ne se conserve plus : $\Delta E_m = W(\vec{f}) = -f\,h$ (force opposée au déplacement). Donc :
\[
\dfrac{1}{2}mv^2 - mgh = -fh
\;\Longrightarrow\;
v = \sqrt{2gh - \dfrac{2fh}{m}} = \sqrt{100 - \dfrac{2\times 0{,}5\times 5{,}0}{0{,}100}} = \sqrt{50} \approx 7{,}1\ \mathrm{m\,s^{-1}}
\]
L'énergie dissipée en chaleur vaut $fh = 2{,}5$ J, soit la moitié de l'énergie mécanique initiale ($mgh = 5{,}0$ J). La bille arrive nettement moins vite : les frottements ne sont pas négligeables.
\end{exoresolu}

\courssub{Non-conservation : frottements et énergie dissipée}

\begin{propriete}[Diminution de l'énergie mécanique en présence de frottements]
Les forces de frottement s'opposent toujours au mouvement : leur travail est résistant, $W(\vec{f}) < 0$. Par conséquent :
\[
\Delta E_m = W(\vec{f}) < 0
\]
L'énergie mécanique \emph{diminue} au cours du mouvement. L'énergie « perdue » n'a pas disparu : elle est convertie en \cle{chaleur} (échauffement de la piste, de l'air, des freins d'une mototaxi). C'est le principe de conservation de l'énergie totale, plus général que la conservation de l'énergie mécanique.
\end{propriete}

\begin{methode}[Calculer une force de frottement moyenne par l'énergie]
Pour déterminer une force de frottement moyenne $f$ supposée constante sur un trajet $AB$ de longueur $L$ :
\begin{enumerate}
\item Faire le bilan des forces et repérer celles qui travaillent.
\item Calculer $E_m(A)$ et $E_m(B)$ à partir des données (vitesses, altitudes).
\item Écrire $\Delta E_m = E_m(B) - E_m(A) = W(\vec{f}) = -fL$.
\item En déduire $f = \dfrac{E_m(A)-E_m(B)}{L}$ et vérifier que $f > 0$.
\end{enumerate}
\end{methode}

\begin{exemplebox}[Freinage d'une mototaxi]
Une mototaxi et son conducteur ($m = 180$ kg) roulent à $v = 54$ km/h $= 15$ m/s sur une route horizontale de Douala. Le conducteur freine et s'immobilise sur $L = 30$ m. Force de freinage moyenne ?
Route horizontale : $E_p$ constante, donc $\Delta E_m = \Delta E_c = 0 - \dfrac{1}{2}mv^2 = -\dfrac{1}{2}\times 180\times 15^2 = -20\,250$ J. D'où :
\[
f = \dfrac{20\,250}{30} = 675\ \mathrm{N}
\]
Cette énergie de 20,25 kJ est dissipée en chaleur dans les freins — d'où l'odeur de plaquettes brûlées dans les descentes de la côte ouest !
\end{exemplebox}

\courssub{Méthode dynamique ou méthode énergétique ?}

Deux grandes stratégies s'offrent à toi pour résoudre un problème de mécanique. Le bon choix fait gagner un temps précieux le jour du Bac.

\begin{methode}[Arbre de choix de la méthode]
\begin{itemize}
\item La question porte sur une \textbf{accélération}, une \textbf{force}, une \textbf{durée}, une \textbf{date} ou une \textbf{équation horaire} $\to$ \textbf{méthode dynamique} : bilan des forces, $\sum\vec{F}=m\vec{a}_G$, projections, intégration.
\item La question porte sur une \textbf{vitesse en un point}, une \textbf{hauteur}, une \textbf{distance}, sans que le temps intervienne $\to$ \textbf{méthode énergétique} : conservation de $E_m$ si aucune force non conservative ne travaille, sinon $\Delta E_m = W(\vec{f})$.
\item En cas de doute : la méthode énergétique est presque toujours plus courte, mais elle ne donne \emph{jamais} ni le temps, ni les forces.
\end{itemize}
\end{methode}

\begin{popfigure}
\begin{center}
\begin{tikzpicture}[
  font=\small,
  decision/.style={diamond, draw=poppurple, fill=poppurpleL, aspect=2.2, align=center, inner sep=1pt, text width=3.6cm},
  bloc/.style={rectangle, rounded corners, draw=popblue, fill=popblueL, align=center, text width=4.6cm, minimum height=1.1cm},
  res/.style={rectangle, rounded corners, draw=popgreen, fill=popgreenL, align=center, text width=4.2cm},
  fleche/.style={->, thick, popdark}
]
\node[decision] (q) at (0,0) {La question fait-elle intervenir le \textbf{temps} ou une \textbf{force} ?};
\node[bloc, align=center] (dyn) at (-4.6,-2.6) {\textbf{Méthode dynamique}\\ bilan des forces, $\sum\vec{F}=m\vec{a}_G$, projections, intégration};
\node[decision] (q2) at (4.6,-2.6) {Des frottements (ou forces non conservatives) travaillent-ils ?};
\node[res] (cons) at (2.2,-5.4) {\textbf{Conservation} : $E_m(A)=E_m(B)$};
\node[res] (noncons) at (7.4,-5.4) {\textbf{Théorème de $E_m$} : $\Delta E_m = W(\vec{f})$};
\draw[fleche] (q) -- node[above left]{oui} (dyn);
\draw[fleche] (q) -- node[above right]{non} (q2);
\draw[fleche] (q2) -- node[above left]{non} (cons);
\draw[fleche] (q2) -- node[above right]{oui} (noncons);
\end{tikzpicture}
\end{center}
\legende{Arbre de décision : le temps et les forces relèvent de Newton ; les vitesses et positions relèvent de l'énergie.}
\end{popfigure}

\begin{aretenir}[L'essentiel de la section]
\begin{itemize}
\item Une force conservative a un travail indépendant du chemin ; on lui associe une énergie potentielle : $E_p = mgz$ (pesanteur), $E_{p,\text{él}} = \tfrac{1}{2}kx^2$ (ressort, complément).
\item Énergie mécanique : $E_m = E_c + E_p$. Théorème : $\Delta E_m = W(\text{forces non conservatives})$.
\item Si seules des forces conservatives (ou ne travaillant pas) agissent : $E_m$ = constante. Sinon, avec frottements : $\Delta E_m = W(\vec{f}) < 0$, l'énergie dissipée devient chaleur.
\item Vitesse de chute libre : $v = \sqrt{2gh}$ ; pendule : $v_{\max} = \sqrt{2g\ell(1-\cos\alpha_0)}$.
\item Question de temps ou de force $\to$ Newton ; question de vitesse ou de position $\to$ énergie.
\end{itemize}
\end{aretenir}

\begin{exercice}[Pour s'entraîner]
Un enfant de masse $m = 30$ kg glisse du haut d'un toboggan de hauteur $h = 3{,}0$ m, sans vitesse initiale. Il arrive en bas avec une vitesse $v = 6{,}0$ m/s. On prend $g = 10\ \mathrm{m\,s^{-2}}$.
\begin{enumerate}
\item Quelle serait sa vitesse d'arrivée en l'absence de frottements ?
\item Calculer la variation d'énergie mécanique et en déduire le travail des frottements.
\item La glissière mesure $L = 5{,}0$ m ; en déduire la force de frottement moyenne.
\end{enumerate}
\end{exercice}

\begin{corrige}[Exercice : toboggan]
\textbf{1.} Sans frottements : $v_0 = \sqrt{2gh} = \sqrt{2\times 10\times 3{,}0} = \sqrt{60} \approx 7{,}7\ \mathrm{m\,s^{-1}}$.

\textbf{2.} $\Delta E_m = \dfrac{1}{2}mv^2 - mgh = \dfrac{1}{2}\times 30\times 6{,}0^2 - 30\times 10\times 3{,}0 = 540 - 900 = -360$ J. Donc $W(\vec{f}) = -360$ J : l'énergie mécanique diminue, cohérent avec des frottements.

\textbf{3.} $W(\vec{f}) = -fL$ donc $f = \dfrac{360}{5{,}0} = 72\ \mathrm{N}$. La vitesse réelle ($6{,}0$ m/s) est bien inférieure à la vitesse idéale ($7{,}7$ m/s).
\end{corrige}

\begin{savaistu}[Le barrage de Nachtigal : l'énergie potentielle à l'échelle d'un pays]
Le barrage hydroélectrique de Nachtigal, sur le fleuve Sanaga, stocke l'eau à une trentaine de mètres au-dessus des turbines. Chaque kilogramme d'eau possède une énergie potentielle $E_p = mgh$ qui se convertit en énergie cinétique dans les conduites forcées, puis en énergie électrique dans les alternateurs. Avec une puissance de 420 MW, c'est littéralement la conservation — et les conversions — de l'énergie mécanique qui éclaire le Cameroun. Les pertes par frottements dans les conduites, elles, sont bien réelles : les ingénieurs appliquent exactement le bilan $\Delta E_m = W(\vec{f})$ que tu viens d'apprendre !
\end{savaistu}

\newpage

\coursec{Activité d'intégration}

\coursec{Lois de Newton et conservation de l'énergie mécanique}

\courssub{Objectifs de l'activité}

À l'issue de cette activité d'intégration, tu seras capable de :
\begin{itemize}
\item mobiliser, sur une \textbf{même situation réelle}, les trois outils de la leçon : le théorème du centre d'inertie (TCI), le théorème de l'énergie cinétique (TEC) et la conservation de l'énergie mécanique ;
\item choisir, en le justifiant, la méthode la plus adaptée (dynamique ou énergétique) selon la question posée ;
\item exploiter un écart entre une valeur théorique et une valeur mesurée pour remonter à une force de frottement ;
\item rédiger une conclusion argumentée comparant les démarches, compétence essentielle de l'épreuve de physique au Baccalauréat.
\end{itemize}

\courssub{Situation}

Pour la fête de la Jeunesse, la commune de Mbalmayo a inauguré une aire de jeux équipée d'un grand toboggan. Le maire, inquiet après qu'un enfant est arrivé « trop vite » en bas, demande au club scientifique du lycée de vérifier la sécurité de l'installation et d'expliquer \emph{pourquoi} la vitesse réelle est plus faible que celle annoncée par le fabricant.

Le toboggan est modélisé par un plan incliné AB de longueur $L = \SI{8,0}{m}$, dont le sommet A est à la hauteur $h = \SI{4,0}{m}$ au-dessus du sol. En B, la glissière rejoint une piste horizontale BC en béton. Un enfant de masse $m = \SI{30}{kg}$ part de A \textbf{sans vitesse initiale} et glisse jusqu'en B. On prendra $g = \SI{10}{N.kg^{-1}}$.

Le fabricant annonce une vitesse maximale en B calculée « dans le cas idéal, sans frottement ». Or, une mesure au radar du club donne $v_B = \SI{7,0}{m.s^{-1}}$. Sur la piste horizontale, les frottements sont modélisés par une force constante $\vec{f}'$ d'intensité $f' = \SI{60}{N}$.

\begin{popfigure}
\begin{tikzpicture}[scale=0.9]
\coordinate (A) at (0,3);
\coordinate (B) at (6,0);
\coordinate (C) at (10.5,0);
\draw[very thick,popblue] (A) -- (B);
\draw[very thick,popblue] (B) -- (C);
\draw[popmuted,dashed] (A) -- (0,0);
\draw[popmuted] (0,0) -- (B);
\draw[<->,popdark] (-0.25,0) -- (-0.25,3) node[midway,left]{$h$};
\fill[poporange] (A) circle (0.09) node[above left,popdark]{A};
\fill[poporange] (B) circle (0.09) node[below right,popdark]{B};
\fill[poporange] (C) circle (0.09) node[below,popdark]{C};
\draw[<->,popdark] (0,-0.45) -- (6,-0.45) node[midway,below]{$L = \SI{8,0}{m}$ (pente)};
\node[popdark] at (8.3,0.35) {piste horizontale rugueuse};
\draw[->,popgreen,thick] (5.2,0.25) -- (6.4,0.25) node[midway,above]{$\vec{v}_B$};
\end{tikzpicture}
\legende{Modélisation du toboggan de Mbalmayo : plan incliné AB puis piste horizontale BC.}
\end{popfigure}

\courssub{Consignes}

Travail en groupes de quatre. Chaque groupe rédige un rapport et remplit le tableau comparatif final.

\begin{enumerate}
\item \textbf{Cas idéal.} En supposant le toboggan parfaitement lisse (frottements négligeables), calculer la vitesse théorique $v_B^{\text{th}}$ de l'enfant en B en utilisant la \textbf{conservation de l'énergie mécanique}.
\item \textbf{Vérification par la dynamique.} Faire le bilan des forces sur le toboggan lisse, appliquer le \textbf{TCI}, établir l'accélération $a$ de l'enfant, puis retrouver $v_B^{\text{th}}$ à l'aide de la relation de cinématique adaptée. Comparer avec la question 1.
\item \textbf{Cas réel.} La mesure donne $v_B = \SI{7,0}{m.s^{-1}}$. En appliquant le \textbf{théorème de l'énergie cinétique} entre A et B, déterminer le travail de la force de frottement moyenne $\vec{f}$ sur le toboggan, puis son intensité $f$.
\item \textbf{Arrêt sur la piste horizontale.} Déterminer, par la méthode de ton choix (à justifier), la distance d'arrêt $d = $ BC sur la piste horizontale où $f' = \SI{60}{N}$. Conclure sur la sécurité si l'aire de jeux dispose de $\SI{3,0}{m}$ de dégagement.
\item \textbf{Synthèse.} Rédiger une conclusion argumentée : pour chaque étape, quelle méthode (TCI, TEC, conservation de $E_m$) était la plus rapide, et pourquoi ? Compléter le tableau comparatif ci-dessous.
\end{enumerate}

\courssub{Ressources à mobiliser}

\begin{aretenir}[Boîte à outils de la leçon]
\begin{itemize}
\item \textbf{TCI} (référentiel terrestre galiléen) : $\sum \vec{F}_{\text{ext}} = m\,\vec{a}_G$.
\item \textbf{TEC} entre A et B : $E_{c,B} - E_{c,A} = \sum W_{A\to B}(\vec{F}_{\text{ext}})$.
\item \textbf{Énergie mécanique} : $E_m = E_c + E_p$ avec $E_p = mgh$ (référence au sol). $E_m$ se conserve si seules des forces conservatives travaillent ; sinon $\Delta E_m = W(\vec{f})$.
\item \textbf{Cinématique} (mouvement rectiligne uniformément varié) : $v_B^2 - v_A^2 = 2a\,L$.
\item Travail d'une force constante : $W = \vec{F}\cdot\overrightarrow{AB} = F\,L\cos\alpha$ ; travail du poids : $W(\vec{P}) = mgh$ (descente) ; travail de la réaction normale : nul car $\vec{R}_N \perp$ déplacement.
\end{itemize}
\end{aretenir}

\courssub{Démarche et corrigé commenté}

\begin{exoresolu}[Le toboggan du parc : dynamique ou énergie ?]
\textbf{Énoncé.} Un enfant ($m = \SI{30}{kg}$) glisse sans vitesse initiale depuis A ($h = \SI{4,0}{m}$) le long d'un toboggan de longueur $L = \SI{8,0}{m}$. (1) Calculer $v_B^{\text{th}}$ sans frottement par l'énergie. (2) Retrouver ce résultat par le TCI. (3) Sachant que $v_B = \SI{7,0}{m.s^{-1}}$ en réalité, trouver la force de frottement moyenne $f$ sur le toboggan. (4) Calculer la distance d'arrêt sur la piste horizontale ($f' = \SI{60}{N}$). (5) Comparer les méthodes.

\tcblower

\textbf{Question 1 — Conservation de l'énergie mécanique.}

Sans frottement, seules le poids (conservatif) et la réaction normale (travail nul) interviennent : $E_m$ se conserve entre A et B. En prenant le sol comme référence ($E_{p,B} = 0$) et $v_A = 0$ :
\[
E_{m,A} = E_{m,B} \;\Longrightarrow\; mgh = \tfrac{1}{2}m\left(v_B^{\text{th}}\right)^2
\;\Longrightarrow\; v_B^{\text{th}} = \sqrt{2gh} = \sqrt{2\times 10\times 4{,}0} = \sqrt{80} \approx \SI{8,9}{m.s^{-1}}.
\]
Remarque : $m$ se simplifie — la vitesse ne dépend pas de la masse, résultat typique de la méthode énergétique.

\textbf{Question 2 — Vérification par le TCI.}

Bilan des forces sur le toboggan lisse : le poids $\vec{P}$ et la réaction normale $\vec{R}_N$. L'angle $\alpha$ du plan vérifie $\sin\alpha = \dfrac{h}{L} = \dfrac{4{,}0}{8{,}0} = 0{,}5$, soit $\alpha = 30^{\circ}$. En projetant le TCI $\vec{P} + \vec{R}_N = m\vec{a}$ sur l'axe du mouvement (descendant) :
\[
mg\sin\alpha = ma \;\Longrightarrow\; a = g\sin\alpha = 10\times 0{,}5 = \SI{5,0}{m.s^{-2}}.
\]
Le mouvement est rectiligne uniformément accéléré ; avec $v_A = 0$ :
\[
v_B^2 = 2aL = 2\times 5{,}0\times 8{,}0 = 80 \;\Longrightarrow\; v_B^{\text{th}} = \sqrt{80} \approx \SI{8,9}{m.s^{-1}}.
\]
Les deux méthodes donnent le même résultat : la physique est cohérente ! Mais la méthode énergétique a évité le calcul de l'angle et de l'accélération.

\textbf{Question 3 — Force de frottement réelle par le TEC.}

Entre A et B, le TEC s'écrit :
\[
E_{c,B} - E_{c,A} = W(\vec{P}) + W(\vec{R}_N) + W(\vec{f})
\;\Longrightarrow\; \tfrac{1}{2}mv_B^2 - 0 = mgh + 0 - fL.
\]
Application numérique :
\begin{align*}
fL &= mgh - \tfrac{1}{2}mv_B^2 = 30\times 10\times 4{,}0 - \tfrac{1}{2}\times 30\times 7{,}0^2 \\
&= 1200 - 735 = \SI{465}{J}
\;\Longrightarrow\; f = \dfrac{465}{8{,}0} \approx \SI{58}{N}.
\end{align*}
L'écart de vitesse ($\SI{8,9}{m.s^{-1}}$ théorique contre $\SI{7,0}{m.s^{-1}}$ mesurée) s'explique par une dissipation de $\SI{465}{J}$ sous forme de chaleur : c'est la \textbf{non-conservation de l'énergie mécanique}, $\Delta E_m = W(\vec{f}) = \SI{-465}{J}$.

\textbf{Question 4 — Distance d'arrêt sur la piste horizontale.}

Sur BC, poids et réaction sont perpendiculaires au déplacement (travail nul) ; seule $\vec{f}'$ travaille. Le TEC entre B et C (arrêt, $v_C = 0$) :
\[
0 - \tfrac{1}{2}mv_B^2 = -f'\,d
\;\Longrightarrow\; d = \dfrac{mv_B^2}{2f'} = \dfrac{30\times 7{,}0^2}{2\times 60} = \dfrac{1470}{120} = \SI{12,25}{m} \approx \SI{12}{m}.
\]
Le dégagement de $\SI{3,0}{m}$ est largement insuffisant : le club recommande d'allonger la zone de réception ou d'augmenter le frottement (tapis de réception).

\textbf{Question 5 — Tableau comparatif.}

\begin{center}
\begin{tabularx}{\linewidth}{|l|X|X|}
\hline
\rowcolor{popblueL} \textbf{Étape} & \textbf{Méthode la plus rapide} & \textbf{Pourquoi ?} \\
\hline
1. Vitesse en B (cas idéal) & Conservation de $E_m$ & Une seule ligne ; pas besoin de l'angle ni de l'accélération. \\
\hline
2. Accélération sur la pente & TCI & L'énergie ne donne pas directement $\vec{a}$ ni les forces. \\
\hline
3. Force de frottement réelle & TEC & Relie directement la variation de $E_c$ au travail de $\vec{f}$. \\
\hline
4. Distance d'arrêt & TEC & Évite de repasser par l'accélération et le temps. \\
\hline
\end{tabularx}
\end{center}
\end{exoresolu}

\courssub{Bilan de l'activité}

Cette activité a permis de mettre en évidence plusieurs idées centrales de la leçon :
\begin{itemize}
\item \textbf{cohérence des lois} : le TCI et la conservation de l'énergie mécanique conduisent au même résultat, car le TEC découle du TCI — ce sont deux regards sur la même réalité ;
\item \textbf{choix stratégique de méthode} : l'énergie est imbattable pour relier vitesses et positions sans connaître le détail du mouvement ; la dynamique reste indispensable dès qu'on cherche une accélération, un temps de parcours ou une force de contact ;
\item \textbf{rôle des frottements} : l'écart entre le modèle idéal ($\SI{8,9}{m.s^{-1}}$) et la mesure ($\SI{7,0}{m.s^{-1}}$) traduit la non-conservation de l'énergie mécanique, $\Delta E_m = W(\vec{f}) < 0$, et permet de quantifier une force difficile à mesurer directement ;
\item \textbf{compétence citoyenne} : le raisonnement physique aboutit à une recommandation concrète de sécurité (allonger la zone de réception), comme le ferait un ingénieur du génie civil ou un contrôleur d'installations.
\end{itemize}

\begin{attention}[Erreurs fréquentes repérées pendant l'activité]
\begin{itemize}
\item Écrire $W(\vec{P}) = -mgh$ alors que l'enfant \emph{descend} : le poids est moteur, $W(\vec{P}) = +mgh$.
\item Faire travailler la réaction normale : $\vec{R}_N \perp$ déplacement, donc $W(\vec{R}_N) = 0$.
\item Oublier que la conservation de $E_m$ exige l'absence de frottements : avec $v_B = \SI{7,0}{m.s^{-1}}$, écrire $mgh = \tfrac{1}{2}mv_B^2$ est faux.
\item Confondre $L$ (distance sur la pente) et $h$ (hauteur verticale) dans $W(\vec{f}) = -fL$.
\item Négliger les chiffres significatifs : les données ($h=4,0$, $L=8,0$, $v_B=7,0$) imposent deux chiffres significatifs sur les résultats ($v_B^{\text{th}}=8,9$, $f=58$, $d=12$).
\end{itemize}
\end{attention}

\newpage

\coursec{Méthodes et exercices résolus}

\coursec{Leçon P05 : Lois de Newton et conservation de l'énergie mécanique}

\courssub{Rappels de cinématique et repère de Frenet}

\begin{definition}[Repères et vecteurs cinématiques]
Dans un référentiel $\mathcal{R}$ muni d'un repère cartésien $(O, \vec{\imath}, \vec{\jmath}, \vec{k})$, le mouvement d'un point matériel $M$ est décrit par :
\begin{itemize}
\item le \textbf{vecteur position} $\overrightarrow{OM}(t) = x(t)\vec{\imath} + y(t)\vec{\jmath} + z(t)\vec{k}$ ;
\item le \textbf{vecteur vitesse} $\vec{v}(t) = \dfrac{d\overrightarrow{OM}}{dt} = \dot{x}\vec{\imath} + \dot{y}\vec{\jmath} + \dot{z}\vec{k}$ ; sa norme $v = \|\vec{v}\|$ est la vitesse (scalaire) ;
\item le \textbf{vecteur accélération} $\vec{a}(t) = \dfrac{d\vec{v}}{dt} = \ddot{x}\vec{\imath} + \ddot{y}\vec{\jmath} + \ddot{z}\vec{k}$.
\end{itemize}
\end{definition}

\begin{propriete}[Repère de Frenet -- Mouvement planaire]
Pour une trajectoire plane, en chaque point $M$ on définit le trièdre mobile $(M, \vec{T}, \vec{N}, \vec{B})$ :
\begin{itemize}
\item $\vec{T}$ : vecteur unitaire \textbf{tangent} à la trajectoire, orienté dans le sens du mouvement ($\vec{v} = v\vec{T}$) ;
\item $\vec{N}$ : vecteur unitaire \textbf{normal principal}, dirigé vers le centre de courbure ($\vec{N} \perp \vec{T}$) ;
\item $\vec{B} = \vec{T} \wedge \vec{N}$ : binormale (constante pour un mouvement planaire).
\end{itemize}
L'accélération se décompose en :
\[ \vec{a} = \frac{dv}{dt}\,\vec{T} + \frac{v^2}{\mathcal{R}}\,\vec{N} \]
où $\mathcal{R}$ est le rayon de courbure. Pour un \textbf{cercle de rayon $R$} : $\mathcal{R} = R$ et $\vec{a} = \dot{v}\,\vec{T} + \dfrac{v^2}{R}\,\vec{N}$.
\end{propriete}

\begin{aretenir}[Composantes de l'accélération]
\begin{itemize}
\item \textbf{Accélération tangentielle} $a_T = \dfrac{dv}{dt}$ : modifie la \emph{norme} de la vitesse (ralentissement ou accélération).
\item \textbf{Accélération normale (centripète)} $a_N = \dfrac{v^2}{R}$ : modifie la \emph{direction} de la vitesse, toujours dirigée vers le centre de courbure.
\end{itemize}
\textbf{Analyse dimensionnelle} : $[v^2/R] = (\text{m}^2\!\cdot\!\text{s}^{-2})/\text{m} = \text{m}\!\cdot\!\text{s}^{-2} = [a]$. Cohérent.
\end{aretenir}

\begin{popfigure}
\begin{tikzpicture}[scale=1.2, >=Stealth]
\draw[thick, popblue] (0,0) circle (2.5);
\coordinate (O) at (0,0);
\coordinate (M) at (45:2.5);
\draw[->, thick, popdark] (O) -- (M) node[midway, above left] {$R$};
\draw[->, thick, poporange] (M) -- ++(135:1.2) node[above left] {$\vec{N}$};
\draw[->, thick, popgreen] (M) -- ++(-45:1.2) node[below right] {$\vec{T}$};
\draw[->, thick, poppurple] (M) -- ++(45:1.0) node[above right] {$\vec{v}$};
\draw[->, thick, poppink] (M) -- ++(135:0.7) node[above left] {$\vec{a}_N$};
\draw[->, thick, poppink] (M) -- ++(-45:0.5) node[below right] {$\vec{a}_T$};
\fill (M) circle (2pt) node[below right] {$M$};
\fill (O) circle (2pt) node[below left] {$O$};
\end{tikzpicture}
\legende{Repère de Frenet sur une trajectoire circulaire : $\vec{T}$ tangent, $\vec{N}$ normal entrant, $\vec{a}_N$ centripète.}
\end{popfigure}

\courssub{Référentiels galiléens et première loi de Newton}

\begin{definition}[Référentiel galiléen]
Un référentiel $\mathcal{R}$ est dit \textbf{galiléen} (ou inertiel) si tout point matériel isolé (soumis à aucune force, ou à des forces dont la résultante est nulle) y a un mouvement rectiligne uniforme (ou est au repos).
\end{definition}

\begin{propriete}[Première loi de Newton -- Principe d'inertie]
Dans un référentiel galiléen, un système matériel reste au repos ou en mouvement rectiligne uniforme tant que la résultante des forces extérieures qui s'exercent sur lui est nulle.
\[ \sum \vec{F}_{\text{ext}} = \vec{0} \quad \Longleftrightarrow \quad \vec{v}_G = \text{cste} \quad (\text{ou } \vec{a}_G = \vec{0}) \]
\end{propriete}

\begin{attention}[Référentiel terrestre]
Le référentiel lié à la Terre (héliocentrique) n'est \emph{pas} rigoureusement galiléen (rotation terrestre, translation orbitale). Pour la mécanique au lycée (durées courtes, vitesses $\ll c$), on le \textbf{suppose galiléen}. Cette approximation est excellente pour les mouvements quotidiens (voitures, projectiles, machines).
\end{attention}

\courssub{Deuxième et troisième lois de Newton}

\begin{propriete}[Deuxième loi -- Théorème du centre d'inertie]
Dans un référentiel galiléen, la dérivée temporelle du vecteur quantité de mouvement $\vec{P} = m\vec{v}_G$ d'un système de masse $m$ est égale à la résultante des forces extérieures :
\[ \sum \vec{F}_{\text{ext}} = \frac{d\vec{P}}{dt} = m\,\vec{a}_G \qquad (\text{masse constante}) \]
\end{propriete}

\begin{propriete}[Troisième loi -- Principe des actions réciproques]
Si un système $A$ exerce une force $\vec{F}_{A \to B}$ sur un système $B$, alors $B$ exerce sur $A$ une force $\vec{F}_{B \to A}$ telle que :
\[ \vec{F}_{A \to B} = - \vec{F}_{B \to A} \]
Ces deux forces ont même droite d'action, même norme, sens opposés. Elles s'appliquent sur des \textbf{systèmes différents} : elles ne s'annulent jamais dans un même bilan de forces.
\end{propriete}

\begin{methode}[Appliquer la deuxième loi de Newton à un point matériel]
\begin{enumerate}
\item \textbf{Définir le système} étudié et le \textbf{référentiel} d'étude (vérifier qu'il est galiléen, ou supposé tel).
\item \textbf{Faire le bilan des forces extérieures} appliquées au système, puis les représenter sur un schéma soigné (origine au centre d'inertie $G$).
\item \textbf{Choisir un repère} adapté : axes $(Ox)$, $(Oy)$ orientés si possible l'un selon le mouvement, l'autre perpendiculaire ; ou repère de Frenet pour un mouvement circulaire.
\item \textbf{Écrire le théorème du centre d'inertie} vectoriellement : $\sum \vec{F}_{\text{ext}} = m\,\vec{a}_G$.
\item \textbf{Projeter} sur chaque axe pour obtenir des équations scalaires.
\item \textbf{Résoudre} : intégrer deux fois pour obtenir $\vec{v}(t)$ puis $\overrightarrow{OG}(t)$, en utilisant les conditions initiales ($\vec{v}_0$, position initiale).
\item \textbf{Exploiter} les équations horaires : nature de la trajectoire, date ou position cherchée, vitesse à un instant donné.
\item \textbf{Vérifier} : homogénéité des résultats, cohérence des signes, ordre de grandeur.
\end{enumerate}
\textbf{Réflexes} : une force constante $\Rightarrow$ mouvement à accélération constante ; si $\sum \vec{F} = \vec{0}$, le mouvement est rectiligne uniforme (ou le repos) ; sur un support sans frottement, la réaction est normale au support et ne travaille pas.
\end{methode}

\begin{exoresolu}[Descente d'une mototaxi sur une pente]
Une mototaxi de masse totale $m = \SI{180}{kg}$ (conducteur compris) descend une route rectiligne inclinée d'un angle $\alpha = 10^{\circ}$ sur l'horizontale, frein moteur coupé. Les frottements sont modélisés par une force unique $\vec{f}$ de norme constante $f = \SI{120}{N}$, opposée au mouvement. La moto part du repos en haut de la pente. On prend $g = \SI{9,8}{m.s^{-2}}$.
\begin{enumerate}
\item Faire le bilan des forces et les représenter.
\item Déterminer l'accélération de la moto.
\item Quelle distance parcourt-elle pour atteindre la vitesse $v = \SI{36}{km.h^{-1}}$ ?
\end{enumerate}
\tcblower
\textbf{1. Bilan des forces.} Système : la mototaxi, assimilée à son centre d'inertie $G$. Référentiel terrestre, supposé galiléen. Trois forces extérieures :
\begin{itemize}
\item le poids $\vec{P} = m\vec{g}$, vertical, vers le bas ;
\item la réaction normale $\vec{R}_N$ de la route, perpendiculaire à la pente ;
\item la force de frottement $\vec{f}$, parallèle à la pente, dirigée vers le haut (opposée au mouvement).
\end{itemize}
On choisit le repère $(G, \vec{\imath}, \vec{\jmath})$ avec $(Gx)$ dirigé vers le bas de la pente et $(Gy)$ perpendiculaire à la pente vers le haut.

\begin{popfigure}
\begin{tikzpicture}[scale=1.2, >=Stealth]
\draw[thick, popdark] (-3,-1) -- (3,1) node[right] {Pente $\alpha$};
\draw[dashed] (-3,0) -- (3,0) node[right] {Horizontal};
\coordinate (G) at (0,0);
\draw[->, thick, poporange] (G) -- ++(-90:1.2) node[below] {$\vec{P}=m\vec{g}$};
\draw[->, thick, popblue] (G) -- ++(90-10:1.2) node[above left] {$\vec{R}_N$};
\draw[->, thick, popgreen] (G) -- ++(-10:1.2) node[below right] {$\vec{f}$};
\draw[->, thick, poppurple] (G) -- ++(170:1.0) node[left] {$\vec{\jmath}$};
\draw[->, thick, poppurple] (G) -- ++(80:1.0) node[above] {$\vec{\imath}$};
\fill (G) circle (2pt) node[below left] {$G$};
\draw (0.5,0) arc (0:10:0.5) node[midway, right, xshift=2mm] {$\alpha$};
\end{tikzpicture}
\legende{Bilan des forces sur la mototaxi. Repère $(Gx)$ vers le bas de la pente.}
\end{popfigure}

\textbf{2. Accélération.} Le théorème du centre d'inertie s'écrit :
\[ \vec{P} + \vec{R}_N + \vec{f} = m\,\vec{a}_G. \]
Le mouvement a lieu selon $(Gx)$, donc $a_y = 0$. Projections :
\begin{align*}
\text{sur } (Gx) : \quad & mg\sin\alpha - f = m\,a_G \\
\text{sur } (Gy) : \quad & R_N - mg\cos\alpha = 0
\end{align*}
D'où :
\[ a_G = g\sin\alpha - \frac{f}{m} = 9{,}8 \times \sin 10^{\circ} - \frac{120}{180}. \]
Avec $\sin 10^{\circ} \approx 0{,}1736$ :
\[ a_G \approx 9{,}8 \times 0{,}1736 - 0{,}6667 \approx 1{,}701 - 0{,}667 = \SI{1,03}{m.s^{-2}}. \]
L'accélération est constante : le mouvement est rectiligne uniformément accéléré.

\textbf{3. Distance parcourue.} Conversion : $v = \SI{36}{km.h^{-1}} = \SI{10}{m.s^{-1}}$. Pour un mouvement uniformément accéléré partant du repos ($v_0=0$) :
\[ v^2 = 2\,a_G\,d \quad \Longrightarrow \quad d = \frac{v^2}{2a_G} = \frac{10^2}{2 \times 1{,}03} \approx \SI{48,5}{m}. \]
\textbf{Conclusion} : la mototaxi atteint $\SI{36}{km.h^{-1}}$ après environ $\SI{49}{m}$ de descente, avec une accélération constante $a_G \approx \SI{1,0}{m.s^{-2}}$.
\end{exoresolu}

\courssub{Méthode 2 : Mettre en équation un mouvement circulaire avec le repère de Frenet}

\begin{methode}[Exploiter le repère de Frenet pour un mouvement circulaire]
\begin{enumerate}
\item Vérifier que la trajectoire est un \textbf{cercle de rayon $R$} et repérer le sens du mouvement.
\item Définir le repère de Frenet $(M, \vec{T}, \vec{N})$ : $\vec{T}$ tangent à la trajectoire dans le sens du mouvement, $\vec{N}$ normal dirigé vers le centre.
\item Écrire l'accélération : $\vec{a} = \dfrac{dv}{dt}\,\vec{T} + \dfrac{v^2}{R}\,\vec{N}$.
\item Appliquer $\sum \vec{F}_{\text{ext}} = m\vec{a}$ et \textbf{projeter sur $\vec{T}$ et $\vec{N}$} :
\[ \sum F_T = m\,\frac{dv}{dt} \qquad ; \qquad \sum F_N = m\,\frac{v^2}{R}. \]
\item Si le mouvement est \textbf{uniforme} ($v$ constante), alors $\dfrac{dv}{dt} = 0$ : la somme des composantes tangentielles est nulle et $\sum F_N = \dfrac{mv^2}{R}$.
\item Résoudre l'équation obtenue (vitesse limite, tension, réaction...).
\end{enumerate}
\textbf{Piège classique} : ne jamais ajouter une « force centrifuge » dans le bilan des forces dans un référentiel galiléen. L'accélération est centripète, fournie par les forces réelles (tension, réaction normale, gravité, frottement).
\end{methode}

\begin{exoresolu}[Vitesse caractéristique d'un virage relevé]
Une voiture de masse $m = \SI{900}{kg}$ décrit un virage circulaire de rayon $R = \SI{50}{m}$ relevé d'un angle $\alpha = 8^{\circ}$ par rapport à l'horizontale. On néglige les frottements. On prend $g = \SI{9,8}{m.s^{-2}}$.
\begin{enumerate}
\item Faire le bilan des forces.
\item Déterminer la vitesse $v_0$ pour laquelle la voiture reste dans le plan du virage sans tendance à déraper.
\end{enumerate}
\tcblower
\textbf{1. Bilan des forces.} Système : la voiture, référentiel terrestre galiléen. Forces : le poids $\vec{P} = m\vec{g}$ (vertical, vers le bas) et la réaction normale $\vec{R}_N$ de la chaussée (perpendiculaire à la piste, inclinée de $\alpha$ par rapport à la verticale).

\begin{popfigure}
\begin{tikzpicture}[scale=1.2, >=Stealth]
\draw[thick, popdark] (-2,-0.5) -- (2,0.5) node[right] {Chaussée};
\coordinate (G) at (0,0);
\draw[->, thick, poporange] (G) -- ++(-90:1.2) node[below] {$\vec{P}=mg$};
\draw[->, thick, popblue] (G) -- ++(90-8:1.2) node[above left] {$\vec{R}_N$};
\draw[dashed] (G) -- ++(-90:0.8);
\draw[dashed] (G) -- ++(90-8:0.8);
\draw (0,-0.3) arc (-90:-90+8:0.3) node[midway, right, xshift=2mm] {$\alpha$};
\draw[->, thick, poppurple] (G) -- ++(-180:1.0) node[left] {$\vec{N}$ (centre)};
\fill (G) circle (2pt) node[below right] {$G$};
\end{tikzpicture}
\legende{Virage relevé : la réaction normale $\vec{R}_N$ a une composante horizontale centripète.}
\end{popfigure}

\textbf{2. Vitesse $v_0$.} Le mouvement est circulaire uniforme dans un plan horizontal ; l'accélération est donc horizontale, centripète : $\vec{a} = \dfrac{v_0^2}{R}\,\vec{N}$ avec $\vec{N}$ horizontal dirigé vers le centre. Le théorème du centre d'inertie donne $\vec{P} + \vec{R}_N = m\vec{a}$. Projections sur axes vertical et horizontal (vers le centre) :
\begin{align*}
\text{Verticale} : \quad & R_N\cos\alpha - mg = 0 \\
\text{Horizontale (vers le centre)} : \quad & R_N\sin\alpha = m\,\frac{v_0^2}{R}
\end{align*}
En divisant membre à membre (élimination de $R_N$ et $m$) :
\[ \tan\alpha = \frac{v_0^2}{Rg} \quad \Longrightarrow \quad v_0 = \sqrt{Rg\tan\alpha}. \]
Numériquement : $\tan 8^{\circ} \approx 0{,}1405$.
\[ v_0 = \sqrt{50 \times 9{,}8 \times 0{,}1405} = \sqrt{68{,}845} \approx \SI{8,30}{m.s^{-1}}. \]
Soit $v_0 \approx 8{,}30 \times 3{,}6 \approx \SI{30}{km.h^{-1}}$.
\textbf{Conclusion} : la vitesse caractéristique (unique sans frottement) est $v_0 \approx \SI{8,3}{m.s^{-1}}$. Remarquons que $v_0$ ne dépend pas de la masse : c'est une propriété géométrique du virage ($R$, $\alpha$).
\end{exoresolu}

\courssub{Dynamique du solide en rotation autour d'un axe fixe}

\begin{definition}[Moment d'inertie]
Pour un solide $S$ en rotation autour d'un axe fixe $(\Delta)$, le \textbf{moment d'inertie} $J_{\Delta}$ est la somme des produits $m_i d_i^2$ pour tous les éléments de masse $m_i$ à distance $d_i$ de l'axe :
\[ J_{\Delta} = \sum_i m_i d_i^2 = \int_S r^2 dm. \]
Unité SI : $\text{kg}\!\cdot\!\text{m}^2$. Il caractérise l'inertie du solide en rotation (analogue de la masse $m$ en translation).
\end{definition}

\begin{aretenir}[Moments d'inertie usuels (axe de symétrie)]
\begin{center}
\begin{tabularx}{\linewidth}{|l|X|c|}\hline
\textbf{Solide} & \textbf{Description} & \textbf{$J_{\Delta}$} \\ \hline
Cylindre / Disque plein & Axe de symétrie & $\frac{1}{2}MR^2$ \\ \hline
Cerceau / Anneau mince & Axe de symétrie & $MR^2$ \\ \hline
Sphère pleine & Diamètre & $\frac{2}{5}MR^2$ \\ \hline
Tige homogène & Axe $\perp$ par un bout & $\frac{1}{3}ML^2$ \\ \hline
Tige homogène & Axe $\perp$ par le centre & $\frac{1}{12}ML^2$ \\ \hline
\end{tabularx}
\end{center}
\textbf{Théorème de Huygens} : $J_{\Delta} = J_G + M d^2$ où $d$ est la distance entre l'axe $(\Delta)$ et l'axe parallèle passant par $G$.
\end{aretenir}

\begin{propriete}[Relation fondamentale de la dynamique de rotation]
Dans un référentiel galiléen, pour un solide en rotation autour d'un axe fixe $(\Delta)$, la somme des moments des forces extérieures par rapport à $(\Delta)$ est égale au moment d'inertie fois l'accélération angulaire :
\[ \sum \mathcal{M}_{\Delta}(\vec{F}_{\text{ext}}) = J_{\Delta}\,\ddot{\theta} \]
où $\theta$ est l'angle de rotation (en radians), $\dot{\theta} = \omega$ la vitesse angulaire, $\ddot{\theta} = \alpha$ l'accélération angulaire.
\end{propriete}

\begin{methode}[Appliquer la relation fondamentale de la dynamique de rotation]
\begin{enumerate}
\item Définir le solide, l'\textbf{axe de rotation fixe} $(\Delta)$ et le référentiel galiléen.
\item Faire le bilan des forces extérieures et calculer le \textbf{moment de chaque force} par rapport à $(\Delta)$ : $\mathcal{M}_{\Delta}(\vec{F}) = \pm F \times d$, où $d$ est le bras de levier (distance de l'axe à la ligne d'action). Choisir un sens positif de rotation. Les forces dont la ligne d'action rencontre l'axe (poids si l'axe passe par $G$, réaction de l'axe) ont un moment nul.
\item Écrire la \textbf{relation fondamentale de la dynamique} :
\[ \sum \mathcal{M}_{\Delta}(\vec{F}_{\text{ext}}) = J_{\Delta}\,\ddot{\theta} \]
où $J_{\Delta}$ est le moment d'inertie du solide par rapport à l'axe.
\item Résoudre l'équation différentielle en $\theta(t)$ : si $\sum \mathcal{M}$ est constant, le mouvement est uniformément varié : $\ddot{\theta} = \text{cste}$, $\dot{\theta}(t) = \ddot{\theta}\,t + \dot{\theta}_0$, $\theta(t) = \tfrac{1}{2}\ddot{\theta}\,t^2 + \dot{\theta}_0 t + \theta_0$.
\item Exploiter : durée d'un tour, vitesse angulaire atteinte, nombre de tours effectués ($\theta$ en radians, $n = \dfrac{\theta}{2\pi}$).
\end{enumerate}
\textbf{Vérification homogénéité} : $[J] = \text{kg}\!\cdot\!\text{m}^2$, $[\mathcal{M}] = \text{N}\!\cdot\!\text{m} = \text{kg}\!\cdot\!\text{m}^2\!\cdot\!\text{s}^{-2}$, donc $[\ddot{\theta}] = \text{s}^{-2}$ (rad$\cdot$s$^{-2}$).
\end{methode}

\begin{exoresolu}[Volant d'une machine au village]
Un volant homogène de masse $M = \SI{40}{kg}$ et de rayon $R = \SI{0,30}{m}$, assimilé à un disque plein, tourne autour de son axe horizontal fixe $(\Delta)$. Un moteur lui applique un couple constant de moment $\mathcal{M}_m = \SI{6,0}{N.m}$. Les frottements de l'axe sont négligés. Le volant part du repos.
\begin{enumerate}
\item Calculer le moment d'inertie $J_{\Delta}$ du volant.
\item Déterminer son accélération angulaire $\ddot{\theta}$.
\item Au bout de combien de temps atteint-il la fréquence de rotation $f = \SI{300}{tr.min^{-1}}$ ? Combien de tours a-t-il alors effectués ?
\end{enumerate}
\tcblower
\textbf{1. Moment d'inertie.} Pour un disque plein autour de son axe de symétrie :
\[ J_{\Delta} = \frac{1}{2}MR^2 = \frac{1}{2}\times 40 \times (0{,}30)^2 = 20 \times 0{,}09 = \SI{1,8}{kg.m^2}. \]

\textbf{2. Accélération angulaire.} Bilan des moments par rapport à $(\Delta)$ : le poids et la réaction de l'axe ont des lignes d'action qui rencontrent l'axe, leurs moments sont nuls. Seul le couple moteur intervient (sens positif choisi dans le sens du couple) :
\[ \mathcal{M}_m = J_{\Delta}\,\ddot{\theta} \quad \Longrightarrow \quad \ddot{\theta} = \frac{\mathcal{M}_m}{J_{\Delta}} = \frac{6{,}0}{1{,}8} = \SI{3,33}{rad.s^{-2}}. \]

\textbf{3. Durée et nombre de tours.} Conversion de la fréquence : $f = \SI{300}{tr.min^{-1}} = \frac{300}{60} = \SI{5}{tr.s^{-1}}$. La vitesse angulaire visée est
\[ \dot{\theta} = 2\pi f = 2\pi \times 5 = 10\pi \approx \SI{31,4}{rad.s^{-1}}. \]
Le mouvement est uniformément accéléré depuis le repos ($\dot{\theta}_0 = 0$) : $\dot{\theta} = \ddot{\theta}\,t$, d'où
\[ t = \frac{\dot{\theta}}{\ddot{\theta}} = \frac{10\pi}{3{,}33} \approx \SI{9,43}{s}. \]
Angle balayé : $\theta = \tfrac{1}{2}\ddot{\theta}\,t^2 = \tfrac{1}{2}\times 3{,}33 \times (9{,}43)^2 \approx \SI{148}{rad}$. Nombre de tours :
\[ n = \frac{\theta}{2\pi} = \frac{148}{2\pi} \approx 23{,}6 \text{ tours}. \]
\textbf{Conclusion} : le volant atteint $\SI{300}{tr.min^{-1}}$ en environ $\SI{9,4}{s}$, après avoir effectué environ 24 tours.
\end{exoresolu}

\courssub{Théorème de l'énergie cinétique}

\begin{definition}[Énergie cinétique]
\begin{itemize}
\item \textbf{Translation} (point matériel ou solide en translation) : $E_c = \frac{1}{2}m v_G^2$.
\item \textbf{Rotation} autour d'un axe fixe $(\Delta)$ : $E_c = \frac{1}{2}J_{\Delta} \dot{\theta}^2$.
\item \textbf{Général} (translation + rotation) : $E_c = \frac{1}{2}m v_G^2 + \frac{1}{2}J_G \omega^2$ (théorème de König).
\end{itemize}
Unité SI : le joule ($\text{J} = \text{kg}\!\cdot\!\text{m}^2\!\cdot\!\text{s}^{-2}$).
\end{definition}

\begin{propriete}[Théorème de l'énergie cinétique]
Dans un référentiel galiléen, la variation de l'énergie cinétique d'un système entre deux états A et B est égale à la somme des travaux des forces extérieures (et des couples extérieurs en rotation) :
\[ \Delta E_c = E_c(B) - E_c(A) = \sum W_{AB}(\vec{F}_{\text{ext}}) + \sum W_{AB}(\mathcal{M}_{\text{ext}}). \]
\end{propriete}

\begin{methode}[Utiliser le théorème de l'énergie cinétique]
\begin{enumerate}
\item Définir le système, le référentiel galiléen, l'\textbf{état initial A} et l'\textbf{état final B}.
\item Faire le bilan de \textbf{toutes} les forces extérieures (et des couples, en rotation).
\item Calculer le \textbf{travail de chaque force} entre A et B :
\begin{itemize}
\item Force constante : $W_{AB}(\vec{F}) = \vec{F}\cdot\overrightarrow{AB} = F \times AB \times \cos\alpha$ ;
\item Poids : $W_{AB}(\vec{P}) = mg(z_A - z_B)$ (ne dépend que de la dénivellation) ;
\item Force perpendiculaire au déplacement (réaction normale sans frottement, tension d'un fil inextensible) : travail nul ;
\item Couple de moment $\mathcal{M}$ constant en rotation : $W = \mathcal{M}\,\Delta\theta$ ($\Delta\theta$ en radians).
\end{itemize}
\item Écrire le théorème :
\[ \Delta E_c = E_c(B) - E_c(A) = \sum W_{AB}(\vec{F}_{\text{ext}}) \]
avec $E_c = \tfrac{1}{2}mv_G^2$ (translation) ou $E_c = \tfrac{1}{2}J_{\Delta}\dot{\theta}^2$ (rotation).
\item Isoler la grandeur cherchée (vitesse, distance, norme d'une force de frottement...).
\end{enumerate}
\textbf{Intérêt} : ce théorème relie directement vitesse et déplacement \emph{sans passer par le temps} ; il est idéal quand la durée n'est pas demandée.
\end{methode}

\begin{exoresolu}[Freinage d'un camion sur la route de l'Ouest]
Un camion de masse $m = \SI{12}{t}$ roule à $v_0 = \SI{72}{km.h^{-1}}$ sur une route horizontale. Il freine jusqu'à l'arrêt ; la force de freinage totale $\vec{F}$ est supposée constante, de norme $F = \SI{6,0e4}{N}$. On prend $g = \SI{9,8}{m.s^{-2}}$.
\begin{enumerate}
\item Calculer la distance de freinage $d$.
\item Le camion roule maintenant à $v_0' = \SI{90}{km.h^{-1}}$. Sans refaire tous les calculs, donner la nouvelle distance de freinage.
\end{enumerate}
\tcblower
\textbf{1. Distance de freinage.} Système : le camion ; référentiel terrestre galiléen. Forces : le poids $\vec{P}$ et la réaction normale $\vec{R}_N$, perpendiculaires au déplacement horizontal (travail nul), et la force de freinage $\vec{F}$, opposée au mouvement : $W(\vec{F}) = -F\,d$.
Conversion : $v_0 = \SI{72}{km.h^{-1}} = \SI{20}{m.s^{-1}}$. Théorème de l'énergie cinétique entre le début du freinage (A) et l'arrêt (B, $v_B = 0$) :
\[ 0 - \frac{1}{2}mv_0^2 = -F\,d \quad \Longrightarrow \quad d = \frac{mv_0^2}{2F} = \frac{12\,000 \times 20^2}{2 \times 6{,}0\times 10^4} = \frac{4{,}8\times 10^6}{1{,}2\times 10^5} = \SI{40}{m}. \]

\textbf{2. À $\SI{90}{km.h^{-1}}$.} La distance est proportionnelle à $v_0^2$ :
\[ d' = d\left(\frac{v_0'}{v_0}\right)^2 = 40 \times \left(\frac{90}{72}\right)^2 = 40 \times 1{,}5625 = \SI{62,5}{m}. \]
\textbf{Conclusion} : la distance de freinage est $\SI{40}{m}$ à $\SI{72}{km.h^{-1}}$ et $\SI{62,5}{m}$ à $\SI{90}{km.h^{-1}}$ : une augmentation de vitesse de $25\,\%$ allonge le freinage de plus de $56\,\%$. C'est pourquoi les limitations de vitesse sauvent des vies sur nos routes.
\end{exoresolu}

\courssub{Énergie potentielle, énergie mécanique et conservation}

\begin{definition}[Forces conservatives et énergie potentielle]
Une force est \textbf{conservative} si son travail ne dépend que des positions initiale et finale, pas du chemin. On peut alors définir une \textbf{énergie potentielle} $E_p$ telle que $W_{AB}(\vec{F}) = E_p(A) - E_p(B) = -\Delta E_p$.
\begin{itemize}
\item \textbf{Poids} $\vec{P} = m\vec{g}$ (champ de pesanteur uniforme) : $E_p^{\text{pes}} = mgz$ (origine choisie arbitrairement, $z$ hauteur).
\item \textbf{Force élastique} (ressort de raideur $k$, allongement $x$) : $E_p^{\text{el}} = \frac{1}{2}k x^2$ (origine $x=0$ au repos).
\end{itemize}
\end{definition}

\begin{definition}[Énergie mécanique]
L'énergie mécanique d'un système est la somme de son énergie cinétique et de ses énergies potentielles :
\[ E_m = E_c + \sum E_p. \]
\end{definition}

\begin{propriete}[Conservation de l'énergie mécanique]
Si un système ne subit que des forces \textbf{conservatives} (poids, élastique) et des forces \textbf{qui ne travaillent pas} (réaction normale sans frottement, tension fil inextensible), alors son énergie mécanique est \textbf{constante} :
\[ E_m(A) = E_m(B) \quad \text{pour tous états A, B.} \]
Si des forces non conservatives travaillent (frottements $\vec{f}$), l'énergie mécanique varie :
\[ \Delta E_m = E_m(B) - E_m(A) = W_{AB}(\vec{f}_{\text{nc}}) < 0 \quad (\text{dissipation}). \]
\end{propriete}

\begin{methode}[Utiliser la conservation de l'énergie mécanique]
\begin{enumerate}
\item Définir le système et le référentiel galiléen ; repérer les états A et B.
\item \textbf{Vérifier la condition d'application} : le système est soumis uniquement à des forces conservatives (poids, tension élastique) ou à des forces qui ne travaillent pas (réaction normale sans frottement). Si des frottements travaillent, l'énergie mécanique \emph{ne se conserve pas}.
\item Choisir une \textbf{origine des énergies potentielles} (par exemple $E_p = 0$ au point le plus bas) et l'annoncer clairement.
\item Écrire l'énergie mécanique en chaque état :
\[ E_m = E_c + E_p = \frac{1}{2}mv^2 + mgz \quad (\text{ou } + \tfrac{1}{2}kx^2). \]
\item Écrire la conservation $E_m(A) = E_m(B)$ et isoler l'inconnue.
\item En présence de frottements $f$ sur un trajet AB, écrire la \textbf{variation} :
\[ \Delta E_m = E_m(B) - E_m(A) = W_{AB}(\vec{f}) = -f \times AB. \]
\end{enumerate}
\textbf{Réflexe} : pour un corps lâché sans vitesse d'une hauteur $h$ sans frottement, $v = \sqrt{2gh}$, indépendamment de la forme du trajet (pente, courbe, chute libre).
\end{methode}

\begin{exoresolu}[Le plongeur de la piscine olympique]
Un plongeur de masse $m = \SI{60}{kg}$ s'élance sans vitesse initiale depuis un plongeoir situé à la hauteur $h = \SI{5,0}{m}$ au-dessus de l'eau. On néglige la résistance de l'air et on prend $g = \SI{9,8}{m.s^{-2}}$.
\begin{enumerate}
\item Justifier la conservation de l'énergie mécanique.
\item Calculer la vitesse du plongeur à l'entrée dans l'eau.
\item En réalité, sa vitesse d'arrivée est $v' = \SI{9,5}{m.s^{-1}}$. Calculer le travail de la résistance de l'air.
\end{enumerate}
\tcblower
\textbf{1. Conservation.} Système : le plongeur ; référentiel terrestre galiléen. La seule force qui travaille est le poids, force conservative : l'énergie mécanique se conserve.

\textbf{2. Vitesse d'arrivée.} Origine des énergies potentielles : surface de l'eau ($z = 0$). État A (plongeoir) : $E_m(A) = mgh$ (vitesse nulle). État B (entrée dans l'eau) : $E_m(B) = \tfrac{1}{2}mv^2$.
\[ mgh = \frac{1}{2}mv^2 \quad \Longrightarrow \quad v = \sqrt{2gh} = \sqrt{2\times 9{,}8 \times 5{,}0} = \sqrt{98} \approx \SI{9,90}{m.s^{-1}}. \]

\textbf{3. Travail de la résistance de l'air.} L'énergie mécanique ne se conserve plus : sa variation égale le travail des forces non conservatives :
\[ W(\vec{f}_{\text{air}}) = E_m(B) - E_m(A) = \frac{1}{2}mv'^2 - mgh = \frac{1}{2}\times 60 \times (9{,}5)^2 - 60 \times 9{,}8 \times 5{,}0 \]
\[ W(\vec{f}_{\text{air}}) = 2\,707{,}5 - 2\,940 = \SI{-232,5}{J} \approx \SI{-2,3e2}{J}. \]
\textbf{Conclusion} : sans frottement, $v \approx \SI{9,9}{m.s^{-1}}$ ; la résistance de l'air dissipe environ $\SI{230}{J}$, ce qui ramène la vitesse réelle à $\SI{9,5}{m.s^{-1}}$. Le travail des frottements est négatif, comme attendu.
\end{exoresolu}

\courssub{Méthode 6 : Choisir entre méthode dynamique et méthode énergétique}

\begin{methode}[Choisir la bonne méthode de résolution]
\begin{enumerate}
\item \textbf{Identifier la question posée} :
\begin{itemize}
\item On demande une \textbf{durée}, une \textbf{équation horaire}, une \textbf{trajectoire}, une \textbf{force} (tension, réaction normale), un \textbf{moment} $\Rightarrow$ méthode \textbf{dynamique} (lois de Newton, projections, intégration, relation fondamentale rotation).
\item On demande une \textbf{vitesse en un point}, une \textbf{distance}, une \textbf{hauteur}, sans que le temps intervienne $\Rightarrow$ méthode \textbf{énergétique} (théorème de l'énergie cinétique ou conservation de $E_m$), beaucoup plus rapide.
\end{itemize}
\item Vérifier les conditions d'application : référentiel galiléen dans les deux cas ; conservation de $E_m$ seulement sans travail de frottements.
\item En rotation, même logique : relation $\sum \mathcal{M} = J\ddot{\theta}$ (dynamique, accès au temps) ou $\Delta(\tfrac{1}{2}J\dot{\theta}^2) = \sum W$ (énergétique).
\item Les deux méthodes peuvent se \textbf{contrôler mutuellement} : si le temps le permet, vérifier un résultat énergétique par la dynamique (ou inversement).
\end{enumerate}
\textbf{Réflexe Bac} : rédiger la méthode choisie en l'annonçant explicitement (« Appliquons le théorème de l'énergie cinétique entre A et B... ») : le correcteur valorise la démarche autant que le calcul.
\end{methode}

\begin{exoresolu}[Boule sur piste : les deux méthodes face à face]
Une boule de masse $m = \SI{0,50}{kg}$ glisse sans frottement sur une piste. Elle part du repos au point A situé à la hauteur $h = \SI{1,25}{m}$ et arrive au point B situé au niveau du sol. On prend $g = \SI{10}{m.s^{-2}}$.
\begin{enumerate}
\item Déterminer la vitesse en B par la méthode énergétique.
\item La portion AB est un plan incliné de longueur $L = \SI{2,5}{m}$. Retrouver la vitesse en B par la méthode dynamique, puis calculer la durée du parcours AB.
\end{enumerate}
\tcblower
\textbf{1. Méthode énergétique.} Référentiel terrestre galiléen. Forces : poids $\vec{P}$ (conservatif) et réaction normale $\vec{R}_N$ (perpendiculaire au déplacement, travail nul). L'énergie mécanique se conserve. Origine de $E_p$ au sol ($z=0$) :
\[ E_m(A) = E_m(B) \iff mgh = \frac{1}{2}mv_B^2 \iff v_B = \sqrt{2gh} = \sqrt{2\times 10 \times 1{,}25} = \sqrt{25} = \SI{5,0}{m.s^{-1}}. \]

\textbf{2. Méthode dynamique.} Le plan incliné fait un angle $\alpha$ avec l'horizontale tel que $\sin\alpha = \dfrac{h}{L} = \dfrac{1{,}25}{2{,}5} = 0{,}5$, donc $\alpha = 30^{\circ}$. Sur l'axe $(Ax)$ dirigé vers le bas de la pente, le théorème du centre d'inertie projeté donne :
\[ mg\sin\alpha = m\,a \quad \Longrightarrow \quad a = g\sin\alpha = 10 \times 0{,}5 = \SI{5,0}{m.s^{-2}}. \]
Mouvement uniformément accéléré depuis le repos : $v_B^2 = 2aL$, d'où
\[ v_B = \sqrt{2 \times 5{,}0 \times 2{,}5} = \sqrt{25} = \SI{5,0}{m.s^{-1}}. \]
On retrouve le même résultat, ce qui conforte la cohérence des deux lois. Durée du parcours :
\[ v_B = a\,t \quad \Longrightarrow \quad t = \frac{v_B}{a} = \frac{5{,}0}{5{,}0} = \SI{1,0}{s}. \]
\textbf{Conclusion} : la méthode énergétique donne $v_B = \SI{5,0}{m.s^{-1}}$ en une ligne, sans connaître la forme de la piste ; seule la méthode dynamique permet d'accéder à la durée $t = \SI{1,0}{s}$. Le choix de la méthode dépend de la grandeur recherchée.
\end{exoresolu}

\begin{savaistu}[Applications au Cameroun : Barrages et Énergie]
\begin{itemize}
\item \textbf{Barrages hydroélectriques (Edéa, Song Loulou, Lom Pangar, Nachtigal)} : L'eau du réservoir possède une énergie potentielle de pesanteur $E_p = mgh$. En tombant dans les conduites forcées, $E_p$ se convertit en énergie cinétique $E_c$ (vitesse de l'eau), qui actionne les turbines (rotation). Les alternateurs transforment l'énergie mécanique de rotation en énergie électrique transportée par le réseau SONATREL/ENEO. Le rendement global est de l'ordre de $85\text{--}90\,\%$ ; les pertes (frottements fluides, échauffement Joule) illustrent la non-conservation de $E_m$ dans les parties réelles.
\item \textbf{Télécommunications (MTN, Orange, CAMTEL)} : Les satellites géostationnaires (orbite circulaire $R \approx 42\,000$ km) sont maintenus par la force de gravitation (force centrale conservative). Leur vitesse orbitale $v = \sqrt{GM_T/R} \approx \SI{3,07}{km.s^{-1}}$ résulte de l'équilibre dynamique $\frac{GM_T m}{R^2} = m\frac{v^2}{R}$ (2ème loi + repère de Frenet).
\item \textbf{Sécurité routière} : L'exemple du camion montre que l'énergie cinétique (donc la distance de freinage) croît avec le \emph{carré} de la vitesse. C'est la base physique des limitations de vitesse sur nos axes Yaoundé-Douala, Bamenda-Enugu, etc.
\end{itemize}
\end{savaistu}

\begin{attention}[Les 5 erreurs qui coûtent des points au Bac]
\begin{enumerate}
\item \textbf{Oublier le bilan des forces et le référentiel.} Toute application des lois de Newton ou du théorème de l'énergie cinétique doit commencer par : système, référentiel galiléen, bilan des forces extérieures (avec schéma). Une équation $\sum\vec{F} = m\vec{a}$ écrite sans bilan est sanctionnée.
\item \textbf{Inventer une « force centrifuge » ou projeter le poids avec un mauvais angle.} Dans un référentiel galiléen, il n'y a pas de force centrifuge : c'est l'accélération qui est centripète. Sur un plan incliné d'angle $\alpha$, la composante motrice du poids est $mg\sin\alpha$ (et non $mg\cos\alpha$) : vérifier avec le cas limite $\alpha = 0$ (horizontal, composante nulle).
\item \textbf{Appliquer la conservation de l'énergie mécanique en présence de frottements.} Si des frottements travaillent, $E_m$ ne se conserve pas : il faut écrire $\Delta E_m = W(\vec{f})$. Toujours justifier la conservation avant de l'utiliser (« les seules forces qui travaillent sont conservatives »).
\item \textbf{Confondre les unités et les grandeurs de rotation.} Vitesses en $\si{km.h^{-1}}$ non converties ($\div 3{,}6$), fréquence en $\si{tr.min^{-1}}$ non convertie en vitesse angulaire ($\dot{\theta} = 2\pi f$ avec $f$ en $\si{tr.s^{-1}}$), angle en degrés au lieu de radians dans $W = \mathcal{M}\theta$ : trois sources classiques de résultats faux d'un facteur $3{,}6$, $2\pi$ ou $60$.
\item \textbf{Mal calculer un travail ou un moment.} Le travail du poids est $W = mg(z_A - z_B)$ : signe positif si le corps descend, négatif s'il monte ; une force perpendiculaire au déplacement ne travaille pas ; le moment d'une force dont la ligne d'action coupe l'axe est nul. Toujours conclure par une phrase avec la valeur numérique, son unité SI et un nombre cohérent de chiffres significatifs (généralement 2 ou 3).
\end{enumerate}
\end{attention}

\newpage

\coursec{Exercices}

\courssub{Je vérifie mes connaissances}

\begin{exercice}[QCM : les lois de Newton]
Pour chaque affirmation, choisir la (ou les) bonne(s) réponse(s) en justifiant brièvement.
\begin{enumerate}
\item Un solide pseudo-isolé dans un référentiel galiléen :
\begin{enumerate}
\item est nécessairement au repos ;
\item a un mouvement rectiligne uniforme ou est au repos ;
\item peut avoir un mouvement circulaire uniforme.
\end{enumerate}
\item Le théorème du centre d'inertie s'écrit, dans un référentiel galiléen :
\begin{enumerate}
\item $\sum \vec{F}_{\mathrm{ext}} = \vec{0}$ ;
\item $\sum \vec{F}_{\mathrm{ext}} = m\,\vec{a}_G$ ;
\item $\sum \vec{F}_{\mathrm{ext}} = \dfrac{d\vec{v}_G}{dt}$.
\end{enumerate}
\item L'unité SI du moment d'inertie est :
\begin{enumerate}
\item $\mathrm{kg\cdot m}$ ;
\item $\mathrm{kg\cdot m^2}$ ;
\item $\mathrm{kg\cdot m^2\cdot s^{-1}}$.
\end{enumerate}
\item L'énergie mécanique d'un système se conserve si :
\begin{enumerate}
\item le système est soumis uniquement à des forces conservatives (ou à des forces dont le travail est nul) ;
\item la vitesse du système est constante ;
\item le système est isolé de toute interaction.
\end{enumerate}
\end{enumerate}
\end{exercice}

\begin{exercice}[Vrai ou faux]
Répondre par vrai ou faux en justifiant chaque réponse par une loi ou un contre-exemple.
\begin{enumerate}
\item Si la somme des forces extérieures appliquées à un solide est nulle, alors le solide est immobile.
\item Dans le repère de Frenet, l'accélération normale d'un point en mouvement circulaire uniforme est nulle.
\item D'après la troisième loi de Newton, la force qu'exerce la Terre sur un corps a même norme que la force qu'exerce ce corps sur la Terre.
\item Le travail du poids d'un corps dépend du chemin suivi entre deux points.
\item En présence de frottements solides, l'énergie mécanique d'un système diminue au cours du mouvement.
\end{enumerate}
\end{exercice}

\begin{exercice}[Définitions et expressions]
\begin{enumerate}
\item Définir un référentiel galiléen et citer deux référentiels utilisés en mécanique terrestre, en précisant leur caractère galiléen approché.
\item Énoncer le principe des actions réciproques (troisième loi de Newton).
\item Donner l'expression du moment d'inertie $J_\Delta$ d'un point matériel de masse $m$ situé à la distance $r$ d'un axe $\Delta$, puis celle d'un cerceau de masse $M$ et de rayon $R$ par rapport à son axe.
\item Donner l'expression de l'énergie cinétique d'un solide en rotation autour d'un axe fixe à la vitesse angulaire $\dot{\theta}$.
\end{enumerate}
\end{exercice}

\begin{exercice}[Calculs directs]
\begin{enumerate}
\item Une mototaxi de masse totale $m = \SI{180}{kg}$ passe de $\SI{0}{}$ à $\SI{54}{km.h^{-1}}$ en $\SI{6,0}{s}$ sur une route rectiligne horizontale. Calculer son accélération moyenne, puis la résultante des forces extérieures supposée constante.
\item Un point matériel décrit un cercle de rayon $R = \SI{2,0}{m}$ à la vitesse constante $v = \SI{6,0}{m.s^{-1}}$. Calculer son accélération normale.
\item Un volant de moment d'inertie $J_\Delta = \SI{0,40}{kg.m^2}$ tourne à $\omega = \SI{20}{rad.s^{-1}}$. Calculer son énergie cinétique.
\item Un sac de ciment de masse $m = \SI{50}{kg}$ est hissé de $\SI{12}{m}$. Calculer la variation de son énergie potentielle de pesanteur. On prendra $g = \SI{10}{N.kg^{-1}}$.
\end{enumerate}
\end{exercice}

\courssub{Je m'entraîne}

\begin{exercice}[Chariot sur plan incliné]
Un chariot de masse $m = \SI{2,0}{kg}$ glisse sans frottement sur un plan incliné d'un angle $\alpha = 30^{\circ}$ sur l'horizontale. On prend $g = \SI{10}{N.kg^{-1}}$.
\begin{enumerate}
\item Faire le bilan des forces extérieures et les représenter sur un schéma.
\item En appliquant le théorème du centre d'inertie, déterminer l'accélération du chariot.
\item Le chariot part du repos. Quelle distance parcourt-il en $\SI{2,0}{s}$ ? Quelle est alors sa vitesse ?
\item Retrouver cette vitesse par le théorème de l'énergie cinétique.
\end{enumerate}
\end{exercice}

\begin{exercice}[Pendule simple]
Un pendule simple est constitué d'une bille de masse $m = \SI{100}{g}$ suspendue à un fil inextensible de longueur $\ell = \SI{0,80}{m}$, de masse négligeable. On écarte le fil d'un angle $\theta_0 = 60^{\circ}$ par rapport à la verticale et on le lâche sans vitesse initiale. On néglige les frottements et on prend $g = \SI{10}{N.kg^{-1}}$.
\begin{enumerate}
\item Exprimer, en fonction de $\ell$ et $\theta_0$, la dénivellation $h$ entre le point de lâcher et le point le plus bas de la trajectoire.
\item En appliquant la conservation de l'énergie mécanique, calculer la vitesse $v_B$ de la bille au point le plus bas.
\item En projetant le théorème du centre d'inertie sur l'axe normal du repère de Frenet au point le plus bas, calculer la tension $T$ du fil à cet instant.
\end{enumerate}
\end{exercice}

\begin{exercice}[Volant entraîné par un couple]
Un volant de machine, assimilé à un disque homogène de masse $M = \SI{20}{kg}$ et de rayon $R = \SI{0,30}{m}$, tourne autour de son axe fixe $\Delta$. Son moment d'inertie est $J_\Delta = \dfrac{1}{2}MR^2$. On lui applique un couple moteur de moment constant $\mathcal{M}_m = \SI{5,4}{N.m}$ ; les frottements sont négligés.
\begin{enumerate}
\item Calculer $J_\Delta$.
\item En appliquant la relation fondamentale de la dynamique de rotation, déterminer l'accélération angulaire $\ddot{\theta}$ du volant.
\item Le volant part du repos. Quelle est sa vitesse angulaire après $\SI{10}{s}$ ? Exprimer cette vitesse en tours par minute.
\item Retrouver cette vitesse angulaire à l'aide du théorème de l'énergie cinétique, en calculant d'abord l'angle total balayé.
\end{enumerate}
\end{exercice}

\begin{exercice}[Virage circulaire]
Une voiture de masse $m = \SI{1000}{kg}$ aborde un virage circulaire horizontal de rayon $R = \SI{50}{m}$. Le coefficient de frottement statique entre les pneus et la chaussée est $\mu_s = 0{,}60$ : la force de frottement maximale vaut $f_{\max} = \mu_s\,N$, où $N$ est la réaction normale de la route. On prend $g = \SI{10}{N.kg^{-1}}$.
\begin{enumerate}
\item Faire le bilan des forces et montrer que c'est la force de frottement qui joue le rôle de force centripète.
\item En utilisant l'expression de l'accélération normale dans le repère de Frenet, établir l'expression de la vitesse maximale $v_{\max}$ avant dérapage.
\item Calculer $v_{\max}$ en $\mathrm{m.s^{-1}}$ puis en $\mathrm{km.h^{-1}}$. Commenter la valeur obtenue au regard de la signalisation routière.
\end{enumerate}
\end{exercice}

\courssub{Je me prépare au Bac}

\begin{exercice}[Type Bac : solide lancé sur un plan incliné avec frottements]
Sur un chantier de construction à Yaoundé, un ouvrier lance vers le haut d'un plan incliné un bloc de masse $m = \SI{5,0}{kg}$ avec une vitesse initiale $v_0 = \SI{6,0}{m.s^{-1}}$. Le plan est incliné d'un angle $\alpha = 30^{\circ}$ sur l'horizontale. Les frottements sont équivalents à une force unique $\vec{f}$, constante, de norme $f = \SI{8,0}{N}$, opposée au mouvement. On prend $g = \SI{10}{N.kg^{-1}}$.
\begin{enumerate}
\item Représenter sur un schéma les forces extérieures appliquées au bloc pendant la montée.
\item En appliquant le théorème du centre d'inertie, établir l'expression littérale de l'accélération $a$ du bloc, puis calculer sa valeur. Préciser la nature du mouvement.
\item Déterminer, par le théorème de l'énergie cinétique, la distance $d$ parcourue par le bloc sur le plan avant de s'arrêter.
\item En déduire l'altitude $h$ atteinte par le bloc au-dessus de son point de lancement.
\item Retrouver la valeur de $h$ en écrivant la variation d'énergie mécanique du bloc entre le point de lancement et le point d'arrêt, et en la reliant au travail des frottements.
\item Le bloc redescend ensuite. Montrer, sans nouveau calcul complet, que sa vitesse au bas du plan est inférieure à $v_0$, et proposer une méthode énergétique pour la calculer.
\end{enumerate}
\end{exercice}

\begin{exercice}[Type Bac : poulie et masse descendante]
Une machine de levage utilisée sur le chantier du barrage de Nachtigal est modélisée par une poulie homogène de rayon $R = \SI{0,20}{m}$ et de moment d'inertie $J_\Delta = \SI{0,10}{kg.m^2}$, mobile sans frottement autour de son axe horizontal fixe $\Delta$. Un fil inextensible, de masse négligeable, enroulé sur la poulie sans glisser, supporte à son extrémité une charge de masse $m = \SI{10}{kg}$. On abandonne le système sans vitesse initiale et la charge descend. On prend $g = \SI{10}{N.kg^{-1}}$.
\begin{enumerate}
\item Faire le bilan des forces appliquées à la charge, puis le bilan des moments appliqués à la poulie.
\item Écrire le théorème du centre d'inertie pour la charge et la relation fondamentale de la dynamique de rotation pour la poulie, en notant $T$ la tension du fil.
\item En utilisant la condition de non-glissement du fil sur la poulie ($a = R\,\ddot{\theta}$), établir l'expression de l'accélération $a$ de la charge en fonction de $m$, $g$, $J_\Delta$ et $R$. Calculer sa valeur.
\item En déduire la tension $T$ du fil. Comparer $T$ au poids de la charge et commenter.
\item La charge descend d'une hauteur $h = \SI{4,0}{m}$. Calculer sa vitesse $v$ en bas de la course par le théorème de l'énergie cinétique appliqué à la charge.
\item Vérification énergétique : en considérant le système \{charge + poulie + Terre\}, écrire la conservation de l'énergie mécanique entre le départ et la fin de la descente, et retrouver la valeur de $v$.
\end{enumerate}
\end{exercice}

\begin{exercice}[Type Bac : enfant sur toboggan et freinage d'un volant]
\textbf{Partie A : descente d'un toboggan.} Dans un jardin d'enfants de Douala, un enfant de masse $m = \SI{25}{kg}$ descend un toboggan dont le sommet $A$ est à la hauteur $h = \SI{3,0}{m}$ au-dessus du sol. Il part de $A$ sans vitesse initiale et arrive au point bas $B$ avec la vitesse $v_B = \SI{5,0}{m.s^{-1}}$. On prend $g = \SI{10}{N.kg^{-1}}$.
\begin{enumerate}
\item Calculer l'énergie mécanique de l'enfant en $A$ et en $B$ (on prendra le niveau du sol comme référence des énergies potentielles).
\item L'énergie mécanique se conserve-t-elle au cours de la descente ? Justifier.
\item Calculer le travail $W(\vec{f})$ de la résultante des forces de frottement entre $A$ et $B$.
\item La longueur de la glissière est $L = \SI{6,0}{m}$. En assimilant les frottements à une force constante $\vec{f}$ colinéaire au déplacement, calculer la norme $f$ de cette force.
\item Quelle serait la vitesse en $B$ en l'absence totale de frottements ? Commenter l'écart avec la valeur réelle.
\end{enumerate}
\textbf{Partie B : freinage d'un volant.} Le volant d'un groupe électrogène, de moment d'inertie $J_\Delta = \SI{2,0}{kg.m^2}$, tourne à la vitesse angulaire initiale $\omega_0 = \SI{100}{rad.s^{-1}}$. À la coupure du moteur, il n'est plus soumis qu'à un couple de freinage de moment constant $\mathcal{M}_f = -\SI{10}{N.m}$.
\begin{enumerate}
\setcounter{enumi}{5}
\item En appliquant la relation fondamentale de la dynamique de rotation, déterminer l'accélération angulaire du volant et la durée du freinage jusqu'à l'arrêt.
\item À l'aide du théorème de l'énergie cinétique, calculer l'angle total $\theta$ balayé par le volant pendant le freinage, puis le nombre de tours effectués avant l'arrêt.
\end{enumerate}
\end{exercice}



\newpage

\coursec{Corrigés des exercices}

\courssub{Je vérifie mes connaissances}

\begin{corrige}[Exercice 1 --- QCM : les lois de Newton]
\begin{enumerate}
\item \textbf{Réponse b.} D'après le principe d'inertie (et sa réciproque), si $\sum \vec{F}_{\mathrm{ext}} = \vec{0}$, alors $\vec{v}_G = \overrightarrow{\mathrm{cste}}$ : le centre d'inertie est au repos ou en mouvement rectiligne uniforme. La réponse c est exclue : un mouvement circulaire uniforme possède une accélération normale $a_N = v^2/R \neq 0$, donc exige une résultante de forces non nulle.
\item \textbf{Réponse b.} La réponse c est dimensionnellement fausse : $\dfrac{d\vec{v}_G}{dt}$ est une accélération ($\mathrm{m.s^{-2}}$), pas une force ; il manque le facteur $m$ pour obtenir des newtons. Rappelons l'analyse dimensionnelle : $[F] = \mathrm{kg \cdot m \cdot s^{-2}}$.
\item \textbf{Réponse b.} Pour un point matériel de masse $m$ situé à la distance $r$ de l'axe, $J_\Delta = m r^2$, d'où $[J_\Delta] = \mathrm{kg \cdot m^2}$.
\item \textbf{Réponse a.} La conservation de l'énergie mécanique exige que seules des forces conservatives travaillent, ou que les forces non conservatives aient un travail nul (exemple : tension d'un fil dont le point d'application est fixe, ou réaction normale d'un support fixe). La réponse b est un piège classique : un parachutiste en chute à vitesse limite constante a une énergie cinétique constante mais une énergie potentielle qui diminue : son énergie mécanique diminue, dissipée par les frottements de l'air.
\end{enumerate}
\end{corrige}

\begin{corrige}[Exercice 2 --- Vrai ou faux]
\begin{enumerate}
\item \textbf{Faux.} D'après le principe d'inertie, le centre d'inertie est au repos \emph{ou} en mouvement rectiligne uniforme. Contre-exemple : un palet lancé sur une patinoire horizontale (frottements négligeables) garde sa vitesse.
\item \textbf{Faux.} C'est l'accélération \emph{tangentielle} qui est nulle, car la vitesse est constante ($a_T = \dot{v} = 0$). L'accélération normale vaut $a_N = v^2/R \neq 0$ : c'est elle qui courbe la trajectoire.
\item \textbf{Vrai.} Le principe des actions réciproques impose $\vec{F}_{\mathrm{Terre/corps}} = -\vec{F}_{\mathrm{corps/Terre}}$ : mêmes normes, même droite d'action, sens opposés. La Terre, de masse considérable, acquiert une accélération indétectable.
\item \textbf{Faux.} Le poids est une force conservative : $W_{AB}(\vec{P}) = mg(z_A - z_B)$ ne dépend que de la dénivellation entre $A$ et $B$, jamais du chemin suivi.
\item \textbf{Vrai.} Le travail des frottements solides est résistant ($W(\vec{f}) < 0$) et $\Delta E_m = W(\vec{f}) < 0$ : l'énergie mécanique diminue, convertie en énergie thermique.
\end{enumerate}
\end{corrige}

\begin{corrige}[Exercice 3 --- Définitions et expressions]
\begin{enumerate}
\item Un référentiel galiléen est un référentiel dans lequel le principe d'inertie est vérifié : tout point matériel pseudo-isolé y est au repos ou en mouvement rectiligne uniforme. Le référentiel héliocentrique (de Copernic) en est une excellente approximation ; le référentiel terrestre (lié au sol du laboratoire) est galiléen de façon approchée pour des mouvements de courte durée, la rotation de la Terre étant négligeable à cette échelle.
\item Si un solide $A$ exerce sur un solide $B$ la force $\vec{F}_{A/B}$, alors $B$ exerce sur $A$ la force $\vec{F}_{B/A}$ telle que $\vec{F}_{A/B} = -\vec{F}_{B/A}$ : ces deux forces ont même droite d'action, même norme et des sens opposés. Attention : elles s'exercent sur des corps différents et ne s'annulent donc jamais dans un bilan de forces portant sur un seul système.
\item Point matériel : $J_\Delta = m r^2$. Cerceau (toute la masse à la distance $R$ de l'axe) : $J_\Delta = M R^2$. Cylindre plein (ou disque homogène) : $J_\Delta = \dfrac{1}{2}MR^2$.
\item $E_c = \dfrac{1}{2} J_\Delta\, \dot{\theta}^2$, avec $\dot{\theta}$ en $\mathrm{rad.s^{-1}}$. Vérification dimensionnelle : $[\tfrac{1}{2}J_\Delta\dot{\theta}^2] = \mathrm{kg \cdot m^2 \cdot s^{-2}} = \mathrm{J}$.
\end{enumerate}
\end{corrige}

\begin{corrige}[Exercice 4 --- Calculs directs]
\begin{enumerate}
\item Conversion : $v = 54/3{,}6 = \SI{15}{m.s^{-1}}$. Accélération moyenne : $a = \dfrac{\Delta v}{\Delta t} = \dfrac{15}{6{,}0} = \SI{2,5}{m.s^{-2}}$. Théorème du centre d'inertie : $\sum F = ma = 180 \times 2{,}5 = \SI{4,5e2}{N}$.
\item $a_N = \dfrac{v^2}{R} = \dfrac{6{,}0^2}{2{,}0} = \SI{18}{m.s^{-2}}$, dirigée vers le centre du cercle (accélération centripète).
\item $E_c = \dfrac{1}{2} J_\Delta \omega^2 = \dfrac{1}{2} \times 0{,}40 \times 20^2 = \SI{80}{J}$.
\item $\Delta E_p = mg\,\Delta z = 50 \times 10 \times 12 = \SI{6,0e3}{J}$. La variation est positive : l'énergie potentielle augmente lorsque le sac monte.
\end{enumerate}
\end{corrige}

\courssub{Je m'entraîne}

\begin{corrige}[Exercice 5 --- Chariot sur plan incliné]
\begin{enumerate}
\item Bilan des forces extérieures appliquées au chariot : le poids $\vec{P} = m\vec{g}$ (vertical, vers le bas) et la réaction $\vec{R}$ du plan (normale au plan, car les frottements sont négligés). Le référentiel terrestre est supposé galiléen.
\item Projection du théorème du centre d'inertie $\vec{P} + \vec{R} = m\vec{a}$ sur l'axe $(Ox)$ parallèle au plan, orienté vers le bas de la pente : $mg\sin\alpha = ma$, d'où
\[ a = g\sin\alpha = 10 \times \sin 30^{\circ} = \SI{5,0}{m.s^{-2}}. \]
L'accélération est constante : le mouvement est rectiligne uniformément accéléré.
\item $x(t) = \dfrac{1}{2}at^2 = \dfrac{1}{2} \times 5{,}0 \times (2{,}0)^2 = \SI{10}{m}$ ; $v = at = 5{,}0 \times 2{,}0 = \SI{10}{m.s^{-1}}$.
\item Entre le départ et $t = \SI{2,0}{s}$, seul le poids travaille ($\vec{R} \perp$ déplacement) : $W(\vec{P}) = mg\,x\sin\alpha$. Le théorème de l'énergie cinétique donne
\[ \frac{1}{2}mv^2 - 0 = mg\,x\sin\alpha \quad\Rightarrow\quad v = \sqrt{2gx\sin\alpha} = \sqrt{2 \times 10 \times 10 \times 0{,}5} = \SI{10}{m.s^{-1}}. \]
Les deux méthodes donnent le même résultat : elles sont cohérentes. La méthode énergétique est ici plus rapide car elle évite le calcul de l'accélération et l'intégration.
\end{enumerate}
\end{corrige}

\begin{corrige}[Exercice 6 --- Pendule simple]
\begin{enumerate}
\item La bille descend d'une hauteur $h = \ell - \ell\cos\theta_0 = \ell(1-\cos\theta_0)$. Numériquement : $h = 0{,}80 \times (1 - 0{,}5) = \SI{0,40}{m}$.
\item La tension du fil, perpendiculaire à la trajectoire, ne travaille pas ; seul le poids (force conservative) travaille : l'énergie mécanique se conserve. Entre le point de lâcher $A$ ($v_A = 0$) et le point bas $B$ :
\[ mgh = \frac{1}{2}mv_B^2 \quad\Rightarrow\quad v_B = \sqrt{2gh} = \sqrt{2 \times 10 \times 0{,}40} = \sqrt{8} \approx \SI{2,8}{m.s^{-1}}. \]
\item Au point le plus bas, l'axe normal $\vec{n}$ du repère de Frenet est vertical, dirigé vers le haut (vers le point de suspension). Projection du théorème du centre d'inertie sur $\vec{n}$ :
\[ T - mg = m\,a_N = \frac{mv_B^2}{\ell} \quad\Rightarrow\quad T = m\left(g + \frac{v_B^2}{\ell}\right). \]
Numériquement : $T = 0{,}100 \times \left(10 + \dfrac{8}{0{,}80}\right) = 0{,}100 \times 20 = \SI{2,0}{N}$. On remarque que $T > mg$ : le fil est plus tendu au passage au point bas qu'à l'équilibre statique, car il doit en plus fournir la force centripète.
\end{enumerate}
\end{corrige}

\begin{corrige}[Exercice 7 --- Volant entraîné par un couple]
\begin{enumerate}
\item $J_\Delta = \dfrac{1}{2}MR^2 = \dfrac{1}{2} \times 20 \times (0{,}30)^2 = \SI{0,90}{kg.m^2}$.
\item Le poids et la réaction de l'axe ont un moment nul par rapport à $\Delta$ (leurs droites d'action rencontrent l'axe). La relation fondamentale de la dynamique de rotation s'écrit :
\[ \mathcal{M}_m = J_\Delta\,\ddot{\theta} \quad\Rightarrow\quad \ddot{\theta} = \frac{\mathcal{M}_m}{J_\Delta} = \frac{5{,}4}{0{,}90} = \SI{6,0}{rad.s^{-2}}. \]
\item $\omega = \ddot{\theta}\,t = 6{,}0 \times 10 = \SI{60}{rad.s^{-1}}$. En tours par minute : $n = \dfrac{60\,\omega}{2\pi} = \dfrac{60 \times 60}{2\pi} \approx \SI{5,7e2}{tr.min^{-1}}$.
\item Angle balayé : $\theta = \dfrac{1}{2}\ddot{\theta}\,t^2 = \dfrac{1}{2} \times 6{,}0 \times 10^2 = \SI{3,0e2}{rad}$. Travail du couple : $W = \mathcal{M}_m\,\theta = 5{,}4 \times 300 = \SI{1,62e3}{J}$. Le théorème de l'énergie cinétique donne
\[ \frac{1}{2}J_\Delta\,\omega^2 - 0 = W \quad\Rightarrow\quad \omega = \sqrt{\frac{2W}{J_\Delta}} = \sqrt{\frac{2 \times 1620}{0{,}90}} = \sqrt{3600} = \SI{60}{rad.s^{-1}}. \]
Résultat identique : la méthode énergétique confirme la méthode dynamique.
\end{enumerate}
\end{corrige}

\begin{corrige}[Exercice 8 --- Virage circulaire]
\begin{enumerate}
\item Forces appliquées à la voiture : le poids $\vec{P}$ (vertical), la réaction normale $\vec{N}$ de la route (verticale, vers le haut) et la force de frottement $\vec{f}$ (horizontale, dirigée vers le centre du virage). Sur la verticale, pas de mouvement : $N = mg$. La trajectoire étant circulaire, l'accélération possède une composante normale $a_N = v^2/R$ dirigée vers le centre : seule $\vec{f}$, horizontale et centripète, peut la fournir. C'est donc le frottement qui joue le rôle de force centripète.
\item Projection du théorème du centre d'inertie sur l'axe normal : $f = \dfrac{mv^2}{R}$. L'adhérence impose $f \leq f_{\max} = \mu_s N = \mu_s mg$. La limite de dérapage est atteinte pour
\[ \frac{mv_{\max}^2}{R} = \mu_s mg \quad\Rightarrow\quad v_{\max} = \sqrt{\mu_s\,g\,R}. \]
On remarque que la masse n'intervient pas : une mototaxi et un camion ont la même vitesse limite théorique dans ce virage.
\item $v_{\max} = \sqrt{0{,}60 \times 10 \times 50} = \sqrt{300} \approx \SI{17}{m.s^{-1}}$, soit $17{,}3 \times 3{,}6 \approx \SI{62}{km.h^{-1}}$. C'est cohérent avec la limitation à $\SI{60}{km.h^{-1}}$ souvent signalée avant les virages serrés ; sur chaussée mouillée, $\mu_s$ diminue fortement, donc $v_{\max}$ aussi : il faut ralentir.
\end{enumerate}
\end{corrige}

\courssub{Je me prépare au Bac}

\begin{corrige}[Exercice 9 --- Type Bac : solide lancé sur un plan incliné avec frottements]
\begin{enumerate}
\item Forces pendant la montée : le poids $\vec{P}$ (vertical, vers le bas), la réaction normale $\vec{N}$ (perpendiculaire au plan) et la force de frottement $\vec{f}$ (parallèle au plan, dirigée vers le bas de la pente, car opposée au mouvement ascendant). Le référentiel terrestre est supposé galiléen.
\item Axe $(Ox)$ orienté vers le haut de la pente. Projection du théorème du centre d'inertie :
\[ -mg\sin\alpha - f = ma \quad\Rightarrow\quad a = -\left(g\sin\alpha + \frac{f}{m}\right). \]
Numériquement : $a = -\left(10 \times 0{,}5 + \dfrac{8{,}0}{5{,}0}\right) = -(5{,}0 + 1{,}6) = \SI{-6,6}{m.s^{-2}}$. L'accélération est constante et de sens opposé à la vitesse : le mouvement est rectiligne uniformément décéléré.
\item Entre le lancement (vitesse $v_0$) et l'arrêt (vitesse nulle) après la distance $d$, le théorème de l'énergie cinétique s'écrit :
\[ 0 - \frac{1}{2}mv_0^2 = -mg\,d\sin\alpha - f\,d \quad\Rightarrow\quad d = \frac{v_0^2}{2\left(g\sin\alpha + f/m\right)} = \frac{36}{2 \times 6{,}6} \approx \SI{2,7}{m}. \]
\item $h = d\sin\alpha = 2{,}73 \times 0{,}5 \approx \SI{1,4}{m}$.
\item La variation d'énergie mécanique égale le travail des forces non conservatives :
\[ \Delta E_m = W(\vec{f}) \quad\Rightarrow\quad mgh - \frac{1}{2}mv_0^2 = -f\,d = -f\,\frac{h}{\sin\alpha}. \]
D'où $h\left(mg + \dfrac{f}{\sin\alpha}\right) = \dfrac{1}{2}mv_0^2$, soit
\[ h = \frac{\tfrac{1}{2}mv_0^2}{mg + f/\sin\alpha} = \frac{0{,}5 \times 5{,}0 \times 36}{50 + 16} = \frac{90}{66} \approx \SI{1,4}{m}. \]
On retrouve la même valeur : les deux approches sont équivalentes.
\item Pendant la descente, les frottements sont encore résistants : sur l'aller-retour, $W(\vec{f}) = -2fd < 0$, donc l'énergie mécanique au retour est inférieure à celle du départ. Comme le bloc revient à la même altitude ($E_p$ identique), son énergie cinétique, donc sa vitesse, est plus faible : $v_{\mathrm{retour}} < v_0$. Calcul par le théorème de l'énergie cinétique sur la descente : $\dfrac{1}{2}mv^2 = mgh - f\,d$, d'où
\[ v = \sqrt{2gh - \frac{2fd}{m}} = \sqrt{2 \times 10 \times 1{,}36 - \frac{2 \times 8{,}0 \times 2{,}73}{5{,}0}} \approx \sqrt{27{,}3 - 8{,}7} \approx \SI{4,3}{m.s^{-1}} < v_0. \]
\end{enumerate}
\textbf{Barème indicatif (sur 6 points) :} schéma des forces : 0,5 ; équation de la dynamique et valeur de $a$ : 1 ; distance $d$ : 1 ; altitude $h$ : 0,5 ; retrouver $h$ par $\Delta E_m = W(\vec{f})$ : 1,5 ; argument énergétique et calcul de $v$ : 1,5.

\textbf{Points de vigilance :} citer explicitement le référentiel terrestre supposé galiléen avant d'appliquer le théorème du centre d'inertie ; orienter l'axe et justifier le signe de chaque projection ; préciser que $\vec{N}$ ne travaille pas ; pour la question 5, écrire $\Delta E_m = W_{\mathrm{non\,conservatives}}$ et non $\Delta E_c$ ; conserver les valeurs exactes intermédiaires et arrondir seulement le résultat final.
\end{corrige}

\begin{corrige}[Exercice 10 --- Type Bac : poulie et masse descendante]
\begin{enumerate}
\item Charge : le poids $\vec{P} = m\vec{g}$ (vers le bas) et la tension $\vec{T}$ du fil (vers le haut). Poulie : son poids et la réaction de l'axe ont des moments nuls par rapport à $\Delta$ (leurs droites d'action coupent l'axe) ; seule la tension du fil exerce un moment moteur $\mathcal{M} = TR$.
\item Charge (axe vertical orienté vers le bas, sens du mouvement) : $mg - T = ma$. Poulie : $TR = J_\Delta\,\ddot{\theta}$.
\item Avec $\ddot{\theta} = a/R$ (non-glissement du fil sur la poulie), la seconde équation donne $T = \dfrac{J_\Delta\,a}{R^2}$. En reportant :
\[ mg - \frac{J_\Delta}{R^2}\,a = ma \quad\Rightarrow\quad a = \frac{mg}{m + J_\Delta/R^2}. \]
Numériquement : $\dfrac{J_\Delta}{R^2} = \dfrac{0{,}10}{0{,}20^2} = \SI{2,5}{kg}$, d'où $a = \dfrac{100}{12{,}5} = \SI{8,0}{m.s^{-2}}$. On vérifie que $a < g$ : une partie du poids sert à mettre la poulie en rotation.
\item $T = m(g-a) = 10 \times (10 - 8{,}0) = \SI{20}{N}$. Le poids vaut $mg = \SI{100}{N}$ : $T < mg$, cohérent avec une charge accélérée vers le bas (la résultante des forces doit être dirigée vers le bas).
\item Théorème de l'énergie cinétique appliqué à la charge entre le départ et la fin de la descente :
\[ \frac{1}{2}mv^2 - 0 = (mg - T)\,h = mah \quad\Rightarrow\quad v = \sqrt{2ah} = \sqrt{2 \times 8{,}0 \times 4{,}0} = \SI{8,0}{m.s^{-1}}. \]
\item Le système \{charge + poulie + Terre\} n'est soumis qu'au poids (conservatif) et à des forces de travail nul (réaction de l'axe ; tension intérieure au système) : l'énergie mécanique se conserve. La perte d'énergie potentielle alimente les énergies cinétiques de translation et de rotation :
\[ mgh = \frac{1}{2}mv^2 + \frac{1}{2}J_\Delta\,\omega^2, \quad \text{avec } \omega = \frac{v}{R}. \]
D'où $v = \sqrt{\dfrac{2mgh}{m + J_\Delta/R^2}} = \sqrt{\dfrac{2 \times 10 \times 10 \times 4{,}0}{12{,}5}} = \sqrt{64} = \SI{8,0}{m.s^{-1}}$. Vérification : $mgh = \SI{400}{J}$ ; $\tfrac{1}{2}mv^2 = \SI{320}{J}$ ; $\tfrac{1}{2}J_\Delta\omega^2 = \tfrac{1}{2} \times 0{,}10 \times 40^2 = \SI{80}{J}$ ; $320 + 80 = \SI{400}{J}$ : le bilan est exact.
\end{enumerate}
\textbf{Barème indicatif (sur 7 points) :} bilans de forces et de moments : 1 ; deux équations de la dynamique : 1 ; expression littérale et valeur de $a$ : 1,5 ; tension $T$ et comparaison : 1 ; vitesse $v$ par l'énergie cinétique : 1 ; conservation de $E_m$ et vérification : 1,5.

\textbf{Points de vigilance :} ne pas écrire $T = mg$ (la charge n'est pas en équilibre) ; justifier le moment nul du poids de la poulie (droite d'action passant par l'axe) ; utiliser la condition de non-glissement $a = R\ddot{\theta}$ et non $\omega = v/R$ dans les équations du mouvement ; dans le bilan énergétique, ne pas oublier le terme de rotation $\tfrac{1}{2}J_\Delta\omega^2$.
\end{corrige}

\begin{corrige}[Exercice 11 --- Type Bac : enfant sur toboggan et freinage d'un volant]
\textbf{Partie A.}
\begin{enumerate}
\item En $A$ : $E_m(A) = 0 + mgh = 25 \times 10 \times 3{,}0 = \SI{7,5e2}{J}$. En $B$ ($z_B = 0$) : $E_m(B) = \dfrac{1}{2}mv_B^2 = \dfrac{1}{2} \times 25 \times 5{,}0^2 = \SI{312,5}{J} \approx \SI{3,1e2}{J}$.
\item $E_m(B) < E_m(A)$ : l'énergie mécanique \textbf{ne se conserve pas}. Il existe donc des forces non conservatives de travail non nul : les frottements de l'enfant sur la glissière.
\item $\Delta E_m = W(\vec{f})$, d'où $W(\vec{f}) = 312{,}5 - 750 = \SI{-437,5}{J} \approx \SI{-4,4e2}{J}$. Le travail est négatif : les frottements sont résistants.
\item $W(\vec{f}) = -fL$, d'où $f = \dfrac{|W(\vec{f})|}{L} = \dfrac{437{,}5}{6{,}0} \approx \SI{73}{N}$.
\item Sans frottements, conservation de l'énergie mécanique : $mgh = \dfrac{1}{2}mv_B'^2$, soit $v_B' = \sqrt{2gh} = \sqrt{60} \approx \SI{7,7}{m.s^{-1}}$. Les frottements font perdre environ $35\,\%$ de la vitesse (et près de $60\,\%$ de l'énergie) : ils jouent ici un rôle de sécurité en limitant la vitesse d'arrivée.
\end{enumerate}
\textbf{Partie B.}
\begin{enumerate}
\setcounter{enumi}{5}
\item Relation fondamentale de la dynamique de rotation : $\mathcal{M}_f = J_\Delta\,\ddot{\theta}$, d'où $\ddot{\theta} = \dfrac{-10}{2{,}0} = \SI{-5,0}{rad.s^{-2}}$ : rotation uniformément décélérée. Arrêt lorsque $\omega = 0$ : $t = \dfrac{-\omega_0}{\ddot{\theta}} = \dfrac{100}{5{,}0} = \SI{20}{s}$.
\item Théorème de l'énergie cinétique entre la coupure et l'arrêt :
\[ 0 - \frac{1}{2}J_\Delta\,\omega_0^2 = \mathcal{M}_f\,\theta \quad\Rightarrow\quad \theta = \frac{J_\Delta\,\omega_0^2}{2|\mathcal{M}_f|} = \frac{2{,}0 \times 100^2}{2 \times 10} = \SI{1,0e3}{rad}. \]
Nombre de tours : $N = \dfrac{\theta}{2\pi} \approx 159$ tours. Le volant tourne longtemps après la coupure : c'est son rôle de réservoir d'énergie cinétique de rotation, qui régularise le fonctionnement du groupe électrogène.
\end{enumerate}
\textbf{Barème indicatif (sur 7 points) :} énergies mécaniques en $A$ et $B$ : 1 ; conclusion sur la conservation : 0,5 ; travail des frottements : 1 ; norme $f$ : 0,5 ; vitesse sans frottements et commentaire : 1 ; accélération angulaire et durée : 1,5 ; angle balayé et nombre de tours : 1,5.

\textbf{Points de vigilance :} préciser la référence des énergies potentielles ($z = 0$ au sol) ; pour conclure à la non-conservation, comparer explicitement $E_m(A)$ et $E_m(B)$ ; écrire $\Delta E_m = W(\vec{f})$ avec le signe négatif du travail résistant ; dans la partie B, le travail d'un couple constant est $W = \mathcal{M}\,\theta$ avec $\theta$ en radians ; convertir ensuite en tours via $N = \theta/(2\pi)$.
\end{corrige}

\newpage

\coursec{Fiche bilan}

\begin{popfigure}
\begin{tikzpicture}[
  mindmap,
  grow cyclic,
  every node/.style={concept, text=white, font=\small\bfseries},
  root concept/.append style={concept color=popdark, font=\normalsize\bfseries},
  level 1/.append style={level distance=3.4cm, sibling angle=60},
  level 2/.append style={level distance=2.2cm, sibling angle=45, font=\scriptsize},
  concept connection/.append style={color=popmuted}
]
\node[root concept]{Lois de Newton et énergie mécanique}
  child[concept color=popblue]{ node {Cinématique}
    child { node {repère de Frenet} }
    child { node {$\vec{a}_T$, $\vec{a}_N$} }
  }
  child[concept color=popgreen]{ node {1\textsuperscript{re} loi}
    child { node {référentiel galiléen} }
    child { node {principe d'inertie} }
  }
  child[concept color=poporange]{ node {2\textsuperscript{e} et 3\textsuperscript{e} lois}
    child { node {$\sum\vec{F}=m\vec{a}_G$} }
    child { node {actions réciproques} }
  }
  child[concept color=popgold]{ node {Rotation axe fixe}
    child { node {moment d'inertie $J_\Delta$} }
    child { node {$\sum M_\Delta=J_\Delta\ddot{\theta}$} }
  }
  child[concept color=poppurple]{ node {Énergie cinétique}
    child { node {$E_c=\frac12 mv^2$} }
    child { node {$\Delta E_c=\sum W(\vec{F})$} }
  }
  child[concept color=poppink]{ node {Énergie mécanique}
    child { node {$E_m=E_c+E_p$} }
    child { node {conservation} }
    child { node {frottements} }
  };
\end{tikzpicture}
\legende{Carte mentale de la leçon : les six piliers de la dynamique newtonienne et de l'approche énergétique.}
\end{popfigure}

\begin{aretenir}[L'essentiel de la leçon]

\textbf{1. Définitions clés}
\begin{itemize}
  \item \cle{Référentiel galiléen} : référentiel dans lequel le principe d'inertie est vérifié (ex. : référentiel terrestre pour des mouvements de courte durée et de faible amplitude).
  \item \cle{Repère de Frenet} $(M,\vec{T},\vec{N})$ : $\vec{T}$ tangent à la trajectoire (orienté dans le sens du mouvement), $\vec{N}$ normal, dirigé vers le centre de courbure de la trajectoire.
  \item \cle{Moment d'inertie} $J_\Delta$ (en $\mathrm{kg\cdot m^2}$) : grandeur positive qui mesure la résistance d'un solide à la mise en rotation autour de l'axe fixe $\Delta$ ; il dépend de la masse du solide et de sa répartition par rapport à l'axe.
  \item \cle{Moment d'une force} par rapport à l'axe $\Delta$ : $M_\Delta(\vec{F})=\pm F\,d$, où $d$ est le bras de levier (distance de l'axe à la ligne d'action de la force) ; le signe dépend du sens de rotation choisi comme positif.
  \item \cle{Énergie mécanique} : $E_m=E_c+E_p$, somme de l'énergie cinétique et de l'énergie potentielle du système.
\end{itemize}

\textbf{2. Formules et théorèmes à connaître par c\oe ur}
\begin{center}
\begin{tabularx}{\linewidth}{|l|X|}
\hline
\rowcolor{popblueL} \textbf{Grandeur / loi} & \textbf{Expression} \\
\hline
Accélération (Frenet) & $\vec{a}=\dfrac{dv}{dt}\,\vec{T}+\dfrac{v^2}{R}\,\vec{N}$, avec $R$ rayon de courbure \\
\hline
Principe d'inertie (et réciproque) & $\sum\vec{F}_{\mathrm{ext}}=\vec{0}\ \Longleftrightarrow\ \vec{v}_G=\overrightarrow{\mathrm{cste}}$ \\
\hline
2\textsuperscript{e} loi de Newton & $\sum\vec{F}_{\mathrm{ext}}=m\,\vec{a}_G$ (référentiel galiléen, $m$ constante) \\
\hline
3\textsuperscript{e} loi de Newton & $\vec{F}_{A/B}=-\vec{F}_{B/A}$ \\
\hline
RFD de la rotation & $\sum M_\Delta(\vec{F}_{\mathrm{ext}})=J_\Delta\,\ddot{\theta}$ \\
\hline
Énergie cinétique & translation : $E_c=\dfrac12 m v_G^2$ \quad rotation : $E_c=\dfrac12 J_\Delta\,\dot{\theta}^2$ \\
\hline
Théorème de l'énergie cinétique & $\Delta E_c=E_{c,B}-E_{c,A}=\sum W_{A\to B}(\vec{F})$ \\
\hline
Énergie potentielle de pesanteur & $E_p=m g z$ (à une constante additive près), $z$ altitude sur un axe vertical ascendant \\
\hline
Énergie potentielle élastique & $E_p=\dfrac12 k x^2$ (à une constante additive près), $x$ allongement algébrique du ressort \\
\hline
Conservation de $E_m$ & si seules des forces conservatives travaillent : $E_m=\mathrm{cste}$ \\
\hline
Frottements & $\Delta E_m=W(\vec{f})<0$ : l'énergie mécanique diminue \\
\hline
\end{tabularx}
\end{center}

\textbf{3. Méthodes express}
\begin{itemize}
  \item \cle{Méthode dynamique} : (1) définir le système et le référentiel galiléen ; (2) faire le bilan des forces extérieures et un schéma ; (3) appliquer $\sum\vec{F}=m\vec{a}_G$ (ou $\sum M_\Delta=J_\Delta\ddot{\theta}$) ; (4) projeter sur des axes bien choisis ; (5) intégrer en tenant compte des conditions initiales.
  \item \cle{Méthode énergétique} : idéale pour relier \textbf{vitesse et position} sans calculer le temps : $E_{c,B}-E_{c,A}=\sum W_{A\to B}(\vec{F})$, ou $E_{m,A}=E_{m,B}$ si les frottements sont négligeables.
  \item \cle{Choix de la méthode} : on demande une durée, une accélération, une force $\Rightarrow$ méthode dynamique ; on demande une vitesse en un point, une hauteur atteinte $\Rightarrow$ méthode énergétique.
  \item \cle{Travail du poids} : $W_{A\to B}(\vec{P})=mg(z_A-z_B)$, indépendant du chemin suivi (le poids est une force conservative).
  \item \cle{Forces qui ne travaillent pas} : toute force constamment perpendiculaire à la vitesse (réaction normale sans frottement, tension d'un fil inextensible en mouvement circulaire, force magnétique de Lorentz).
\end{itemize}
\end{aretenir}

\begin{aretenir}[Checklist avant le Bac]
\begin{itemize}
  \item $\square$ Je sais définir un référentiel galiléen et citer le principe d'inertie (et sa réciproque).
  \item $\square$ Je connais les deux composantes de l'accélération dans le repère de Frenet : $a_T=\dfrac{dv}{dt}$ et $a_N=\dfrac{v^2}{R}$.
  \item $\square$ Je fais toujours le \textbf{bilan des forces extérieures} avec un schéma avant d'écrire la 2\textsuperscript{e} loi.
  \item $\square$ Je sais projeter $\sum\vec{F}=m\vec{a}_G$ sur des axes et vérifier l'homogénéité des expressions.
  \item $\square$ Je sais énoncer la 3\textsuperscript{e} loi : $\vec{F}_{A/B}=-\vec{F}_{B/A}$, forces de même nature appliquées à des corps différents.
  \item $\square$ Je sais écrire la RFD de la rotation $\sum M_\Delta=J_\Delta\ddot{\theta}$ avec les bons signes pour les moments.
  \item $\square$ Je connais les moments d'inertie usuels (cerceau $J_\Delta=mR^2$, disque ou cylindre $J_\Delta=\tfrac12 mR^2$, tige autour de son centre $J_\Delta=\tfrac{1}{12}mL^2$).
  \item $\square$ J'écris le théorème de l'énergie cinétique avec $\Delta E_c=E_{c,\mathrm{final}}-E_{c,\mathrm{initial}}$ (jamais l'inverse).
  \item $\square$ Je sais que $W(\vec{P})=mg(z_A-z_B)$ ne dépend pas du chemin suivi.
  \item $\square$ Je vérifie si les frottements existent avant d'écrire la conservation de $E_m$ ; sinon $\Delta E_m=W(\vec{f})$.
  \item $\square$ Je choisis la méthode adaptée : dynamique pour accélérations/durées, énergie pour vitesses/positions.
  \item $\square$ Je donne chaque résultat avec son unité SI et un nombre cohérent de chiffres significatifs.
\end{itemize}
\end{aretenir}

\findecours

\end{document}
